% =====================================================================================
%  PRESENT BIAS IS A PROPERTY OF THE CHOICE, NOT THE CHOOSER  --  the writing sample
%  Alessandro Facchini -- 15 September 2026
%  Build (from this folder):  pdflatex main_short ; bibtex main_short ; pdflatex main_short ; pdflatex main_short
%  Main text is sections 1-9 (target: 25 pages including figures and tables); the
%  references and the Online Appendix follow in the same file and are to be split off.
%  The long paper (../main.tex) holds every number, table and robustness check cited here.
% =====================================================================================
\documentclass[11pt,a4paper]{article}

\usepackage[utf8]{inputenc}
\usepackage[T1]{fontenc}
\usepackage{lmodern}
\usepackage[a4paper,margin=2.25cm]{geometry}
\usepackage{amsmath,amssymb,amsthm}
\usepackage{enumitem}
\usepackage{booktabs}
\usepackage{array}
\usepackage{graphicx}
\usepackage{xcolor}
\usepackage{float}
\usepackage{multicol}
\usepackage{threeparttable}
\usepackage{microtype}
\usepackage{caption}
\usepackage[hidelinks]{hyperref}
\usepackage[round]{natbib}
\bibliographystyle{plainnat}
\graphicspath{{../figs/}{../figs/diagrams/}}

\captionsetup{font=small,labelfont=bf}
\linespread{1.12}
\setlength{\parskip}{1.5pt}

\theoremstyle{plain}
\newtheorem{proposition}{Proposition}
\theoremstyle{definition}
\newtheorem{definition}{Definition}

\newcommand{\E}{\mathbb{E}}
\newcommand{\Var}{\mathrm{Var}}
\newcommand{\R}{\mathbb{R}}
\newcommand{\beff}{\beta_{\mathrm{eff}}}


\title{\textbf{Beyond the Trait:\\Present Bias as a Reaction Function}\\[4pt] % Beyond Fixed Traits
\large Revealed Preference, %Foundations, 
Empirical Evidence, and State-Dependent Self-Control}

%\title{\textbf{From Trait to State:\\Present Bias as a Reaction Function}\\[4pt]
%\large Revealed Preference, %Foundations, 
%Recentered Identification, and State-Dependent Self-Control}

%\title{\textbf{Present Bias Is a Property of the Choice,\\ Not the Chooser}\\[4pt]
%\large A Revealed-Preference Audit of the Present-Bias Parameter}


\author{Alessandro Facchini
%\thanks{Writing sample, 01 October 2026. All estimates come from a single replication pipeline built on public replication packages; every number in the text traces to a named code block, and a longer version of the paper with every specification, certificate and robustness table is available on request. The Online Appendix referred to below accompanies this file. I am grateful to the authors of the replication packages reanalyzed here: in the headline case the identifying variation is theirs, collected for a different purpose, and the paper is explicit throughout about what is reproduced and what is new. \textcolor{red}{to check}}
}

\date{10 October 2026}

\begin{document}
\maketitle

\begin{abstract}
\noindent
Present bias is often estimated, reported and calibrated as a personal parameter: a common $\beta$ across states, goods, and settings.
%This paper tests the empirical validity of this unified representation.
Re-examining data from three independent experimental designs, I test the invariance restrictions imposed by a single $\beta$ and find that several fail.
%A single $\beta$ imposes a bundle of invariance restrictions; re-examining data across three independent experimental setups, I show that core conditions systematically fail.
%%%
\emph{First}, in a real-effort panel \citep{augenblick2019}, randomizing whether choices precede or follow a mandatory work block shifts the estimated $\beta$ from $0.88$ (before) to $0.74$ (after) within subjects.
This exogenous state variation reduces immediate effort while leaving choices and forecasts about future work unchanged, rejecting the cross-vantage equality assumed by pooled specifications.
These findings are robust to design-recentering on the treatment's conditional expectation under exact randomization inference.
%%%
\emph{Second}, applying a bivariate measurement-error model to the data of \citet{ans2015} to account for sampling noise reveals genuine dispersion in present bias over effort, but none over money, among the same individuals.
This contradicts the stable-trait hypothesis, since no correlation can arise between a variable and a constant.
%%%
\emph{Third}, using a common protocol across rewards \citep{cheung2022}, %I likewise find stronger present bias over money than over food.
I confirm their finding of stronger present bias over money than over food.
However, differences in $\beta$ across snack subgroups exceed the money--food gap, challenging commodity-level explanations.
Individual-level correlations suggest trait-like stability within trade-off structures, but not across them.
Present bias thus appears to depend on the person--reward pairing, conditional on trade-off structure and state, rather than on either component alone. 
%\emph{Third}, in a protocol that varies only the reward \citep{cheung2022}, choices exhibit stronger present bias over money than over food. 
%However, I show that commodity categories mask the underlying structure, as within-food dispersion in $\beta$ exceeds the overall money--food gap.
%Individual-level correlations suggest that present bias may behave like a trait within a given trade-off structure, yet it fails to generalize across structures. Indeed, it attaches to the specific person--reward combination rather than to either component alone.
%%%
Uncovering these patterns requires an appropriate testing framework: the level statistics conventionally used to document heterogeneity in $\beta$ cannot distinguish genuine state dependence from stochastic variation under a state-independent model.
I derive a valid class of test statistics using recentered functions of the assigned state and provide a revealed-preference characterization of state-dependent quasi-hyperbolic discounting that extends \citet{eis2020}, nesting exponential discounting and the fixed-trait specification as special cases.
I also establish coordinate-level conditions for eliminating non-linear transformation bias and develop nonparametric tests that relax functional-form restrictions on utility and cost functions.
%%%
To establish common ground across normative perspectives, I conservatively define welfare losses as the unanticipated component of the plan-versus-action wedge, showing empirically that it constitutes the majority of the observed wedge.
Accordingly, commitment demand cannot serve as a reliable diagnostic of bias, as it tracks the accurate component of people's forecasts rather than their prediction errors.
%%%
Finally, I present a framework of state-weighted valuation systems that formalizes present bias as a state-dependent reaction function, recovering traditional models under specific restrictions.
%%%
Overall, these results broaden the scope of policy beyond individual targeting and voluntary self-correction to encompass the choice environment, recognizing that the decision context can embed the very cues that shape present bias.


\iffalse %OLD claude
Present bias is estimated, reported and calibrated as a personal parameter: one $\beta$ per person, carried across goods, states and settings. This paper asks whether that object exists. A single $\beta$ imposes a bundle of invariances, and three designs, sharing no data and no identification strategy, each break a different one. \emph{First}, in the canonical real-effort panel, a within-person randomization of whether choices are made before or after a block of mandatory work moves the present-bias factor from $0.88$ to $0.74$ in the same hundred people: the state lowers the act and leaves the plan and the forecast unchanged, and the effect survives design-recentering of the randomized order under exact randomization inference. \emph{Second}, once sampling noise is modeled, individual present bias in the same people is a variable over effort and a constant over money. \emph{Third}, in a protocol that varies only the reward, money is the \emph{most} present-biased of three goods, and within food $\beta$ follows the person who chooses the snack. These answers required repairing the test: the level statistic used to document heterogeneity in $\beta$ rejects with probability one under every null in which $\beta$ does not depend on the state. I characterize the valid class -- recentered functions of the assigned state -- and give a revealed-preference characterization of state-dependent quasi-hyperbolic discounting that extends \citet{eis2020} and nests the trait model. Counting as welfare-relevant only the part of the wedge that people fail to anticipate, I find that this part is most of it, and that demand for commitment tracks the accurate part of people's forecasts rather than their errors. A framework of state-weighted valuation systems derives this reaction function and nests the standard models. Present bias is less a trait of the chooser than a response to the moment of choice, which moves the focus of policy from the person to the environment in which choices are made.
\fi

%\medskip
%\noindent\textbf{JEL:} C12, C14, C26, C91, D01, D15, D91. \quad \textbf{Keywords:} present bias, state dependence, quasi-hyperbolic discounting, revealed preference, recentered instruments, commitment.
%%% %randomization inference, misprediction, 
\end{abstract}












\clearpage
% =====================================================================================
\section{Introduction}
\label{sec:intro}



Why do people act against their own plans -- procrastinating, saving less than intended, %undersaving
or making different choices about food, substances, or screens than they had planned? A standard answer in economics is \emph{present bias} (PB): a systematic over-weighting of the present that makes today's self sabotage yesterday's plans. Formalized as quasi-hyperbolic discounting \citep{laibson1997,odonoghue1999}, PB has been estimated in a large body of empirical work, incorporated into models of optimal policy and behavioral interventions, and typically treated as a \emph{personal parameter}: a single number, $\beta$, that characterizes an individual. This paper asks whether that object exists.

According to \citet{lockwood2020}, estimates of $\beta$ range from $0.34$ to $1.03$, a spread that no single study's confidence interval encompasses. 
While such dispersion is commonly attributed to imprecision, this paper advances a different possibility: the spread is not merely noise around a common value, but a trace of genuine heterogeneity in the object being measured.
Before asking whether $\beta<1$, I therefore ask a more fundamental question: does a single $\beta$ exist in the first place, and what exactly is it a property of -- a person, a domain, a moment, or something else?

In this sense, the analysis develops a set of ideas that have been partly present in the theoretical literature but have not been fully carried through to empirical measurement.
Specifically, \citet{loewenstein1996} argued that the weight placed on the present increases with ``visceral states,'' and subsequent work developed cue-triggered models in which PB varies with the salience of such cues \citep{laibson2001,bernheimrangel2004}. %More recently, 
\citet{pennesi2020} derived from axioms a quasi-hyperbolic representation where the PB factor is a function of cue intensity, while \citet{odonoghue2015} identify heterogeneity across domains and states among the key open questions for the literature.
%%%
Yet empirical studies typically estimate $\beta$ as if it were constant. Even when the literature attempts to document heterogeneity in $\beta$, the evidence usually invoked -- a profile of level estimates across groups, goods, or conditions %, each significantly below one 
-- is methodologically unsuited to distinguish it from a fixed trait.
This paper quantifies the inferential cost of assuming a constant $\beta$ and provides an econometric framework to identify state-contingent PB.


\paragraph{Main findings.} A single $(\beta,\delta)$ per person imposes a bundle of invariance restrictions: for instance, the same $\beta$ must govern discounting for a dated payment and dated effort, for cash and a cigarette, for a rested and an exhausted person, and between an advance plan and a proximal choice.
%Each is a separate restriction, testable through a targeted contrast that isolates it. 
%Three independent designs each reject a distinct invariance restriction (Figure~\ref{fig:three}).
Each of the three independent designs rejects a distinct invariance restriction (Figure~\ref{fig:three}).




\emph{The state at the moment of choice.} 
\citet{augenblick2019} elicit real-effort allocations from $100$ subjects over seven weeks. Each session requires ten mandatory ``unpleasant'' transcription tasks, and the authors \emph{randomize} whether work preferences are stated before or after this first block to test for the presence of projection bias. That randomization induces within-person variation in \emph{visceral state} at the moment of choice, %orthogonally 
crossed with choices over dated effort--wage allocations.
Under this interpretation, conditioning on the %visceral 
state at decision time reveals that the same $100$ subjects exhibit $\beta=0.88$ %$[0.80,0.97]$ 
before the initial set of tasks, against $\beta=0.74$ %$[0.65,0.82]$ 
after it ($\Delta\log\beta=-0.180$, $p=.008$).
The published estimate for this panel ($\beta=0.83$) is therefore a state-weighted average. Notably, the state manipulation has a sharply asymmetric effect: under the original paper's cost-function specification, %$\log(e+10)$
it shifts choices about \emph{immediate} work by $-0.14$ log points ($p<.01$), while leaving choices about \emph{future} work unchanged ($+0.01$, $p=.64$).
However, the projection-bias model posited by the authors applies a single uniform shift to every decision made after the mandatory work block. In other words, it assumes that visceral fatigue distorts an advance plan for a future date just as much as an immediate choice for today -- imposing cross-vantage equality as an identifying restriction. The data %themselves 
reject this restriction ($p<.007$), revealing a distinction that the original specification does not capture.
Finally, because the experiment randomizes the \emph{order} of a session's components rather than state labels directly, decisions are unequally exposed to the draw, given the fixed session structure. The visceral state is thus a function of an exogenous shock and of predetermined exposure, in the sense of \citet{borusyak2023}.
Recentering the treatment on its conditional expectation purges this imbalance while leaving the estimated effects virtually unchanged, yielding design-recentered PB factors of $\beta_{\mathrm{before}}=0.85$ and $\beta_{\mathrm{after}}=0.75$.

%AF from LM: L'esperimento non randomizza direttamente un'etichetta "hot" o "cold" sulle singole decisioni. Randomizza invece l'ordine delle componenti all'interno di ciascuna sessione giornaliera. Una sessione contiene fino a 5 set di decisione — il lavoro di oggi ($C$), 1 o 2 set sul lavoro futuro ($F_1, F_2$), 1 o 2 set di predizioni ($G_1, G_2$) — più 1 blocco obbligatorio di 10 task di trascrizione.Una decisione è definita "hot" ($s_k = 1$) se viene presa dopo aver completato il blocco obbligatorio (cold altrimenti).
%L'esposizione è diseguale perché la struttura delle sessioni ($w$) segue una sequenza prestabilita, in cui certi tipi di decisioni tendono ad occupare sistematicamente posizioni iniziali o finali nella sessione. Eg, il lavoro di oggi ($C$) tende a essere piazzato all'inizio della sessione; di conseguenza, si trova nello stato "hot" solo nel 46% dei casi ($\mu = 0.46$).Il primo set sul futuro ($F_1$) si trova nello stato "hot" nel 60% dei casi ($\mu = 0.60$). Il secondo set sul futuro ($F_2$) viene posizionato quasi sempre in fondo alla sessione, quindi si trova nello stato "hot" ben l'83% delle volte ($\mu = 0.83$). etc, see BH paper and \label{fig:exposure} below (and the part around).



\emph{The person.}
Differences in average $\beta$ across domains are %potentially %\emph{a priori} 
consistent with a level shift that preserves the ranking of individuals. 
The relevant question is therefore whether relative differences persist, rather than whether %parameter 
levels coincide. I address it using a bivariate measurement-error model that separates genuine preference heterogeneity from sampling noise. 
%Therefore, the question of individual stability becomes sharper, and this paper answers it after netting out each subject's own noise through a bivariate measurement-error model. 
In the data of \citet{ans2015}, %where the same people allocate both effort and money, 
this method reveals that the standard deviation of true %individual 
$(\log\beta)$ is $0.24$ over effort and $0.00$ over money, within the same people. %That is, PB over money has no residual variation across these subjects.
This is fundamentally inconsistent with the stable-trait hypothesis, as a single parameter cannot be a variable in one setting and a constant in another. 
In contrast, when comparing choices across different rewards within the same intertemporal trade-off structure -- using the three-good panel of \citet{cheung2022} for money, healthy food, and unhealthy food -- the same model recovers a high true cross-reward correlation ($0.71$--$0.92$). PB thus behaves like a trait \emph{within} a trade-off structure, but fails to remain a single object \emph{across} structures.


\emph{The good.} 
A widely cited finding in this field is that the same individuals exhibit PB over effort but not over monetary rewards \citep{ans2015}. This is conventionally interpreted as a structural property of commodities: money is fungible, storable, and ``cold'', %(hence uncorrupted by immediate temptation), 
whereas effort is experienced as it is exerted and generates immediate visceral friction. 
However, the data collected by \citet{cheung2022} point in a different direction. The authors hold potential confounders fixed by employing an identical protocol across three goods -- money, healthy food, and unhealthy food -- for $697$ subjects. %, maintaining the same convex-budget price grid, the same three-week delay, and the same advance versus on-the-day vantages. 
While the traditional commodity reading predicts money to be the \emph{least} present-biased of the three, 
% re-estimating their data while relaxing their baseline functional-form constraints reveals it to be the \emph{most}, corroborating their finding.
their own estimates show it to be the \emph{most}, and re-estimating their data with the scale-free estimator of Section~\ref{sec:coord} confirms this ordering. 
In particular, allowing for good-specific utility curvature, PB over money ($\beta_{\mathrm{money}}=0.82$) is significantly stronger %($p<.004$) 
than over food ($\beta_{\mathrm{healthy}}=0.85$ and $\beta_{\mathrm{unhealthy}}=0.86$, with the two categories statistically indistinguishable).
Crucially, my analysis reveals that a commodity ranking is the wrong theoretical object. Within each food category, subjects selected one of three specific snacks; estimating $\beta$ across these six snack-choice subgroups shows that PB spans $0.82$--$0.89$ under the single label ``food'' --- a range wider than the entire money--food gap. 
PB is therefore not an immutable property of a commodity class, but a reaction function shaped by choice architecture, immediacy, and the person--reward pair.
%PB therefore attaches not to a commodity class but to the pair formed by a person and a reward


\begin{figure}[H]\centering
\caption{\footnotesize%
Same choices, three readings: single-parameter estimation and broken invariances.}
\includegraphics[width=1\linewidth]{d01_three_layers_wide.pdf}
%\caption{\footnotesize%
%One parameter, three readings of the same choices. \emph{Left:} the published estimates of $\beta$ the field has collected, running from $0.34$ to $1.03$. \emph{Right:} the three invariances the paper breaks. \textbf{State:} in the same $100$ people, the randomized block of work moves the choice about today's work (the \emph{act}) and not the choice about future-dated work (the \emph{plan}). Measured against the plan made in the same state, the act gives $\beta=0.88$ cold and $0.74$ hot; the plan made hot, measured on the same scale against the plan made cold, is $1.04$. The dotted line is the single number an estimate that ignores the state returns, $0.80$ --- the object the source paper publishes as $0.83$ on its $72$-subject estimation sample. \textbf{Structure:} in the same $101$ people, individual $\log\beta$ has a true standard deviation of $0.20$ over effort and $0.00$ over money. \textbf{Good:} in the same $697$ people, money is the most present-biased of three rewards elicited on one price grid.
%by state (the same $100$ people before and after an assigned block of work), by structure (the same $101$ people over effort and over money), by good (the same $697$ people over three rewards on one price grid).}
\label{fig:three}
\end{figure}
% FIGURE SOURCE (AF: edit here to change the figure)
%   ../figs/diagrams/d01_three_layers_wide_body.tex   <- the drawing (edit this one)
%   ../figs/diagrams/d01_three_layers_wide.tex        <- wrapper, palette "af" (used here)
%   ../figs/diagrams/d01_three_layers_wide_saffron.tex / _meadow.tex  <- other palettes
%   build:  cd ../figs/diagrams  &&  python make_palette_variants.py
%           (or: pdflatex d01_three_layers_wide)
%   Palettes: ../figs/diagrams/dxpal_af.tex, dxpal_saffron.tex, dxpal_meadow.tex;
%   side-by-side comparison: ../figs/diagrams/palette_gallery.pdf. To switch palette,
%   change the file name above AND Figure 2's (d02_axes_edited_<palette>.pdf).
%   The header of the body file lists every number in the figure and the pipeline block
%   (summary CSV, block, row name) each one comes from.
%   Earlier versions: ../figs/diagrams/_backup_fig1_0924/ and d01_three_layers.tex.


\paragraph{Methodological prerequisite: repairing the test.} Uncovering these empirical patterns requires first repairing the econometric toolkit, as the statistic conventionally used to evaluate state dependence cannot validly test it.
Under a null hypothesis in which $\beta$ is drawn from an arbitrary, heterogeneous, and serially dependent process that is \emph{independent of the state}, a test built on the \emph{level} of the PB wedge within state cells rejects the null of exponential discounting ($\beta = 1$) with probability one. Crucially, this approach can only detect the presence of PB, not its dependence on the state: while cell-by-cell estimates provide valid tests of time consistency, a \emph{profile} of such estimates cannot establish that $\beta$ varies with the conditioning variable.
%%%
Section~\ref{sec:stat} formalizes this size distortion and characterizes the class of valid test statistics. 
To evaluate state dependence, %cleanly, 
any valid test statistic must project the decision-level wedge $g_k$ (the plan-versus-action gap) onto an instrumental vector $z_k$ that is relevant to the assigned state $s_k$ while remaining orthogonal to baseline PB levels and predetermined confounds $w$ (including subject fixed effects and session structure) through a  %cross-cell pooled
balance condition, $E[z_k \mid w] = 0$.
In particular, when state assignment is governed by experimental structure, the unique class of weighting functions %(i.e., instruments) 
that delivers exact size %control 
against every state-independent null without ad-hoc error assumptions consists of %design-
recentered functions of the %assigned 
state -- formalized by the formula instruments of \citet{borusyak2023,borusyakhull2026}.%
\footnote{For example, an individual setting daily screen-time goals and facing a visceral state $s_k \in \{0,1\}$ (where $s_k=1$ denotes late-night exhaustion and $s_k=0$ daytime alertness) exhibits a plan-versus-action wedge $g_k$ %$= \log(\text{scrolling}_{\text{action},k}) - \log(\text{scrolling}_{\text{plan},k})$ 
for each decision $k$. 
%Traditional %practice would run 
%per-cell level tests against exponential discounting ($\beta=1$) could find, for instance, $\hat{\beta}_{\text{day}} = 0.85 < 1$ and $\hat{\beta}_{\text{night}} = 0.65 < 1$. 
Under any state-independent null where PB is an arbitrary stochastic process ($\beta_k \perp s_k$), a test of whether cell-level estimates differ from $1$ rejects asymptotically with probability one (e.g., it could find that $\hat{\beta}_{\text{day}} = 0.85 < 1$ and $\hat{\beta}_{\text{night}} = 0.65 < 1$).
A comparison of the two estimates is informative only if day and night decisions are otherwise alike: when the state marks a kind of person or follows a schedule, the difference also picks up whatever else differs between the cells. %(evaluating $\beta_{\text{night}} \neq \beta_{\text{day}}$ under a null defined only by $\beta_k < 1$).
To test state dependence cleanly, one must project $g_k$ onto an instrumental weighting vector $z_k$ that is relevant to the assigned state $s_k$ but orthogonal to baseline PB levels and predetermined confounds via $E[z_k \mid w] = 0$. Intuitively, $z_k$ acts as a contrast weight that cancels out baseline individual heterogeneity -- isolating whether scrolling behavior fluctuates specifically in response to state shifts, rather than reflecting the individual's time-invariant tendency to over-scroll.
When state assignment $s_k$ is governed by experimental schedules $w$, a recentered instrument $z_k = s_k - \mu_k$ subtracts each decision's conditional exposure probability $\mu_k = E[s_k \mid w]$. This measures whether an observation is \emph{unexpectedly} in a visceral state relative to its schedule expectation, providing a clean exogenous shock that delivers exact size 
%control at the nominal 5\% level 
against every state-independent null.
}

This result underpins the diagnostic ladder of Section~\ref{sec:object}: its rungs translate nested structural restrictions (from unconstrained state dependence down to exponential discounting) into specific balance conditions on the weighting vector, mapping each testable step to a formal hypothesis, a valid test statistic, and an empirical verdict.



\paragraph{Contributions.} 
\emph{Empirically}, the PB factor is shown not to be a personal constant across three independent designs. 
\emph{Theoretically}, a revealed-preference characterization of state-dependent quasi-hyperbolic discounting is developed, extending \citet{eis2020} and nesting a ladder of restrictions including the trait model.
\emph{Econometrically}, this paper demonstrates that traditional level statistics fail to evaluate state dependence, and characterizes the valid class of design-recentered formula instruments. %alongside its efficient member. 
In addition, robust testing procedures are developed, evaluated via Monte Carlo simulations, and applied to empirical panels. These include an audit of the estimator's coordinate establishing the conditions to eliminate non-linear transformation bias, a bivariate measurement-error model that separates true preference heterogeneity from sampling variance in $\beta$, and distribution-free nonparametric tests that relax functional-form assumptions on utility and cost functions. 
\emph{From a welfare perspective}, this study adopts a conservative criterion that grounds welfare losses strictly in the unanticipated component of the plan-versus-outcome wedge -- establishing a baseline that can command broad agreement across normative standpoints. 
Across diverse settings, including high-stakes addiction, most of the observed wedge stems from systematic misprediction. Consequently, commitment demand cannot serve as a reliable diagnostic of bias, as it tracks the accurate component of people's forecasts rather than their prediction errors.
Overall, these results broaden the scope of policy beyond individual targeting and voluntary self-correction to encompass the choice architecture, recognizing that decision environments can embed the very cues that shape PB.
Finally, the Multi-System Control framework provides a \emph{behavioral micro-foundation} for PB as a state-dependent reaction function, unifying existing approaches by nesting standard models -- including quasi-hyperbolic discounting and dual-self architectures -- as restricted special cases.  
%Everything runs on thirty-one public replication packages with one estimator, one coordinate and one inference engine, with every test's size and power certified by simulation before it touched data.

\paragraph{Related literature.} On measurement and market confounders, \citet{andreoni2012a} introduced convex time budgets, while \citet{ans2015} incorporated real effort and identified the domain gap re-examined below. \citet{augenblick2019} constructed the multi-vantage effort panel that serves as the main empirical design here. \citet{balakrishnan2020} demonstrated that PB over money is highly sensitive to the precise timing of immediate delivery, while \citet{deansautmann2021}, \citet{cassidy2019} and \citet{andreoni2012b} documented how credit constraints and risk can distort monetary wedges.

On state dependence, the conceptual foundations trace back to \citet{loewenstein1996}, \citet{lor2003}, and \citet{kruegeluhl2023}. Theoretical models subsequently incorporated state dependence through \citet{laibson2001}, \citet{bernheimrangel2004}, \citet{fudenberg2006}, and \citet{pennesi2020}, whose cue-dependent factor is generalized in Section~\ref{sec:model}. %
Direct empirical evidence is more limited.
Psychology studies manipulate visceral states and measure monetary impatience \citep[e.g.,][]{badger2007,sayette2008,kimzauberman2013,skrynka2019}, but mostly focus on behavioral proxies rather than a structural $\beta$; similarly, cognitive-load studies evaluate static impulse control without a dated trade-off \citep{shiv1999,wardmann2000}.
Where quasi-hyperbolic $\beta$ is estimated across states, the state is typically assigned \emph{between} subjects and exclusively over monetary rewards. Induced sadness lowers estimated $\beta$ from $0.98$ to $0.94$ \citep{lerner2013}, fasting from $1.00$ to $0.95$ \citep{ashton2015}, while a prior self-control task raises it and flattens utility curvature -- the opposite of what depletion models predict \citep{kuhn2017}.
\citet{carvalho2016} find a payday effect on monetary choices that they attribute to liquidity constraints rather than PB, with no effect on real-effort decisions.
Features of the choice itself, by contrast, shift $\beta$ more substantially. Delivering the ``immediate'' reward hours rather than minutes after the session eliminates PB \citep{balakrishnan2020}, while across $220$ published convex-time-budget estimates, reward type and delivery immediacy are the primary sources of heterogeneity \citep{imai2021}. Existing designs, however, lack a within-person structure combining assigned states with dual-vantage elicitation -- a combination needed to distinguish a state-dependent $\beta$ from changes in utility curvature, baseline valuation, and measurement error.

On revealed preference over time, this paper extends the axiomatic 
foundations of \citet{eis2020}, building on a broader nonparametric literature on time consistency \citep{halevy2015,browning1989,crawford2010,adams2014,bbc2021,aguiar2021}. %DAVE-GARP
For model testing and econometric 
inference, \citet{adaocostinotdonaldson2024} motivate the IV framework used to evaluate structural predictions, \citet{arp2023} and \citet{coxshi2023} provide the moment-inequality methods.
Finally, \citet{borusyak2023,borusyakhull2026} 
introduce the recentering approach applied here to PB.


The remainder of the paper is organized as follows. Section~\ref{sec:object} defines the empirical object and core hypotheses; Section~\ref{sec:stat} derives the valid test statistic; Section~\ref{sec:data} describes the data; Sections~\ref{sec:state}--\ref{sec:magn} present the empirical findings and their implications; Section~\ref{sec:model} introduces the theoretical framework that micro-founds state dependence and reconciles standard models; and Section~\ref{sec:limits} concludes.%limitations? or just discussion? or future extension/research?



% =====================================================================================
%\section{The object, and the hypothesis}
\section{Theory: Revealed Preference and the Restriction Ladder}
\label{sec:object}

An intertemporal choice has at least three independent descriptors, and a single $\beta$ must be invariant along all three: the \emph{sign--timing structure} (Axis 1 of Figure~\ref{fig:axes_edited}), the \emph{good} (e.g., money, effort, food, alcohol), and the \emph{state} at the moment of choice (e.g., cued or not, hungry or sated, depleted or rested).
These axes are orthogonal, but existing empirical designs typically are not. The money literature operates almost entirely in \textit{R--R} structures under neutral (cold) visceral states, while the effort literature focuses on \textit{C--C} structures. Consequently, comparing $\beta_{\rm money}$ and $\beta_{\rm effort}$ changes the good as well as the sign--timing structure. Each empirical design examined in this paper instead moves along specific axes while holding the others fixed.

\begin{figure}[H]\centering
\caption{\footnotesize
The three axes and the three experiments.}
\includegraphics[width=.99\linewidth]{diagrams/d02_axes_edited.pdf}
%Columns are sign--timing structures, rows are goods, and the state at the moment of choice is a layer inside a cell. The money literature occupies one cell and the effort literature another, so ``$\beta_{\rm money}$ versus $\beta_{\rm effort}$'' compares structures, not merely goods. Each of the paper's designs moves along exactly one axis with the other two held fixed: Augenblick--Rabin assign the state inside the (C--C, effort) cell; Cheung et al.\ vary the good inside the R--R column; ANS step across structure and good at once.}
\label{fig:axes_edited}
\end{figure}
% "Note on notation: First and second positions denote immediate ($t_0$) and delayed ($t_1$) timing, respectively; superscripts $+$ and $–$ specify positive (reward) and negative (cost) utility flows within mixed trade-offs."
%  * In \\(R-R\\) e \\(C-C\\), entrambe le componenti hanno la stessa valenza (due ricompense o due costi).
 %  * In \\(R^+ C^-\\), l'apice \\(+\\) indica che il flusso a \\(t_0\\) entra con **segno positivo** nella funzione di utilità (\\(+u(x)\\)), mentre l'apice \\(-\\) indica che il flusso a \\(t_1\\) entra con **segno negativo** (\\(-c(y)\\)).
%  * In \\(C^+ R^-\\), l'apice \\(+\\) su \\(C^+\\) indica che l'azione \\(C\\) viene erogata a \\(t_0\\) come sforzo positivo (\\(e > 0\\), che dà disutilità \\(-C(e)\\)), mentre \\(R^-\\) è il beneficio differito a \\(t_1\\).


\noindent
A second distinction used throughout the paper is between two \emph{vantages} (Figure~\ref{fig:vantage}). From sufficiently far ahead, both the cost and the reward lie in the future, so the PB factor multiplies both and cancels. The advance plan therefore reveals the long-run preference shared across selves ($\delta$). On the day the cost is borne, instead, $\beta$ multiplies the delayed reward alone, generating a plan-versus-outcome wedge. The difference between the two first-order conditions yields $\log\beta$ exactly. 
% and it is the reference against which people's own forecasts and choices are compared [...in welfare considerations---(?)] (Section~\ref{sec:magn}.).
This identity is the starting point for every estimate in the paper (Section~\ref{sec:coord}) and explains why the advance vantage and its prediction serve as the welfare benchmark in Section~\ref{sec:magn}.



\begin{figure}[H]\centering
\caption{\footnotesize%
The two vantages: $\beta$ cancels from the plan and shows in the act.}
\includegraphics[width=.99\linewidth]{d03_vantage.pdf}
\label{fig:vantage}
\end{figure}
% fig:vantage is drawn in figs/diagrams/d03_vantage.tex (TikZ; pdflatex d03_vantage); 10-07 redraw in the notation of sec:coord, with the forecast and the state (the old version is figs/diagrams/_backups/d03_vantage_before_1007.tex). Red caption: AF to decide whether and where to keep it.




\iffalse
\begin{figure}[H]\centering
\includegraphics[width=.94\linewidth]{fig_taxonomy.pdf}
\caption{\footnotesize%
The four sign--timing structures. Rows are the sign of the \emph{immediate}
component. A pure $(\beta,\delta)$ model prices all four with the same present-bias
factor, because only the dates of the utils enter. Each cell names the datasets that
populate it here.}
\label{fig:taxonomy}
\end{figure}
\fi

\iffalse
\begin{figure}[H]\centering
\includegraphics[width=.98\linewidth]{figs/diagrams/d03_vantage.pdf}
\caption{The vantage timeline. At the \emph{advance} vantage both the cost and the
reward are in the future, so any present-bias factor multiplies both sides of the
trade-off and cancels: advance choices reveal the long-run preferences every self shares.
At \emph{today}, the cost is immediate and $\beta$ hits one side only. The difference of
the two first-order conditions is $\log\beta$ exactly, and in effort units it is the
log-wedge $\theta$ that every estimate in the paper starts from
(Section~\ref{subsec:scalefreelog}). The welfare wedge $\Gamma$ of Section~\ref{sec:welfare}
is the same person disagreeing with herself between the two vantages, with no outside
standard imported.}
\label{fig:vantage}
\end{figure}
\fi




\iffalse
\subsection{A revealed-preference characterization with a state-dependent $\beta$}
\label{subsec:RPtheory}

A central strength of the nonparametric revealed-preference (RP) approach is that it tests structural predictions without auxiliary parametric assumptions on the utility function, requiring only basic regularity conditions, namely monotonicity and concavity \citep{afriat1967, varian1982, browning1989}. In the domain of intertemporal choice, \citet{eis2020} advance this tradition by establishing necessary and sufficient conditions for a dataset to be rationalized by quasi-hyperbolic discounting (QHD) with a single PB parameter $\bar{\beta}$, invariant for each individual across choices. I generalize this characterization to accommodate a \textit{state-dependent} PB factor $\beta(c)$ that varies across exogenously assigned state cells $c \in \{1,\dots,C\}$. %
Formally, a decision maker makes $K$ choices. In decision $k$, they face prices $p^k = (p_0^k, p_1^k, \dots, p_T^k)$ over dated consumption and choose an affordable stream $x^k = (x_0^k, x_1^k, \dots, x_T^k)$, where $t \in \{0, \dots, T\}$ indexes the delay from the moment of decision. Each observation is assigned to a state cell $c(k)\in\{1,\dots,C\}$.


\begin{definition}[State-Dependent QHD Rationalization]
The data are \emph{SQH-rationalized} by $(\beta(\cdot),\delta,u)$ if there exist one
increasing concave $u$, one $\delta\in(0,1]$ and a function $\beta:\{1,\dots,C\}\to\R_{++}$
such that every $x^{k}$ maximizes $\sum_t D_{c(k)}(t)\,u(x_t)$ over the budget, with
$D_c(0)=1$ and $D_c(t)=\beta(c)\delta^{t}$ for $t>0$. The trait model (QHD) is
$\beta(c)\equiv\bar\beta$; exponential discounting (EDU) is $\bar\beta=1$.
\end{definition}

The utility index and $\delta$ are common across cells, while $\beta$ is allowed to vary by state cell.
Following \citet{eis2020}, a sequence of paired dated consumption entries $(x_{t_i}^{k_i}, x_{t'_i}^{k'_i})_{i=1}^n$ is \emph{balanced} if each observation $k$ appears as often on the left (and, consequently, as numerator $k_i$ of the multiplier ratios $\lambda^{k_i}/\lambda^{k'_i}$) as on the right (as denominator $k'_i$).
\footnote{This balancing ensures that in the \textcolor{red}{FOC product}, $\prod_{i=1}^n \frac{u'(x_{t_i}^{k_i})}{u'(x_{t'_i}^{k'_i})} = \left(\prod_{i=1}^n \frac{\lambda^{k_i}}{\lambda^{k'_i}}\right) \times \left(\prod_{i=1}^n \frac{p_{t_i}^{k_i}}{p_{t'_i}^{k'_i}}\right) \times \left(\prod_{i=1}^n \frac{D_{c(k'_i)}(t'_i)}{D_{c(k_i)}(t_i)}\right)$, the unobserved budget multipliers cancel ($\prod_{i=1}^n \frac{\lambda^{k_i}}{\lambda^{k'_i}} = 1$) -- a necessary condition to construct an empirical RP test.} 
%may need to shorten, eg hiding formula.
%$\frac{u'(x_{\text{left}})}{u'(x_{\text{right}})} = \left(\frac{\lambda_{\text{left}}}{\lambda_{\text{right}}}\right) \left(\frac{p_{\text{left}}}{p_{\text{right}}}\right) \left(\frac{D_{\text{right}}}{D_{\text{left}}}\right)$$
Furthermore, the sequence satisfies the \emph{downward-sloping demand property} if larger quantities on the left ($x_{t_i}^{k_i} \ge x_{t'_i}^{k'_i}$ for all $i$) imply a product of price ratios of at most one: $\prod_{i=1}^n (p_{t_i}^{k_i} / p_{t'_i}^{k'_i}) \le 1$.
\footnote{Intuitively, if a utility-maximizing consumer chooses larger quantities on the left ($x_{t_i}^{k_i} \ge x_{t'_i}^{k'_i}$), those goods must have been relatively cheaper, so the product of their price ratios cannot exceed one.}
Multiplying the first-order conditions along a balanced sequence bounds the product of price ratios:
{\setlength{\abovedisplayskip}{5pt} \setlength{\belowdisplayskip}{5pt}
\[
\prod_{i=1}^n \frac{p_{t_i}^{k_i}}{p_{t'_i}^{k'_i}} 
\;\;\le\;\; 
\delta^{\left[ \sum_i t_i - \sum_i t'_i \right]} 
\;\;\cdot\;\; 
\prod_c \left[ \beta(c) \right]^{\left( n_c^+ - n_c^- \right)}
\]}
where $n_c^+$ and $n_c^-$ count future-dated entries ($t > 0$) of cell $c$ on the left and right, respectively. Two counting conditions eliminate these unobserved parameters: $\sum_i t_i \ge \sum_i t'_i$ neutralizes $\delta$, while requiring $n_c^+ = n_c^-$ \emph{for each cell separately} neutralizes every $\beta(c)$ simultaneously. The requirement that every balanced sequence satisfying both conditions exhibits the downward-sloping demand property constitutes the axiom I call \textbf{SAR-SQH}.
%Strong Axiom of Revealed State-dependent Quasi-Hyperbolic Discounting

% Poiché sappiamo che $\delta \le 1$, l'esponente negativo ci garantisce matematicamente che il lato destro della disuguaglianza è sempre $\le 1$, indipendentemente da quale sia il vero valore sconosciuto di $\delta$. Non abbiamo bisogno di stimare $\delta$: il test empirico richiede semplicemente che il prodotto dei prezzi sia $\le 1$.


\begin{proposition}[Characterization and nesting]
\label{prop:sqh}
The data are SQH-rationalized if and only if they satisfy SAR-SQH. Replacing the per-cell
count by the pooled count $\sum_c n^+_c=\sum_c n^-_c$ gives \citeauthor{eis2020}'s axiom
for a single $\beta$; deleting it gives the axiom for $\beta=1$. Hence
$\text{EDU}\subset\text{QHD}\subset\text{SQH}$, and the empirical content of
state-dependence is exactly the set of datasets satisfying SAR-SQH and violating SAR-QHD.
\end{proposition}

\noindent
The proof (Online Appendix~A) follows \citet{eis2020}, extended to operate cell by cell.
%The proof (Online Appendix~A) follows \citeat{eis2020} cell by cell; feasibility is decided by a linearized Afriat system whose difference constraints reduce to a negative-cycle check, and the implementation is certified on synthetic agents (it returns $\{0.80\}$ for a present-biased agent, $\{1.00\}$ for an exponential one) and reproduces all six of \citeauthor{eis2020}'s published pass rates on their two panels to three decimals. The last part of the proposition is the useful one: a balanced sequence that is not cell-balanced is informative about the trait model and silent about the state-dependent one.
\fi


\noindent
The rest of this section establishes the revealed-preference implications of state-dependent discounting, evaluating intertemporal choice consistency without parametric assumptions on utility.
Section~\ref{subsec:RPtheory2} characterizes choices rationalizable by state-dependent quasi-hyperbolic discounting, nesting the fixed-trait model and exponential discounting as special cases. 
Section~\ref{subsec:rungladder} arranges these specifications into a restriction ladder, showing that an unconstrained $\beta(\cdot)$ carries no falsifiable implications: testability hinges on how $\beta$ varies, rather than whether it does.


\clearpage
%\subsection{A revealed-preference characterization with a state-dependent $\beta$}
\subsection{Revealed Preference: Characterizing State-Dependent Present Bias}
\label{subsec:RPtheory2}


\citet{eis2020} characterize the revealed-preference (RP) restrictions of quasi-hyperbolic discounting (QHD) with a single PB factor $\bar\beta$, assuming only an increasing concave utility \citep{afriat1967, varian1982}. %, browning1989}. 
I extend their characterization to a factor $\beta(c)$ that varies across exogenously assigned state cells $c\in\{1,\dots,C\}$. In decision $k$, in cell $c(k)$, the agent chooses an affordable stream $x^k=(x^k_0,\dots,x^k_T)$ at prices $p^k=(p^k_0,\dots,p^k_T)$, where $t$ is the delay from the decision.\footnote{As in \citet{eis2020}, $x^k$ is committed to when deciding: a date is $t=0$ on the day, $t>0$ in advance.}

\begin{definition}[State-dependent QHD]
The data are \emph{SQH-rationalized} if there exist an increasing concave $u$, a $\delta\in(0,1]$ and a function $\beta:\{1,\dots,C\}\to\R_{++}$ such that every $x^k$ maximizes $\sum_{t} D_{c(k)}(t)\,u(x_t)$ over its budget, with $D_c(0)=1$ and $D_c(t)=\beta(c)\,\delta^{t}$ for $t>0$. QHD is the case $\beta(c)\equiv\bar\beta$; exponential discounting (EDU) is $\bar\beta=1$.
\end{definition}

%\noindent At an optimum, discounted marginal utility equals price times the multiplier $\lambda^k>0$ of the decision,
\noindent At an optimum, discounted marginal utility equals price multiplied by the multiplier $\lambda^k>0$,\footnote{For exposition, $u$ is differentiable and choices interior. A corner turns \eqref{eq:focRP2} into an inequality that preserves every step below; supergradients replace marginal utilities if $u$ has kinks, which is also why ties impose no restriction.}
\begin{equation}\label{eq:focRP2}
D_{c(k)}(t)\,u'(x^k_t)=\lambda^k p^k_t
\qquad\Longleftrightarrow\qquad
u'(x^k_t)=\lambda^k p^k_t\,/\,D_{c(k)}(t) ,
\end{equation}
so prices reveal marginal utilities up to three unknowns: $\lambda^k$, $\delta$ and the $\beta$ of the decision's cell. By concavity, a larger quantity cannot carry a larger marginal utility. Along a sequence of $n$ pairs of dated consumption entries, $(x_{t_i}^{k_i}, x_{t'_i}^{k'_i})_{i=1}^n$, each with a strictly larger quantity on the left than on the right, every ratio of marginal utilities is thus at most one, and so is their product:
\begin{equation}\label{eq:chainRP2}
1\;\ge\;\prod_{i=1}^{n}\frac{u'(x^{k_i}_{t_i})}{u'(x^{k'_i}_{t'_i})}
\;=\;
\underbrace{\prod_{i=1}^{n}\frac{\lambda^{k_i}}{\lambda^{k'_i}}}_{\text{multipliers}}
\cdot
\underbrace{\prod_{i=1}^{n}\vphantom{\frac{\lambda^{k_i}}{\lambda^{k'_i}}}\rho_i}_{\text{prices}}
\cdot
\underbrace{\vphantom{\prod_{i=1}^{n}\frac{\lambda^{k_i}}{\lambda^{k'_i}}}\delta^{-\tau}}_{\text{impatience}}
\cdot
\underbrace{\prod_{c=1}^{C}\vphantom{\frac{\lambda^{k_i}}{\lambda^{k'_i}}}\beta(c)^{\,n^-_c-n^+_c}}_{\text{present bias}},
\end{equation}
where $\rho_i$ is the left-to-right price ratio of pair $i$, $\tau\equiv\sum_i t_i-\sum_i t'_i$, and $n^+_c$ ($n^-_c$) counts the future-dated entries of cell $c$, which alone carry $\beta(c)$, on the left (right). The multipliers cancel if the sequence is \emph{balanced}, that is, if each decision appears equally often on the left and the right. The impatience factor is at least one if (i) $\tau\ge0$, since $\delta\le1$; and the PB factor equals one if (ii) per-cell balance holds: $n^+_c=n^-_c$ in every cell. %, an equality because $\beta(c)$ may exceed one. 
Then \eqref{eq:chainRP2} reduces to $\prod_i\rho_i\le1$: the larger quantities must have been relatively cheaper -- the \emph{downward-sloping demand property}.
The requirement that every balanced sequence satisfying (i) and (ii) exhibit this property defines the axiom \textbf{SAR-SQH}.
%Strong Axiom of Revealed State-dependent Quasi-Hyperbolic Discounting


\begin{proposition}[Characterization and nesting]
\label{prop:sqh2}
The data are SQH-rationalized if and only if they satisfy \emph{SAR-SQH}: every balanced sequence of pairs with $x^{k_i}_{t_i}>x^{k'_i}_{t'_i}$ for all $i$ that satisfies (i) and (ii) has $\prod_i\rho_i\le1$. With (ii) replaced by the pooled count $\sum_c n^+_c=\sum_c n^-_c$ this is the QHD axiom of \citet{eis2020}; without (ii), their EDU axiom. Hence $\text{EDU}\subset\text{QHD}\subset\text{SQH}$.
\end{proposition}

\noindent
The proof (Online Appendix~\ref{app:proofs}) follows \citet{eis2020}, extended to operate cell by cell.

\noindent\emph{Intuition.} The test looks in the data for balanced sequences meeting (i) and (ii), on which no values of $\lambda^k$, $\delta$ or $\beta(c)$ can offset the prices in \eqref{eq:chainRP2}. A single sequence with larger quantities on the left despite a relative-price product greater than one rejects SQH: no concave $u$, $\delta\le1$ and cell-specific $\beta(\cdot)$ fit those choices.
The per-cell count is a new restriction relative to \citet{eis2020}: it excludes sequences whose bound involves a ratio $\beta(c)/\beta(c')$. For a decision $k$ in cell~1 and $k'$ in cell~2, the sequence $\big((x^{k}_{1},x^{k'}_{1}),(x^{k'}_{0},x^{k}_{0})\big)$ has $\prod_i\rho_i=(p^k_1/p^k_0)/(p^{k'}_1/p^{k'}_0)$. A value above one rejects QHD, but under SQH it only implies $\beta(1)/\beta(2)\geq\prod_i\rho_i$. Data of this kind, which satisfy SAR-SQH but violate SAR-QHD, are the empirical content of state dependence.


\clearpage




%\subsection{The restriction ladder, and why its top rung is empty}
\subsection{The Ladder of Restrictions}
\label{subsec:rungladder}


%\textcolor{red}{CAN ADD AS A FOOTNOTE: Throughout, $X$ denotes whatever a test conditions on — the state at the moment of choice, but also the good, the sign--timing structure, the stake or the horizon. ``State-dependence'' in Rungs 0--2 means dependence on that conditioning variable, of which the state at the moment of choice is the leading case and the one the paper assigns.}

Proposition~\ref{prop:sqh2} defines a ladder of restrictions of increasing stringency (Figure~\ref{fig:ladderdiag}). \textbf{Rung 0} (unrestricted $\beta$) allows each state cell an unconstrained parameter $\beta(c)$. \textbf{Rung 1} (state-\emph{independent} $\beta$) permits $\beta$ to vary stochastically across periods and choices, provided its distribution is independent of the state. % (the \emph{random-$\beta$ null}). 
\textbf{Rung 2} (monotone $\beta(\cdot)$) enforces a monotonic ordering across state cells, \textbf{Rung 3} (constant $\beta$) recovers the standard trait model of QHD, and \textbf{Rung 4} ($\beta=1$) reduces to EDU.


\begin{figure}[H]\centering
\caption{The restriction ladder, rung by rung, and the corresponding balance conditions for the test weight.}
\includegraphics[width=.99\linewidth]{d04_ladder.pdf}
\label{fig:ladderdiag}
\end{figure}


\begin{proposition}[Absence of empirical restrictions on Rung 0]
\label{prop:rung0}
With the state discretized into $C$ cells and one cell-specific wedge statistic $m_c$ observed per cell,
Rung~0 imposes no restriction on $(m_1,\dots,m_C)$: each $m_c$ is a strictly monotone
function of $\beta_c$ at fixed $(u,\delta)$, so any observed vector is exactly rationalized
by some $(\beta_1,\dots,\beta_C)$, leaving zero degrees of freedom for empirical testing.
\end{proposition}

\noindent
One might be tempted to treat \textit{state-dependent PB} as the main hypothesis and a \textit{fixed trait} as the alternative. Proposition~\ref{prop:rung0} establishes that unconstrained state dependence is not a testable hypothesis at all: empirical content resides entirely in the structural restrictions placed on how $\beta$ varies across states. Thus, evaluating state dependence requires specifying which structural restriction on parameter variation -- such as those in the ladder -- is actually being tested.

%The deterministic characterization itself, exact as it is, cannot carry the empirical comparison: it needs a finely sampled budget line, interior choices, both vantages and the same person in two states, and no public dataset satisfies all four (the three-good panel offers six bundles per budget; the money panel is $86\%$ corners). What the characterization contributes to every statistical test in the paper is the next section's reading of its counting condition.





% =====================================================================================
%\section{Which statistic answers a state question}
\section{Econometric Framework for Testing State Dependence}
\label{sec:stat}


Section~\ref{sec:object} characterizes state dependence for decisions measured without error. When choices are observed with noise, identifying state dependence requires specifying appropriate test statistics. This section shows that the validity of such tests is governed by the same counting condition that separates SAR-SQH from SAR-QHD, derives the class of valid test statistics, shows why the conventional level statistic cannot answer a state question, defines the PB wedge underlying every empirical test in the paper, and sets out the measurement-error model used for individual factors.


\subsection{Valid Test Statistics and Recentering}
\label{sec:balance}


From an empirical perspective, the condition distinguishing SAR-SQH from SAR-QHD reduces to a statement about \emph{weights}. A balanced sequence assigns a signed weight to each future-dated entry, $+1$ on the left and $-1$ on the right, and condition~(ii) requires these weights to sum to zero within each state cell: $n^+_c-n^-_c=0$. 
The same logic applies to test statistics.
For a PB wedge $g_k$ of observation $k$ (the log-wedge of Section~\ref{sec:coord} or an RP violation index), the test statistics considered here take the form of a weighted average $T=\frac{1}{N}\sum_{k=1}^{N}z_k g_k$. 
When $g_k$ is specified as the log-wedge, it decomposes as $g_k = b \log \beta(s_k) + \eta_k$, where $b>0$ is the wage elasticity of planned effort defined in Section~\ref{sec:coord}, $s_k \in \{1,\dots,C\}$ denotes the state cell of observation $k$, and $\eta_k$ is an error term that may depend on predetermined covariates $w$ (such as the subject, vantage point, or wage grid). %(such as the subject, vantage point, daily structure, and wage grid).

Under any Rung-1 null, %(the random-$\beta$ null),
the distribution of $\beta$ does not depend on the state, so the cell-specific levels collapse to a common 
baseline $\bar g\equiv b\log\bar\beta$, with $\log\bar\beta\equiv\E[\log\beta_k]$, and each draw's deviation from it is absorbed in the error term. This yields $g_k=\bar{g}+\eta_k$ with $\mathbb{E}[\eta_k\mid s_k,w]=\mathbb{E}[\eta_k\mid w]$: under this null, the wedge consists of
a common PB level $\bar g$ plus an error whose conditional expectation, given predetermined variables, does not depend on the state.\footnote{Rung~1 requires the \emph{distribution} of $\beta$, given $w$, to be independent of the state; with an assigned state that leaves the other components of the error unaffected, this implies the mean condition. A weight that is a function of the state alone needs nothing more; a weight that depends on the whole assignment draw uses the independence of Proposition~\ref{prop:class}. In either case the error may vary freely with $w$, allowing subject-specific levels.} 
%The weighting vector $z$ can take three distinct forms.
Whether $T$ is a valid test of this null depends on how $z$ is balanced across the %state 
cells, and the weight can satisfy three balance conditions.

%%% OLD
%From an empirical perspective, the condition distinguishing SAR-SQH from SAR-QHD reduces to a statement about \emph{weights}. A balanced sequence assigns a signed weight to each future-dated entry, $+1$ on the left and $-1$ on the right, and condition~(ii) requires these weights to sum to zero within each state cell, $n^+_c-n^-_c=0$. 
%The same reading applies to test statistics. For a PB wedge $g_k$ of observation $k$ -- the log-wedge of Section~\ref{sec:coord}, $g_k = b \log \beta(c) + \eta_k$, or an RP violation index -- a test statistic is a weighted average $T=\frac{1}{N}\sum_{k=1}^{N}z_k g_k$. The state cell of observation $k$ is $s_k$, while $w$ collects every predetermined variable (such as the subject, the vantage, the day's structure and the wage grid). Under any Rung-1 null, the wedge decomposes as $g_k=\bar g+\eta_k$ with $\E[\eta_k\mid s_k,w]=\E[\eta_k\mid w]$: a common PB level $\bar g$ plus an error whose conditional expectation, given the predetermined variables, does not depend on the state.\footnote{Rung~1 requires the \emph{distribution} of $\beta$ to be independent of the state, which implies this mean condition. A weight that is a function of the state alone needs nothing more; a weight that depends on the whole assignment draw uses the independence assumed in Proposition~\ref{prop:class}. In either case the error may vary freely with $w$, so subject-specific levels are allowed.}The weighting vector $z$ can take three distinct forms.

\emph{Per-cell balance}, $\sum_{k:\,s_k=c}z_k=0$ for every $c$, is condition~(ii) applied to the weight. It removes any cell-specific level, so that $T$ is the same whatever the value of $\beta$ in each cell: a statistic balanced within cells cannot detect state dependence (Proposition~\ref{prop:rung0}).%
\footnote{Specifically, $\mathbb{E}[T] = \mathbb{E}\left[ (1/N) \sum_k z_k g_k \right] = (1/N) \mathbb{E}\left[ \sum_c \bigl( b \log \beta(c) \sum_{k \in c} z_k \bigr) \right] + (1/N) \mathbb{E}\left[ \sum_k z_k \eta_k \right]$. Per-cell balance sets $\sum_{k \in c} z_k = 0$ for each $c$, which zeroes the first term whatever the values of $\beta(c)$. The second term does not involve $\beta$, so $T$ has the same distribution under every hypothesis on Rung~0.}
% NOTA AF: La vera differenza empirica tra un modello a tratto costante (QHD) e un modello a dipendenza dallo stato (SQH) risiede nelle sequenze che NON rispettano la per-cell balance ($n_c^+ \neq n_c^-$ in qualche cella), ma rispettano solo la pooled balance ($\sum_c n_c^+ = \sum_c n_c^-$).
\emph{No balance}, $z_k\equiv1$, yields the level statistic. Its expectation under any Rung-1 null is the mean wedge, $\bar g+\E[\eta_k]$, which differs from zero whenever PB is present and the error has mean zero: the level tests the existence of PB, not its dependence on the state.
%%%
% Can call it: Pooled/Conditional/Recentered Balance
\emph{Conditional balance}, $\mathbb{E}[z_k \mid w] = 0$, centers the weight on the assignment process, %($g \mid w$), 
while allowing it to differ across cells. It removes the common baseline level $\bar{g}$ and every component that depends on $w$ alone (e.g., a %non-linear 
transformation bias, a fixed misprediction, a censoring artifact, the subject's mean wedge), ensuring $\mathbb{E}[T]=0$ under any Rung-1 null while preserving genuine state variation in $\beta(s_k)$. A statistic with conditional balance is therefore a valid test of Rung~1.% (Proposition~\ref{prop:class}).%
\footnote{Specifically, under any Rung-1 null where $g_k = \bar{g} + \eta_k$ with $\mathbb{E}[\eta_k \mid s_k, w] = \mathbb{E}[\eta_k \mid w]$, by law of iterated expectations $\mathbb{E}[T] = (1/N) \sum_k \mathbb{E}\bigl[ \mathbb{E}[z_k g_k \mid s_k, w] \bigr] = (1/N) \sum_k \mathbb{E}\bigl[ \mathbb{E}[z_k \mid w] (\bar{g} + \mathbb{E}[\eta_k \mid w]) \bigr] = 0$. Under the alternative, conditional balance does not force $\sum_{k\in c}z_k=0$ within cells, so the cell levels survive: with a binary state and the recentered weight $z_k=s_k-\mu_k$ defined below, the same steps give $\E[T]=b\,[\log\beta(1)-\log\beta(0)]\,\frac1N\sum_k\E[\Var(s_k\mid w)]$, the state effect scaled by the variance of the state that remains once $w$ is fixed.}
%\emph{Pooled balance}, $\E[z_k\mid w]=0$, centers the weight on the assignment process while allowing it to differ across cells. It removes the common level and every component that depends on $w$ alone -- a Jensen gap, a fixed misprediction, a censoring artefact, the subject's mean wedge -- but not a $\beta$ that moves with the cell, which is therefore what the statistic measures. 
%%%
Conditional balance is \emph{recentering} in the sense of \citet{borusyak2023}. The assigned state can be viewed as a formula instrument, $s_k=f_k(\xi;w)$: a %known 
function of an exogenous draw $\xi$ (in Section~\ref{sec:state}, the randomized order of a session's components) and of the predetermined structure $w$. Its expectation over the assignment process, $\mu_k=\E[f_k(\xi;w)\mid w]$, can be computed, and the recentered weight $z_k=f_k(\xi;w)-\mu_k$ satisfies conditional balance by construction.%\footnote{Recentering keeps only the variation of the state that remains once $w$ is fixed, a share $\rho=\E[\Var(s\mid w)]/\Var(s)$ of its variance.}
\footnote{Recentering isolates the state variation remaining once $w$ is fixed, a share $\rho=\mathbb{E}[\operatorname{Var}(s\mid w)]/\operatorname{Var}(s)$ of its variance.}

\begin{proposition}[The valid class for Rung~1]
\label{prop:class}
Under exogenous assignment, $\xi \perp \eta \mid w$:\\
(i) every recentered weight $z_k=h_k(\xi;w)-\E[h_k(\xi;w)\mid w]$ gives $\E[T]=0$ under every Rung-1 null, for any function $h_k$; (ii) conversely, for any weight that is a function of $(\xi,w)$ but is not recentered, some Rung-1 null yields $\E[T]\neq0$, unless further restrictions are placed on $\eta$; (iii) the level weight $z_k\equiv1$ gives $\E[T]=\bar g+\E[\eta_k]$, non-zero under every Rung-1 null with PB and a mean-zero error.
\end{proposition}

\noindent
Part~(i) follows by the law of iterated expectations given $\mathbb{E}[z_k \mid w]=0$, part~(ii) establishes the necessity of recentering, %\citep{borusyakhull2026}; 
and part~(iii) illustrates the failure of uncentered level statistics ($z_k \equiv 1$), which confound the presence of PB with its dependence on the state (Online Appendix~\ref{app:proofs}).
As a consequence, every state test used below belongs to the valid class: 
the recentered product moment $\mathbb{E}[(s_k-\mu_k) g_k]$, where $(s_k-\mu_k)$ measures the extent to which an observation experiences an unexpected state relative to the experimental schedule.%
\footnote{%this was in the following subsection, earlier
Monte Carlo simulations confirm that uncentered level weights severely overreject under state-independent nulls, whereas the recentered product moment holds size near $.05$ while rejecting in every replication against a genuine reaction function. % -- using the moment inequality test of \citet{arp2023}. 
More details in Online Appendix~\ref{app:mc}.} 
Specifically, when assignment has predetermined structure $w$, the appropriate center is the expected state given that structure, $\mu_k \equiv \mathbb{E}[s_k \mid w]$, computed as detailed in Section~\ref{sec:recentered}.%
\footnote{Each such moment is computed as a regression coefficient: 
%by the Frisch--Waugh--Lovell theorem \citep{frischwaugh1933,lovell1963}, 
since an OLS or IV coefficient is a weighted sum of the outcome \citep{frischwaugh1933,lovell1963} with the residualized regressor (or instrument) as the weight \citep{angristpischke2009}, the choice of fixed effects, controls and instruments in Section~\ref{sec:coord} is the choice of $z$.} 
Finally, the efficient member of the class is the recentered best predictor of the treatment, reweighted by the inverse conditional variance of the error \citep{borusyakhull2026}. %[Theorem~1]
%%%%%%%%%%
Figure~\ref{fig:ladderdiag} summarizes the balance condition that each rung requires. 
The paper takes the rungs to data in turn: Rung~1 in Section~\ref{sec:state}, where state independence is rejected for the assigned state as well as for the stake and proximity margins; Rung~2 in Section~\ref{sec:margins}; Rung~3, the trait model, in Section~\ref{sec:person}, where it is rejected; and Rung~4 in Section~\ref{sec:magn}.



%Three consequences structure the rest of the paper. First, every state test used below belongs to the valid class. The recentered product moment $\E[(s_k-\mu_k)\,g_k]$, the difference between cell means and the adjacent-increment Wald statistic, which stacks $C-1$ such contrasts, differ only in how the state is centered; when assignment has structure, the appropriate center is the expected state given that structure, computed in Section~\ref{sec:recentered}. Second, the efficient member of the class is known: it is the recentered best predictor of the treatment, reweighted by the inverse conditional variance of the error \citep[Theorem~1]{borusyakhull2026}. 
%The optional third step of their Algorithm~1, which residualizes on $w$ and reweights, is precisely what separates the paper's joint within-subject fit from the raw product moment. 
%Third, a design may leave nothing to recenter. The share of the state's variance that survives conditioning on the structure, $\rho=\E[\Var(s\mid w)]/\Var(s)$, is exactly the share that recentering retains, and when $\rho=0$ no valid test exists within the design, however the moments are arranged. The assigned state of Section~\ref{sec:state} has $\rho=0.81$, whereas the position of money in the three-good protocol of Section~\ref{sec:goods} has $\rho=0$.


% AF: tab:ladder kept here but not printed (replaced by fig:ladderdiag and the sentence above).
\iffalse
\begin{table}[H]\centering
\caption{The ladder as three balance conditions on the test weight: what each rung
restricts, the statistic that answers it, and the verdict. \textcolor{red}{would cancel now}}
\label{tab:ladder}
{\small
\begin{tabular}{>{\raggedright\arraybackslash}p{0.7cm}>{\raggedright\arraybackslash}p{3.0cm}>{\raggedright\arraybackslash}p{4.3cm}>{\raggedright\arraybackslash}p{6.1cm}}
\toprule
Rung & Restriction & Weight (balance); statistic & Where, and verdict \\
\midrule
0 & $\beta_t(X)$ free & per cell: nothing survives & Prop.~\ref{prop:rung0}; untestable \\
1 & $\beta\perp$ state: the random-$\beta$ null & conditional $=$ recentered functions of the assigned state; exact RI & \S\ref{sec:state}: \textbf{rejected} for the assigned state (RI $p=.09$ all, $.02$ keep list); rejected for stakes and proximity; bounded null for load and appetite \\
2 & $\beta$ monotone in the state & distance to the monotone cone & \S\ref{sec:margins}: \emph{decreasing} in the stake rejected; the fine shape not robust \\
3 & $\beta$ constant: the trait & within-person contrast across structures or goods; measurement-error model; ME-robust set & \S\ref{sec:person}: \textbf{rejected} --- a variable over effort, a constant over money; money the most present-biased good \\
4 & $\beta=1$ & none: the level, signed in every cell & \S\ref{sec:magn}: rejected on the effort panels, accepted on a consumption panel \\
\bottomrule
\end{tabular}}
\end{table}
\fi

%\subsection{What the level statistic does, in numbers}

%Part~(iii) of Proposition~\ref{prop:class} is immediate; its practical relevance lies in the magnitude. Across $500$ simulated populations of realistic size, the joint test that the wedge is positive in every cell rejects in every replication under three nulls in which $\beta$ does not depend on the state: a constant $\beta$, a $\beta$ redrawn independently each period, and a heterogeneous, serially dependent $\beta$. Under the same nulls, the recentered product moment holds size at $.054$, $.044$ and $.056$, while rejecting in every replication against a genuine reaction function.

%Among recentered statistics, the adjacent-increment Wald serves as the primary test. Its size is $.045$, and it has full power against a response that peaks in the middle cells, against which a linear slope has power $.043$; the price is some power against a monotone response ($.78$, against $.96$ for the slope). Paired-cell moments are valid under \citet{arp2023} but singular by construction, and Cox--Shi inference applied to them rejects a true null $83\%$ of the time. A one-sided test adds power when the direction of the effect is fixed in advance ($.92$ against $.86$), and a two-instrument basis, a slope and a hump, recovers power against a symmetric response that the product moment cannot detect ($.79$ against $.05$; Online Appendix~\ref{app:mc}).
%The random-$\beta$ null also explains why the sharpest results below come from margins where the state is \emph{assigned}: only there does every PB process drawn independently of the state imply exactly zero response.






%\subsection{The coordinate: scale-free, curvature-free, Jensen-free}
%\subsection{The coordinate: scale-free, curvature-free, transformation-robust}
\subsection{Transformation-Robust Measurement: The Log-Wedge Coordinate}
\label{sec:coord}

Estimates below are computed on a single coordinate $(\theta)$. Let $e$ denote the effort chosen, $\kappa(e)$ an increasing, strictly convex cost, $W$ the wage paid on a later date, and $\tilde\delta$ the long-run discount factor between the work date and the payment date. 
At the advance vantage both the cost and the reward lie in the future, so the objective is proportional to $\beta\,[-\kappa(e)+\tilde\delta\,We]$.
The resulting first-order condition, $\kappa'(e_{\rm plan})=\tilde\delta W$, in log-linear form becomes $\log\kappa'(e_{\rm plan})=\log\tilde\delta+\log W$, which does not involve $\beta$.
On the day of work the cost is immediate, the objective becomes $-\kappa(e)+\beta\tilde\delta\,We$, and $\log\kappa'(e_{\rm act})=\log\beta+\log\tilde\delta+\log W$. The difference between the two conditions is exactly $\log\beta$. Since effort rather than marginal cost is observed, I use the log-wedge $\theta\equiv\log e_{\rm act}-\log e_{\rm plan}$ and the wage elasticity of planned effort $b\equiv\partial\log e_{\rm plan}/\partial\log W$, which the advance-vantage FOC condition identifies as the reciprocal of the elasticity of marginal cost.%
\footnote{%
In the advance FOC, $\kappa'(e_{\rm plan}) = \tilde\delta W$ yields $\log W = -\log \tilde\delta + \log \kappa'(e_{\rm plan})$. Differentiating both sides wrt $\log W$ gives $1 = \frac{\partial \log \kappa'(e_{\rm plan})}{\partial \log e_{\rm plan}} \frac{\partial \log e_{\rm plan}}{\partial \log W} = \epsilon_{\kappa'} \cdot b$, where $\epsilon_{\kappa'} \equiv e \kappa''(e)/\kappa'(e)$ is the marginal cost elasticity. Thus, $b = 1/\epsilon_{\kappa'}$.} 


Then $\log\beta = \theta / b$ (exactly under an isoelastic cost, to a first-order approximation otherwise) and every reported $\beta$ is obtained as $\exp(\hat\theta / \hat b)$ from the within-subject regression%
\footnote{Since $\log\beta = \log \kappa'(e_{\rm act}) - \log \kappa'(e_{\rm plan})$, an isoelastic cost $\kappa(e) = e^\gamma/\gamma$ with $\gamma>1$ gives $(\gamma - 1)(\log e_{\rm act} - \log e_{\rm plan}) = \log \beta$, %\implies 
which implies $\theta = b \log \beta$, given $b \equiv \partial \log e_{\rm plan} / \partial \log W = 1/(\gamma - 1)$. Testing state dependence does not require isoelasticity: at a given wage the plan does not involve $\beta$, and since $\kappa'$ is strictly increasing the act is strictly increasing in $\beta$, so $\theta_c=\theta_{c'}$ if and only if $\beta_c=\beta_{c'}$ whenever the cells face the same wages, as they do in Section~\ref{sec:state}.}
\[
\log e_{ik} = a_i + b \log W_{ik} + \theta \, \mathbf{1}\{\text{today}\}_{ik} + \varepsilon_{ik} ,
\]
where $a_i$ denotes subject fixed effects %(absorbing individual time-invariant productivity and cost parameters), 
and $\varepsilon_{ik}$ is an idiosyncratic error term. 
With $\mathbf{1}\{\text{today}\} = 1$ for choices made at the today vantage and $0$ at the advance vantage, $\theta$ captures the average log-wedge between plan and action. Whenever state dependence is tested, the state-cell indicators $\mathbf 1\{s_k=c\}$ enter both on their own and interacted with $\mathbf 1\{\text{today}\}$, so that $\theta_c$ varies across cells while the plan itself may also shift with the state.
In this coordinate, the log-wedge defined in Section~\ref{sec:balance} is decision-specific: for a today-decision $k$, $g_k = \log e_k - (\hat{a}_i + \hat{b}\log W_k)$ measures the log-deviation of the observed action from the planned allocation for the same subject at the same wage (and, when testing state dependence, within the same state). The cell parameter $\theta_c$ is the sample mean of $g_k$ across today-decisions in cell $c$, and the test of state dependence evaluates the contrast $\theta_c - \theta_{c'}$: a weighted average of $g_k$ with opposite-signed weights across the two cells, delivered directly by the interaction coefficient in the panel regression.%
% Un piccolo esempio numericoNella cella hot ($c$) hai 2 decisioni: $g_1 = 4$, $g_2 = 6$. La media è $\theta_c = 5$.Nella cella cold ($c'$) hai 2 decisioni: $g_3 = 2$, $g_4 = 2$. La media è $\theta_{c'} = 2$.Il contrasto è $\theta_c - \theta_{c'} = 5 - 2 = 3$.Se lo vedi come media pesata: $$g_1\left(+\frac{1}{2}\right) + g_2\left(+\frac{1}{2}\right) + g_3\left(-\frac{1}{2}\right) + g_4\left(-\frac{1}{2}\right) = 4(0.5) + 6(0.5) - 2(0.5) - 2(0.5) = 2 + 3 - 1 - 1 = \mathbf{3}$$Se stimi la regressione con l'interazione, il coefficiente $\hat{\beta}_{\text{interazione}}$ restituirà esattamente $3$.
\footnote{%
Subject fixed effects recenter these weights within each respondent, achieving conditional balance at the subject level; when state assignment follows a structured schedule (as in Section~\ref{sec:state}), the state and its interaction are instead instrumented by their recentered versions, $s_k - \mu_k$ and $\mathbf{1}\{\text{today}\}(s_k - \mu_k)$, or $\mu_k$ is included directly as a control -- the two asymptotically equivalent formulations of \citet{borusyak2023}.} 



%%% new par
\noindent
An advantage of the coordinate $\theta$ is that it is estimated directly from observed choices, 
avoiding functional-form assumptions on the cost structure and preserving the nonparametric nature of the test.
%which avoids functional-form assumptions on the cost structure and preserves non-parametricity at the testing stage. 
%Tests run in $\theta$ are linear in the data and never divide by an estimated elasticity.
Inference in $\theta$ is linear in the data and never divides by an estimated elasticity.
Specifically, $\theta$ satisfies three distinct econometric properties.
First, it is \emph{scale-free}: the level difference $e_{\rm act} - e_{\rm plan}$, reported in parts of the literature, grows mechanically with the planned quantity at a fixed $\beta$, whereas the log-wedge does not.
Second, at a common $b$, it is \emph{curvature-free}: since $\theta = b \log \beta$ combines preference with cost curvature, $\theta$ compares $\beta$ across cells whenever $b$ is shared.%
\footnote{This holds across the state cells of Section~\ref{sec:state} (where it is tested and not rejected), but fails across the goods of Section~\ref{sec:goods}, where comparisons are instead made in $\theta / b$.} 
Third, it is \emph{transformation-robust}: tests in $\theta$ are unaffected by the convex map to $\hat\beta=\exp(\hat\theta/\hat b)$, whose expectation exceeds $\beta$ by a factor of approximately $1+\tfrac12\Var(\hat\theta/\hat b)$, a term amplified by $1/b^{2}$. This convexity bias is negligible in aggregate ($+0.07\%$ on the effort panel) but substantial at the individual level. On the data of \citet{ans2015}, the mean of unrestricted per-subject effort $\hat\beta$'s is $1.38$, against $0.92$ once each estimate is corrected: an inflation of $49.9\%$. Imposing their common-curvature restriction ($b_i=b$) reproduces their published mean ($0.916$). %($0.916$, against their $0.91$). 
Crucially, allowing fully heterogeneous per-subject curvature ($b_i$) and correcting for this non-linear transformation bias yields virtually the same result ($0.92$), showing that the finding is robust to relaxing common curvature once the convexity bias is accounted for.
Every statistical test is therefore conducted in $\theta$ (or $\theta/b$), and reported point estimates are converted to $\beta$ solely for ease of interpretation (\textit{two-step} procedure).


%\subsection{Individual factors: a model for the noise}
\subsection{Accounting for Sampling Noise: A Bivariate Error Model}
\label{sec:memodel}

Evaluating whether PB reflects a stable personal trait requires estimating individual parameters -- a setting where sampling instability poses a more severe challenge than non-linear transformation bias.  Specifically, given the limited number of decisions per respondent, estimating a subject-specific elasticity $\hat{b}_i$ introduces extreme noise into the denominator.\footnote{%
Evaluating the ratio $\hat\theta_i/\hat b_i$ subject by subject divides by an elasticity estimated on a small set of choices. %(around 20 per person). 
When a subject barely responds to prices ($\hat b_i \approx 0$), the ratio becomes heavy-tailed: in the effort task of \citet{ans2015}, a per-subject bootstrap reveals that $\hat\beta_i$ has no usable mean for $26$ of $86$ subjects.} 
Following \citet{ans2015}, I resolve this instability by estimating individual factors under a common domain-level elasticity $\hat{b}$ alongside subject-specific today-effects $\hat{\theta}_i$, yielding $\log\hat{\beta}_i = \hat{\theta}_i / \hat{b}$.
%Across subjects this reads a difference in today-effects as a difference in $\beta$: the restriction under which $\theta$ compares cells in Section~\ref{sec:coord}, now applied to people
Under this formulation, %given the common $\hat b$, 
each $\log\hat{\beta}_i$ is linear in the subject's own choices, avoiding heavy-tailed ratios and allowing the remaining sampling noise to be modeled directly. For a subject observed across two domains $d \in \{1,2\}$, the pair of estimated log present-bias factors is specified as:
\[
\hat r_i = r_i + \nu_i, \qquad r_i \sim N(\bar r, \Sigma), \qquad \nu_i \sim N(0, V_i), \qquad V_i = \mathrm{diag}(V_{i1}, V_{i2}),
\]
where $r_i = (\log\beta_{i1}, \log\beta_{i2})'$ contains the subject's true preference parameters, $\Sigma$ is the population variance-covariance matrix of interest, and $V_{id}$ is the sampling variance of the estimate in domain $d$.\footnote{%
Where microdata are available, $V_{id}$ is computed from the subject's residual variance and number of choices at each vantage, divided by $\hat b^2$; when relying on published %individual 
estimates, $V_{id}=(\mathrm{se}_{id}/\hat\beta_{id})^2$, the delta-method variance of $\log\hat\beta_{id}$.} %
Assuming sampling errors independent across domains and of the true preferences $r_i$, the marginal distribution of observed estimates is $\hat{r}_i \sim N(\bar{r}, \Sigma + V_i)$, and $(\bar{r}, \Sigma)$ can be estimated via maximum likelihood.\footnote{%
Through the inverse-variance weighting $(\Sigma + V_i)^{-1}$, imprecisely estimated subjects receive lower weight in the likelihood function, preserving sample size while preventing noisy estimates from distorting inference.} %
The diagonal elements of $\Sigma$ represent the true population variance of $\log\beta$ within each domain, purged of sampling noise. Isolating these diagonal entries directly addresses the prerequisite that any cross-domain correlation presupposes: whether genuine preference dispersion exists within a given domain in the first place. Section~\ref{sec:person} applies this model to effort and money in \citet{ans2015} and to the three goods of \citet{cheung2022}; Online Appendix~\ref{app:me} provides further details on estimation, tests, and robustness.

%AF, from Claude. "Per-subject bootstrap" (nella frase prima di quella che citi: "*a per-subject bootstrap finds that β̂_i has no usable mean for 26 of 86 subjects*") è un bootstrap fatto **dentro ogni singolo soggetto**, uno alla volta (`matched_v1._boot_inflation`).
%- **Cosa resta fermo:** il soggetto (si usano solo le sue ~20 decisioni sullo sforzo), la specificazione della sua regressione individuale (log-rapporto delle allocazioni su costante, indicatore "oggi" e log del tasso) e il suo stimatore θ̂_i/b̂_i.
%- **Cosa cambia:** quali decisioni entrano. Si ricampionano le sue decisioni con reinserimento, 399 volte. Ogni decisione porta con sé la sua scelta, il suo tasso e la sua vantage. Ogni volta si rifà la regressione e si ricalcola il rapporto θ*/b*.
%- **Cosa si guarda:** quanto la media di exp(θ*/b*) supera la stima exp(θ̂_i/b̂_i). Se il rapporto supera 5, la media è dominata dalla coda, cioè dalle estrazioni in cui b* è vicino a zero. Allora quel soggetto "non ha una media utilizzabile", e succede per 26 soggetti su 86.
%Non va confuso con il **"subject bootstrap"** delle tabelle, che ricampiona i *soggetti* interi e tiene fisse le decisioni di ciascuno.
%La frase che citi descrive proprio la risposta a questo problema: con un'unica elasticità b per dominio, stimata su tutti i soggetti, non si divide più per un b̂_i individuale instabile. Così θ̂_i/b̂ è lineare nelle scelte del soggetto e senza code pesanti.

% Notes from LM (not sure they're correct) in: "C:\Users\alefa\OneDrive 1TB A\OneDrive\MASTER BMELA - BA\Ph\2026\SOP_ideas\FINAL_WRITING\old\LMnoiseMEmodel.md"




% =====================================================================================
\section{Data}
\label{sec:data}

Table~\ref{tab:timing} summarizes the experimental designs of the three primary datasets supporting the empirical findings below.
In \citet{augenblick2019}, hereafter \textit{AR2019}, $100$ subjects allocate transcription tasks across seven sessions over six weeks under piece rates ranging from $1$ to $35$ cents, making choices for the same dated work between $0$ and $26$ days in advance alongside incentivized predictions ($11{,}405$ decisions in total). Each session features ten mandatory tasks, with the ordering of session components block-randomized across subjects.
In \citet{ans2015}, hereafter \textit{ANS2015}, the same $101$ subjects complete intertemporal allocation tasks over both real effort and money within separate, parallel protocols conducted over seven weeks.
In \citet{cheung2022}, hereafter \textit{CTW2022}, $697$ %high-school
students complete convex time budget elicitations across three reward domains -- money, healthy food, and unhealthy food -- under an identical price grid, time horizon, and delivery protocol.
Measured over money, PB estimates may confound the discount factor with other features of dated payments -- risk and trust in the experimenter, transaction costs, and the possibility of borrowing or saving outside the experiment \citep[e.g.,][]{andreoni2012b,deansautmann2021}. %,cohen2020}. 
The three experimental setups neutralize these channels across the trade-offs analyzed in this paper.
\footnote{In AR2019 and in the effort task of ANS2015, all work is paid on one fixed date, the same for the plan and the act, so payment risk, trust, transaction costs and liquidity enter both choices alike and cancel from the wedge; effort, unlike money, also cannot be moved across dates by outside borrowing or lending. In CTW2022 the three rewards share dates, delivery and presentation, and Online Appendix~\ref{app:goods} shows that arbitrage, trust and the scale of delivery do not account for the ordering of goods. Moreover, any remaining factor that does not depend on the randomized order is balanced across the assigned states, and recentering removes every component that depends on the predetermined structure alone (Section~\ref{sec:balance}).}
Nineteen %Twenty-eight (in the long paper version?)
further datasets provide supplementary results in subsequent sections %throughout the paper 
and are documented comprehensively in Online Appendix~\ref{app:data}.

 
%Everything is public, and nothing was collected. One distinction stays visible throughout: what is \emph{reproduced} from the source paper, to anchor against it, and what is \emph{new}; every reproduced benchmark is attributed where it is used.

\begin{table}[H]\centering
\caption{The three designs.}%: when the decision is made, when the cost and the reward land, and what varies.}
\label{tab:timing}
\includegraphics[width=0.99\linewidth]{d06_designs.pdf}
% tab:timing is drawn in figs/diagrams/d06_designs.tex (TikZ; pdflatex d06_designs); the tabular below is the earlier version, kept but not printed
\iffalse
{\small
\begin{tabular}{>{\raggedright\arraybackslash}p{2.3cm}>{\raggedright\arraybackslash}p{3.2cm}>{\raggedright\arraybackslash}p{2.0cm}>{\raggedright\arraybackslash}p{2.4cm}>{\raggedright\arraybackslash}p{0.9cm}>{\raggedright\arraybackslash}p{2.4cm}}
\toprule
Dataset & Decision made & Cost borne & Reward & Cell & What varies \\
\midrule
AR2019%\citet{augenblick2019}
 & $0$--$20$ days before the work date, and on it & the work date & one fixed date, 7 weeks out & C--C & \textbf{just completed 10 mandatory tasks (randomized order)}; wage \\
ANS2015%\citet{ans2015}
 & week 1 (advance) and week 2 (on the day) & the work dates (effort); --- (money) & fixed end date (effort); two dated payments (money) & C--C; R--R & the structure and the good (same people) \\
 CTW2022%\citet{cheung2022}
 & two sessions, three weeks apart & --- & dated money or food, delivered identically & R--R & the good; vantage \\
\bottomrule
\end{tabular}}
\fi
\end{table}

\bigskip

% =====================================================================================
%\section{Evidence I: the state at the moment of choice}
\section{State Dependence}
\label{sec:state}


%there was a \textcolor{red}{very nice fig (\ref{fig:betaeff}) might fit the start/end of results on state-dep; but check/update it to reflect actual final results...}
\iffalse
\begin{figure}[H]\centering
\includegraphics[width=.99\linewidth]{figs/diagrams/d08_betaeff_margins.pdf}
\caption{The object the paper tests for, margin by margin. In every panel the dashed gray
line is the trait model (one number) and the colored profile is the estimated response on
that margin, with the verdict of the ladder test of Section~\ref{sec:ladderresults}. The
proximity margin is drawn apart, because it is a profile across \emph{future}-dated
decisions: a test of the quasi-hyperbolic cliff at ``now'' rather than a value of $\beta$,
which is the distinction the previous subsection draws.}
\label{fig:betaeff}
\end{figure}
\fi

\subsection{Assigned Decision States Move $\beta$}
\label{sec:matched}

A recurring challenge in the literature on visceral influences is that a ``state'' is often confounded with person-level attributes. Hungry, fatigued, or intoxicated individuals may differ along many unobserved dimensions, so that comparing across states inadvertently conflates state dependence with individual heterogeneity. Section~\ref{sec:stat} sharpens the test of state dependence into two requirements: only an assigned state can reject the random-$\beta$ null, and only a recentered statistic can isolate the effect of the state. The experimental design of \citet{augenblick2019} allows both criteria to be met, yet had not previously been exploited for this purpose. %
Crossing the assigned state (deciding before versus after the mandatory task block) with the decision vantage (a future-dated plan versus a today-action) generates a $2\times2$ factorial design on a common wage grid within the same subjects: a plan formed in a cold state, a plan formed in a hot state, an action taken in a cold state, and an action taken in a hot state. Estimating the two-step specification with state and vantage fully interacted within subject, evaluated via a subject bootstrap, yields the following results.\footnote{%
Throughout, ``subject bootstrap'' refers to resampling whole subjects with replacement, re-estimating the full two-step procedure across $B$ resamples, and deriving confidence intervals and two-sided $p$-values directly from the resulting empirical distribution; this accounts for arbitrary dependence across choices made by the same individual, within whom the state is assigned, session by session.}


\begin{table}[H]
\centering
\caption{Present bias is state-dependent: it intensifies in the hot condition.}
\label{tab:matched}

\begin{threeparttable}
{\footnotesize
\renewcommand{\arraystretch}{1.12}
\begin{tabular}{@{}l@{\hspace{16pt}}r@{\hspace{6pt}}c@{\hspace{26pt}}r@{\hspace{6pt}}c@{}}
\toprule
 & \multicolumn{2}{c@{\hspace{26pt}}}{\textbf{All $\mathbf{100}$ subjects}} & \multicolumn{2}{c@{}}{AR2019 subset ($N=72$)} \\
\cmidrule(lr){2-3}\cmidrule(l){4-5}
wage elasticity $b$
    & $0.549$ & $[\phantom{-}0.489,\phantom{-}0.615]$
    & $0.654$ & $[\phantom{-}0.602,\phantom{-}0.709]$ \\
hot state on a \emph{future}-dated choice %(the plan)
    & $+0.021$ & $[-0.020,+0.058]$
    & $+0.006$ & $[-0.039,+0.046]$ \\
\addlinespace[3pt]
$\beta$, state held \emph{cold}
    & $0.882$ & $[\phantom{-}0.802,\phantom{-}0.966]$
    & $0.904$ & $[\phantom{-}0.833,\phantom{-}0.981]$ \\
$\beta$, state held \emph{hot}
    & $0.736$ & $[\phantom{-}0.649,\phantom{-}0.822]$
    & $0.748$ & $[\phantom{-}0.654,\phantom{-}0.841]$ \\
\addlinespace[3pt]
$\Delta\log\beta$ across the assigned state
    & $\mathbf{-0.180}$ & $[-0.310,-0.062]$
    & $-0.191$ & $[-0.326,-0.060]$ \\
\quad $p$-value
    & $.008$ &
    & $<.004$ & \\
\bottomrule
\end{tabular}
}

\begin{tablenotes}[flushleft]
\footnotesize
\item \textit{Note:} Within-subject estimates; subject bootstrap ($B=499$); $95\%$ intervals. Main results use the full sample; results on the original paper's estimation subsample are shown alongside.
\end{tablenotes}

\end{threeparttable}
\end{table}


\begin{figure}[H]\centering
\caption{\footnotesize%
The hot state cuts immediate effort while leaving future-work decisions unchanged; pooling the states, as AR2019 do, returns a state-weighted average.}
\includegraphics[width=.99\linewidth]{fig_matched_state.pdf}
%The assigned state, priced in $\beta$ units, all $100$ subjects. \emph{Left:} the hot state leaves choices about \emph{future} work untouched ($+0.021$) and cuts choices about work \emph{now} by $7.8$ log points (the total effect; $95\%$ subject-bootstrap intervals). \emph{Right:} with the state held fixed on each side, $\beta$ is $0.88$ cold and $0.74$ hot in the same people; the conventional two-step that pools the states returns $0.80$, a state-weighted average of the two.}
\label{fig:matched}
\end{figure}
% fig:matched is generated by writing_sample/TIME_RATIO/code/final/make_figs_v5.py (python make_figs_v5.py matched)

\clearpage
\noindent
The pooled $\beta$ for this panel, like the published one, is a state-weighted average (Figure~\ref{fig:matched}): a study whose observation protocol sampled hot and cold states in different proportions would have inferred a different parameter. Because the state is exogenously assigned within subject, this finding rejects not merely a constant $\beta$, but the entire class of state-independent present-bias processes (Rung~1 Wald test of state equality: $\chi^2_1 = 9.0$, $p = .003$ on AR2019's $72$-subject subsample). 
While AR2019 restrict their main analysis to the $72$ subjects for whom structural estimation succeeds, my two-step procedure retains all 100 subjects with a common domain elasticity $\hat{b}$, leaving the core finding intact ($\Delta\log\beta = -0.180$, $p = .008$).
%Importantly, while the original study restricted its individual analysis to those $72$ subjects due to convergence failures in subject-by-subject MLE, estimating a common domain elasticity $\hat{b}$ within a two-step procedure avoids heavy-tailed denominators and preserves all $100$ subjects without sample selection, leaving the core finding intact ($\Delta\log\beta = -0.180$, $p = .008$).
Under minimal functional-form assumptions, the direction of this finding is corroborated: matching each today-choice directly into the advance schedule bounds $\beta$ nonparametrically between wage grid points. Computed separately by state, the cold-state bracket lies above the hot-state bracket at both wage binnings, meeting only at the boundary (Online Appendix~\ref{app:robust}).
% Computed separately by state on the $96$ subjects observed in both states, the cold-state bracket lies above the hot-state bracket at both usable binnings, the two touching at a single point (Online Appendix~\ref{app:robust}); once sampling error is allowed for, the brackets overlap, so this check corroborates the direction without establishing it.
Finally, state assignment is orthogonal to the wage grid ($16.2$ cents in cold sessions and $16.3$ in hot sessions), confirming that state dependence is not driven by price composition.



%\subsection{What is randomized: exposure, recentering, randomization inference}
\subsection{Unequal Exposure: Design Recentering and Randomization Inference}
\label{sec:recentered}

% AF/LM: uno strato $w$ è definito dall'incrocio tra il giorno di calendario dell'esperimento e la composizione del menu (day $\times$ composition). Usare $(\text{day} \times \text{composition})$ rende lo strato più "fine" e controllato. In questo modo, quando sottrai $\mu_k$, pulisci l'indicatore di stato $s_k$ sia dallo sbilanciamento strutturale dei blocchi decisionali sia da qualsiasi trend o shock di calendario, rendendo lo strumento recenterato $\tilde{s}_k = s_k - \mu_k$ del tutto inattaccabile nei test di falsificazione.

The experiment randomizes the \emph{order} of a session's components, rather than assigning a state label directly to each choice. 
A session day comprises up to five decision sets (today's work, one or two sets of future work, and one or two sets of predictions) alongside one block of ten mandatory tasks; a decision is \textit{hot} if its set was administered after this block. Whether a decision falls into the hot state is a deterministic function of the random order draw $\xi$ and the day's predetermined menu structure $w$: a formula instrument, $s_k=f_k(\xi;w)$, characterized by non-random exposure to exogenous shocks. %\citep{borusyak2023}. 
Exposure to the hot state varies systematically across decision types: today's work is hot $46\%$ of the time, the first future set $60\%$, and the second future set $83\%$. To address this, the expected state $\mu_k \equiv \mathbb{E}[s_k \mid w]$ is computed as the share of subject-days within stratum $w$ (calendar day $\times$ menu composition) in which choice $k$ fell after the mandatory block ($s_k=1$).%
\footnote{%%%%%%%%%%
An unadjusted indicator $s_k \in \{0, 1\}$ would conflate state dependence with session timing: hot choices would over-represent future dates ($83\%$), whereas cold choices would disproportionately involve present work ($54\%$). In addition, shifts in available time horizons across experimental weeks introduce further mechanical variation. Subtracting $\mu_k$ constructs a recentered instrument $\tilde{s}_k \equiv s_k - \mu_k$ with zero conditional expectation ($\mathbb{E}[\tilde{s}_k \mid w] = 0$), eliminating the spurious over-representation of future choices in late-session slots.} 
This expected treatment $\mu_k$ explains $19\%$ of the variance of the raw state indicator $s_k$ and accounts for the correlations of $s_k$ with the decision vantage ($-0.22$), the distance to the work date ($+0.15$), and the %calendar 
decision day ($-0.12$).%
\footnote{%%%%%%%%%%
%Correlations between the raw indicator $s_k$ and design covariates 
These correlations do not reflect behavioral preferences, but rather an artifact of the experimental software's order-generation protocol that must be purged to isolate the exogenous variation in state assignment.} 
Subtracting $\mu_k$ drives all three correlations to essentially zero (Figure~\ref{fig:exposure}), redefining the instrument as the realized (exogenous) state assignment relative to expected exposure rather than a composite of random shocks and structural baseline variation.%
\footnote{%
Estimating $\mu_k$ relies on the assumption that realized orderings are exchangeable within each stratum -- an empirically testable condition. As reported below, the recentered instrument satisfies the joint falsification test ($p = .28$), while identifying the distance to the work date as the sole covariate requiring explicit control.}

Hence, the state effect is estimated via 2SLS using the recentered instrument (instrumenting the state indicator and its interaction with vantage by their recentered counterparts), via a control-function specification (adding $\mu_k$ and $\mu_k \times \text{today}$ as controls), and in reduced form. 
All three approaches, which \citet{borusyak2023} show to be valid under non-random exposure, yield virtually identical results.%($\Delta\log\beta$ of $-0.122$, $-0.122$ and $-0.121$).
%As established by \citet{borusyak2023}, all three approaches yield virtually identical results.%
\footnote{%
Under the exchangeability assumption, component orderings are equally likely within stratum $w$. Under the null (the state affects no decision), %observed 
choices remain fixed across counterfactual orderings; $p$-values are obtained by redrawing orderings $999$ times within strata and comparing the observed contrast against the %resulting 
randomization distribution.}
%AF on RI:
%Per 999 volte, chiedi al computer di rimescolare a caso l'ordine dei blocchi decisionali all'interno dello stesso strato $w$ (stesso giorno di calendario e stessa composizione).
%Ad ogni rimescolamento, assegni a ciascuna decisione un nuovo stato fittizio (hot o cold finto).
% Mantieni le scelte reali (congelate) dei soggetti. [under the null, assumiamo che lo stato viscerale (hot vs cold) non abbia alcun effetto sulle decisioni].
% Ricalcoli la stima (mettiamo eg sia $\Delta\log\beta$ tra state hot vs cold).
%%% queste sono stime fittizie under the null: ciò [eg qui $\Delta\log\beta$ tra state hot vs cold] che otterresti per puro caso se lo stato non contasse nulla e la sorte determinasse l'assegnazione dello stato.
% Ripetendo questo processo 999 volte, costruisci la distribuzione empirica della nulla (cioè la mappa di "quanto grande può questo effetto dello stato per puro caso")
% Infine prendi la stima vera (quella calcolata sui dati reali dell'esperimento, es. $-0.122$). Se, ad esempio, nei 999 rimescolamenti casuali solo 91 volte il caso ha prodotto un valore così estremo, il tuo $p$-value esatto di RI sarà: $\text{exact RI } p = \frac{91}{999} \approx \mathbf{.091}$.




\begin{samepage}%%% 3 figures

\begin{table}[H]\centering
\caption{The state effect, unadjusted and design-recentered.}
%Subject bootstrap for intervals and $p$; exact randomization inference (RI) over $999$ redraws of the orderings. The act row is the total effect of the state on today's decisions (plan effect plus interaction).}
\label{tab:recentered}
\begin{threeparttable}
{\footnotesize
\renewcommand{\arraystretch}{1.12}
\begin{tabular}{@{}l@{\hspace{12pt}}r@{\hspace{3pt}}l@{\hspace{12pt}}r@{\hspace{3pt}}l@{\hspace{22pt}}r@{\hspace{3pt}}l@{\hspace{12pt}}r@{\hspace{3pt}}l@{}}
\toprule
 & \multicolumn{4}{c@{\hspace{22pt}}}{\textbf{All $\mathbf{100}$ subjects}} & \multicolumn{4}{c@{}}{AR2019 subset ($N=72$)} \\
\cmidrule(lr){2-5}\cmidrule(l){6-9}
 & \multicolumn{2}{c@{\hspace{12pt}}}{unadjusted} & \multicolumn{2}{c@{\hspace{22pt}}}{recentered} & \multicolumn{2}{c@{\hspace{12pt}}}{unadjusted} & \multicolumn{2}{c@{}}{recentered} \\
\midrule
hot state on the plan & $+0.021$ & ($p=.27$) & $-0.010$ & ($p=.68$) & $+0.006$ & ($p=.79$) & $-0.030$ & ($p=.20$) \\
hot state on the act & $-0.078$ & ($p=.02$) & $-0.077$ & ($p=.01$) & $-0.119$ & ($p<.004$) & $-0.136$ & ($p<.004$) \\
\quad exact RI $p$ & \multicolumn{2}{c@{\hspace{12pt}}}{---} & \multicolumn{2}{c@{\hspace{22pt}}}{$\mathbf{.018}$} & \multicolumn{2}{c@{\hspace{12pt}}}{---} & \multicolumn{2}{c@{}}{$\mathbf{.001}$} \\
\addlinespace[3pt]
$\beta$ cold / hot & \multicolumn{2}{c@{\hspace{12pt}}}{$0.882$ / $0.736$} & \multicolumn{2}{c@{\hspace{22pt}}}{$0.847$ / $0.750$} & \multicolumn{2}{c@{\hspace{12pt}}}{$0.904$ / $0.748$} & \multicolumn{2}{c@{}}{$0.881$ / $0.750$} \\
$\Delta\log\beta$ & $-0.180$ & ($p=.008$) & $-0.122$ & ($p=.12$) & $-0.191$ & ($p<.004$) & $-0.161$ & ($p=.02$) \\
\quad exact RI $p$ & \multicolumn{2}{c@{\hspace{12pt}}}{---} & \multicolumn{2}{c@{\hspace{22pt}}}{$\mathbf{.091}$} & \multicolumn{2}{c@{\hspace{12pt}}}{---} & \multicolumn{2}{c@{}}{$\mathbf{.021}$} \\
\bottomrule
\end{tabular}}

\begin{tablenotes}[flushleft]
\footnotesize
\item \textit{Note:} Subject-bootstrap $p$ in parentheses ($B=499$). Exact randomization inference (RI) over $999$ redraws of the session orderings within strata. The act row is the total effect of the state on today's decisions.
\end{tablenotes}
\end{threeparttable}
\end{table}
% $\text{exact RI } p = \mathbf{.018}$.%%%%%%%%%%%%%
% INTUITION: Il calo di $\beta$ (o del lavoro) che osserviamo sotto stato hot è un vero effetto causale dello stato viscerale, oppure è solo un artefatto fortuito dovuto a come il computer ha rimescolato l'ordine delle schermate quel giorno?"
% Significato: Sotto l'ipotesi che lo stato non sposti alcuna decisione, la probabilità di osservare per pura fortuna dell'estrazione dell'ordine dei blocchi un calo del lavoro immediato pari a $-0.077$ log points (sui 100 soggetti) o $-0.135$ log points (sui 72 soggetti) è solo dell'1.8% nel campione totale e inferiore allo 0.1% nel campione di Augenblick & Rabin (2019). Questo conferma che l'effetto dello stato sull'azione immediata è altamente significativo e non è un artefatto dell'ordine delle task.

\begin{figure}[H]\centering
\caption{\footnotesize%
Purging non-random exposure by recentering the instrument.}
\includegraphics[width=.99\linewidth]{d05_exposure.pdf}
%\emph{Left:} two draws of the same subject-day; the same component is hot in one and cold in the other, with odds that differ by component. \emph{Right:} the act effect is invariant to the adjustment, the plan effect is null under both, and the $\beta$ gap follows the plan-side zero.}
\label{fig:exposure}
\end{figure}
% fig:exposure is drawn in figs/diagrams/d05_exposure.tex (TikZ; from that folder: pdflatex d05_exposure). 09-30 aligned version; the previous one is figs/diagrams/_backup_fig1_0924/d05_exposure_before_0930.tex

\begin{figure}[H]\centering
\caption{\footnotesize%
Recentering moves the plan-side zero, not the act: the hot state lowers today's allocation under both estimators and leaves the plan untouched, and $\beta$ is lower when hot under both.}
\includegraphics[width=.99\linewidth]{fig_recentred.pdf}
\label{fig:recentered}
\vspace{5pt}
\begin{minipage}{.99\linewidth}
\raggedright
\scriptsize
\textit{Note:} The state effect, unadjusted and design-recentered. 
\emph{(A):} the effect of the hot state on the plan and on the act, %(total)
log points, $95\%$ subject-bootstrap
intervals. %;open markers unadjusted, filled markers recentered. 
\emph{(B):} $\beta$ with the state held cold and held hot, under each estimator and in each sample, $95\%$ subject-bootstrap intervals.
\emph{(C), (D):} the exact randomization-inference null distributions on all $100$ subjects (the recentered coefficient under $999$ redraws of the session orderings within strata), the critical values, and $p$-values.
% -- the recentered coefficient under $999$ redraws of the session orderings within strata -- with the observed value (solid), the $5\%$ and $10\%$ critical values (dashed and dotted) and the rejection regions (shaded); the exact $p$ on the AR2019 subset ($N=72$) is reported alongside.
\end{minipage}
\end{figure}
% fig:recentered is generated by writing_sample/TIME_RATIO/code/final/make_figs_v5.py, function fig_recentered (python make_figs_v5.py recentered), from output/recentere_v1_summary.csv, recentere_v2_summary.csv and ri_draws_v4.csv. 10-01 wide version, drawn at its printed size for .99\linewidth; the 2x2 version is function fig_recentered_2x2 (fig_recentered_2x2.pdf); earlier figures in figs/diagrams/_backups/

% Claude on RI: RI (randomization inference), in short: we redraw the random session orderings 999 times, recompute the estimate each time, and see where the real estimate falls among those fakes. It relies only on the design's randomization, so it's the strongest check we have.
	%P%anel C (the act) is in our favour: the observed −0.077 is beyond the 5% critical value of ±0.064, p = .018.
	%%Panel D (the β gap on all 100) shows p = .091: beyond the 10% critical value, not the 5% one.
	%%I'd keep D. Panel B no longer shows the gap's interval, so dropping D as well would show only the favourable half.



\end{samepage}%%% 3 figures





\noindent
Table~\ref{tab:recentered} and Figure~\ref{fig:recentered} report the results. The effect of the hot state on the act is invariant to recentering: today's allocation falls by $7.8$ log points unadjusted and by $7.7$ recentered, with an exact randomization $p$ of $.018$ on all subjects and at most $.001$ on the AR2019 subset. The plan remains unaffected under both specifications ($+0.021$ and $-0.010$; randomization $p=.73$). 

What recentering changes is the counterfactual against which the act is measured. The raw $+0.02$ on future-dated decisions reflects exposure rather than the state, since the second future set refers to a later date and is hot more than 80\% of the time. 
Once this exposure is removed, the gap between the cold and hot $\beta$ narrows from $0.15$ to $0.10$ ($0.85$ against $0.75$) 
%, with a randomization $p$ of $.09$ on all subjects and $.02$ on the AR2019 subset:  -- driven not by a shift in the act effect, but by an adjustment in the plan estimate, which remains a statistical zero.%
-- driven by the adjustment in the plan estimate, which nonetheless remains a statistical zero.
\footnote{%
The recentered state is orthogonal to the wage drawn, the net distance to payment, the decision day and the expected treatment itself (randomization $p=.58$, $.53$, $.16$ and $.82$; joint $p=.28$). It fails only on the distance to the work date ($+0.085$, $p=.011$, against $+0.55$ for the raw indicator), which records the future dates a day's decision sets refer to; the distance is therefore added as a control, and moves neither estimate of the act by more than $0.001$. The stakes margin of Section~\ref{sec:margins} requires no such correction: the wage draw is flat across components, days, vantages and states, so its expected instrument is a constant.}
%Both estimates are recentered statistics in the sense of Section~\ref{sec:balance}, differing only in what the state is centered on: the unadjusted one is what the design was read as delivering, and the recentered one is the estimator that \citet{borusyak2023} show to be robust to non-random exposure.

% previous version of this paragraph (before 09-30), kept but not printed:
\iffalse
\noindent
\textbf{The effect of the state on the act is invariant to the adjustment}
(Figure~\ref{fig:recentered}): $7.8$ log
points unadjusted, $7.7$ recentered, exact randomization $p=.018$ on all subjects and below
$.001$ on the source sample; distance to the work date as a control moves neither by more
than $0.001$. \textbf{The plan is untouched under both}: $+0.021$ and $-0.010$,
randomization $p=.73$. \textbf{What moves is the counterfactual}: the raw $+0.02$ on
future decisions was exposure --- the second future set concerns a farther date and is hot
five times in six --- and once it is centered the $\beta$ gap is measured against a plan
effect of $-0.01$ rather than $+0.02$. The gap shrinks from $0.15$ to $0.10$ of $\beta$,
with randomization $p=.09$ on all subjects and $.02$ on the keep list, not because the act
moved but because the zero did. Both columns are recentered statistics in the sense of
Section~\ref{sec:balance}, differing in what the state is centered on; the paper leads with
the first because it is what the design was read as delivering, and reports the second
because it is the one \citeauthor{borusyak2023} show to be safe. Falsification licenses
it: the recentered state is orthogonal to the wage drawn, the net distance to payment, the
decision day and the expected treatment itself (randomization $p$ of $.58$, $.53$, $.16$,
$.82$; joint $p=.28$), and fails only on distance to the work date ($+0.085$, $p=.011$,
against $+0.55$ for the raw indicator) --- the trace of which future dates a day's sets
refer to, which is why distance enters as a control. The stakes margin below needed no
repair: the wage draw is flat across components, days, vantages and states, so its
expected instrument is a constant.
\fi




%\subsection{The state moves the act and leaves the plan alone}
\subsection{The State Shifts the Act but Leaves the Plan Unchanged}
\label{sec:vantage}

AR2019 randomize session order to examine projection bias, but focus on pooled effects and leave differences across decision vantages uninterpreted.
Their projection-bias model posits that agents extrapolate current marginal disutility to all subsequent choices, imposing a uniform shift across post-work decisions (present work, future work, and predictions). \textit{Cross-vantage equality} -- an identifying restriction of this structural parameter -- is tested and rejected in Table~\ref{tab:vant}.

\begin{table}[H]\centering
\caption{The cross-vantage test, unadjusted and design-recentered.}
\label{tab:vant}
\begin{threeparttable}
{\footnotesize
\renewcommand{\arraystretch}{1.12}%altezza colonne
\setlength{\tabcolsep}{10pt}%larghezza colonne
\begin{tabular}{lcc}
\toprule
effect of the hot state on & unadjusted & recentered \\
\midrule
\emph{present} decisions & $-0.141$ $[-0.222,-0.053]$ & $\mathbf{-0.160}$ $[-0.254,-0.067]$ \\
\emph{future} decisions & $+0.008$ $[-0.030,+0.048]$ & $-0.031$ $[-0.074,+0.015]$ \\
\emph{predictions} & $+0.015$ $[-0.029,+0.062]$ & $-0.013$ $[-0.069,+0.043]$ \\
present $-$ future & $-0.149$ $[-0.246,-0.062]$, $p<.007$ & $\mathbf{-0.129}$ $[-0.227,-0.035]$, $p=.007$ \\
\bottomrule
\end{tabular}}

\begin{tablenotes}[flushleft]
\footnotesize
\item \textit{Note:}  Source paper's coordinate and estimation sample ($8{,}049$ decisions, $72$ subjects); $95\%$ intervals.
\end{tablenotes}
\end{threeparttable}
\end{table}
%  In their own coordinate $\log(e+10)$ on their own estimation sample:

\noindent
The effect AR2019 attribute to projection bias is estimated from a pooled contrast whose identifying variation is in fact driven almost entirely by immediate decisions.
This matters critically for interpretation: projection in the sense of \citet{lor2003} has no scope for immediate choices,\footnote{%
Because the work is performed immediately, the state at the time of choice coincides with the state at execution.} 
and under their model the bias should be flat across time horizons, whereas in the data it falls from $-0.14$ for work chosen minutes ahead to about zero for work chosen days ahead.
\footnote{%
The total raw wedge ($\theta = -0.18$) decomposes as follows: the state's effect on the plan accounts for $-3\%$ $[-34, 18]$ (indistinguishable from zero), its differential effect on the act (the state $\times$ today interaction) accounts for $68\%$ $[26, 115]$, and the remaining $36\%$ $[11, 69]$ reflects PB with the state held fixed. The cost-curvature parameter $b$ is virtually identical across the cold and hot states ($0.657$ vs.\ $0.654$), %, a ratio of $0.997$), 
which ensures that comparisons in effort wedges ($\theta$) directly mirror shifts in time preferences ($\beta$).}


%%% Old backup (something removed)
%Their parameter is identified from a pooled contrast whose entire action sits in the immediate decisions, and the rejection survives the recentering. That matters for interpretation, not only for fit. An immediate decision is the one place where projection in the sense of \citet{lor2003} is not even defined --- the chosen work follows within minutes, so the elevated state \emph{is} the state of the work --- and under their model the contamination should be horizon-flat, whereas it falls from $-0.14$ at a horizon of minutes to zero at a horizon of days. The honest rival to ``$\beff$ fell'' is that ten tasks genuinely raise the marginal cost of the next few; these data cannot separate the two, and the claim is built to survive both: the published $\beta$ is a state-weighted average, and the object estimated is not stable across the state composition of the sample. Because the state is randomized inside the sample, the raw wedge also decomposes: of the whole unmatched wedge ($\theta=-0.184$), the state's movement of the plan is $-3\%$ $[-34,18]$, its movement of the wedge $68\%$ $[26,115]$, and present bias with the state held fixed $36\%$ $[11,69]$ --- the state share is bounded away from zero and the plan share is not, and the split is reported with its interval rather than as a headline. Two more restrictions, tested with the same machinery, \emph{survive}: the cost curvature $b$ is one object ($0.668$ $[0.60,0.73]$ on future-dated choices and $0.633$ $[0.57,0.69]$ on today's when estimated separately, $0.657$ cold and $0.654$ hot; in the joint model every interaction of $\log W$ with the vantage and the state is null, and the two denominators the two-step would use are $0.658$ and $0.656$, a ratio of $0.997$), which is what licenses the comparison in $\theta$; and a randomized payment medium \citep{hoel2022} leaves the intertemporal switching point alone. After a sequence of rejections a reader is entitled to suspect a machine that rejects whatever it is pointed at; the answer is restrictions that pass.


\clearpage
\subsection{Stakes, Proximity, and Other Margins}
\label{sec:margins}

Applying this framework across other state margins and datasets yields the results summarized in the upper section of Figure~\ref{fig:marginsall}. 
For \emph{stakes}, flatness of $\beta$ across wage quintiles is uniformly rejected ($p \le .03$ in every specification and binning), a pattern not driven by any single subject ($72$ leave-one-out refits, all $p < .001$). 
The monotonicity restriction is also rejected at every binning, confirming that $\beta$ does not decline with the stake. 
Evaluating the \emph{temporal proximity} of future choices reveals that allocated effort for a given work date systematically declines as that date approaches ($p = .0004$).
\citet{marchgluth2025} manipulate \textit{hunger} within seventy subjects; in trials featuring a taste--health conflict, the tastier option is chosen over the healthier one in $66.7\%$ of decisions when hungry, compared to $61.4\%$ when sated ($+4.5$\,pp, $p = .035$).
In contrast, \emph{cross-sectional markers} fail to persist under within-person comparison: in \citet{schilbach2019}, intoxicated drivers take up commitment $43\%$ of the time versus $81\%$ when sober ($-0.38$, $t = -8$), 
yet comparing the same driver's choices when sober and when intoxicated at the interview, the effect shrinks to $-0.11$ ($t = -1.0$). 
Likewise, binning workers in \citet{bessone2021} by measured sleep (evaluating \emph{executive load}) reveals a flat trajectory for $\beff$ -- an estimated null that bounds the effect against dual-self predictions. In the same vein, appetite variations in \citet{cheung2022} fail to shift the food-minus-money $\beta$ gap significantly.
Overall, Rung~1 is rejected on margins where the state is exogenously assigned within a person, whereas it is not rejected when state variation is driven by cross-subject characteristics. % or unmanipulated noise.
%Overall, Rung~1 is rejected on every margin where the state is assigned or manipulated within a person, and is not rejected where the state is a person-level marker, a measured rather than assigned state, or a placebo.

% AF: [MAY ADD A CONCLUSION ON (I) placebo check of panel 8 (II) last row represents finding that will be exposed in next sections.]



\iffalse
\begin{figure}[H]\centering
\caption{\footnotesize%
Other state margins}
%\textcolor{green!45!black}{Where each state margin stops on the ladder: the margins assigned or manipulated within a person reject state independence; person-level markers, measured states and the control row do not.}}
\includegraphics[width=.99\linewidth]{fig_margins.pdf}
\label{fig:margins}
\vspace{0pt}
\begin{minipage}{.99\linewidth}
\raggedright
\scriptsize
\textcolor{green!45!black}{\textit{Note:} One panel per row of Table~\ref{tab:margins}, drawn from the estimates. Each tag is the verdict on Rung~1 (state independence): red, rejected; amber, consistent with rejection; gray, a bounded null or a marker that collapses within person; teal, the control row, which survives. Intervals are $95\%$; the stakes, proximity and sleep panels show point estimates by bin. The stakes and proximity panels use the AR2019 subset ($N=72$), on which the ladder is run.}
\end{minipage}
\end{figure}
% fig:margins is generated by writing_sample/TIME_RATIO/code/final/make_figs_v5.py, function fig_margins (python make_figs_v5.py margins); 09-30 attempt, to decide against Table tab:margins
\fi

\begin{figure}[H]\centering
\caption{\footnotesize%
Where $\beta$ moves and where it stops} 
%the state margins of Section~\ref{sec:state}, and the person, the structure and the good of Section~\ref{sec:person}.}
\includegraphics[width=.8\linewidth]{fig_margins_all.pdf}
\label{fig:marginsall}
\vspace{0pt}
\iffalse
\begin{minipage}{.99\linewidth}
\raggedright
\scriptsize
\textcolor{green!45!black}{\textit{Note:} The first two rows are Figure~\ref{fig:margins}. In the last row: the structure compares $\beta$ over effort and over money in the same $101$ people; the person is the dispersion of true $\log\beta$ from the measurement-error model of Section~\ref{sec:memodel}, with profile-likelihood intervals; the good is $\beta$ by good at good-specific elasticities; the snack panel plots, for each of the six snack-choice groups, $\beta$ over the chosen snack against $\beta$ over money for the same people (marker size proportional to the group; dotted line, the diagonal). There, red marks a restriction rejected by a test and amber a pattern reported without a formal test.}
\end{minipage}
\fi
\end{figure}
% fig:marginsall is generated by writing_sample/TIME_RATIO/code/final/make_figs_v5.py, function fig_margins_all (python make_figs_v5.py margins_all); 10-01 extended version of fig:margins, AF to keep one of the two

%"Randomizzando il mezzo di pagamento (contanti vs mobile) su 338 persone, il punto di scelta intertemporale non si sposta (effetto nullo di $+0.52$ righe su 7, $p=.14$, sia nei dati grezzi sia in quelli corretti per il censuramento). Questa restrizione del modello $(\beta, \delta)$ sopravvive, fornendo la riga di controllo che dimostra la validità e specificità dei nostri test."

\iffalse
\begin{table}[H]\centering
\caption{Where each margin stops on the ladder. Rung~2 is tested by the distance-to-the-cone
statistic, whose null is that the profile \emph{is} monotone.}
\label{tab:margins}
{\small
\begin{tabular}{>{\raggedright\arraybackslash}p{3.9cm}>{\raggedright\arraybackslash}p{5.8cm}>{\raggedright\arraybackslash}p{4.4cm}}
\toprule
Margin & Rung 1 (state-independence) & Rung 2 (monotone) \\
\midrule
Visceral, \textbf{assigned} (just-worked) & \textbf{rejected}: $\beta$ $.88\to.74$; RI $p=.09$ / $.02$ recentered & n/a (binary) \\
Stakes (wage quintile) & \textbf{rejected}, $\chi^2_4=27.8$, $p<10^{-4}$; every binning, every subject & \emph{decreasing} rejected; \emph{increasing} survives at coarse binnings \\
Proximity, within work date & \textbf{rejected}, $p=.0004$ & not rejected \\
Hunger, \textbf{induced} (attribute analogue) & consistent with rejection, $p=.035$ & n/a \\
Intoxication, person marker & collapses within person & n/a \\
Executive load / sleep & bounded null & n/a \\
Appetite $\times$ good, measured & bounded null: $\le0.042$ per s.d. & n/a \\
Payment medium (control row) & survives & n/a \\
\bottomrule
\end{tabular}}
\end{table}
\fi


%The pattern is a rule, not a coincidence: every margin that rejects Rung~1 is one where the state is assigned or manipulated within a person; every margin that stops there is a person-level marker, an untested dose, or a control row. It is the empirical form of the requirement that Section~\ref{sec:balance} derived. A very large natural shift fits it too: across three waves straddling the onset of COVID-19 --- about a thousand responses, $126$ people in all three --- structural time preferences did not move \citep{drichoutis2022}; a common shock that nobody is assigned changes the date, not the state at the moment of choice.

% =====================================================================================
%\section{Evidence II: the person, the structure, and the good}
\section{Present Bias Across Individuals, Goods, and Structures}
\label{sec:person}


%\subsection{One $\beta$ per person: a model, not a correlation}
%\subsection{Rejecting the stable-trait hypothesis}
\subsection{Failure of the Stable-Trait Hypothesis}

Comparing time preferences across domains, \citet{ans2015} elicit real-effort and monetary discounting within the same $101$ subjects, finding substantial PB over effort but little over money.
Re-estimating the aggregate data recovers the same disparity ($\beta_{\text{effort}} = 0.88$ versus $\beta_{\text{money}} = 0.99$). The gap is fully robust to accounting for monetary corner choices via a %two-limit 
Tobit model and to expressing the gap in $\log \beta$, where non-linear transformation adjustments shift the estimate by only $1\%$.
To evaluate whether $\beta$ is stable across domains at the person level, I apply the measurement-error framework in Section~\ref{sec:memodel}, using domain-specific elasticities (Online Appendix~\ref{app:me}), and find:
\iffalse
\begin{center}
\emph{effort:} $\mathrm{SD}(\text{true }\log\beta)=0.24$ $[0.16,0.32]$ \qquad\qquad
\emph{money:} $\mathrm{SD}(\text{true }\log\beta)=0.00$ $[0.00,0.04]$.
\end{center}
\fi
\[
\begin{aligned}
\text{Effort:}\quad & \mathrm{SD}(\text{true }\log\beta) = 0.24 \quad [0.16,\,0.32],\\
\text{Money:}\quad  & \mathrm{SD}(\text{true }\log\beta) = 0.00 \quad [0.00,\,0.04].
\end{aligned}
\]
For effort, the plan-versus-act wedge exhibits genuine population heterogeneity beyond sampling noise,\footnote{%
The hypothesis of zero dispersion is rejected ($p < .00001$), %; and a leave-one-out analysis keeps the estimated SD bounded between $0.18$ and $0.24$.
even under leave-one-out robustness tests.} 
whereas PB over money displays no such variation.%within the same subjects.
\footnote{%
Among the half of subjects with the most precise monetary estimates, the observed dispersion is $0.059$ against a median sampling standard error of $0.077$, indicating that the entire spread reflects sampling noise.} 
Because population PB over money is effectively constant while that over effort varies across individuals, population correlation is undefined; empirically, the sample correlation between individual point estimates is just $+0.03$ $[-0.10, +0.15]$. These results are inconsistent with PB as a stable person-level trait across domains.


Applying the measurement-error model within a given sign--timing structure using data on monetary, healthy, and unhealthy food rewards from \citet{cheung2022} reveals a different pattern. In that dataset, both goods in each pairwise comparison exhibit genuine population dispersion ($\text{true SD}\ \sigma \in [0.20, 0.23]$), and high cross-reward correlations.\footnote{%
Specifically: money--healthy $+0.92$, %$[0.81, 1.00]$, 
money--unhealthy $+0.89$, and healthy--unhealthy $+0.71$ %$[0.55, 0.87]$ 
(with the caveat that all three are elicited in the same sittings, preventing the removal of a common-occasion component).}
Taken together, these results suggest that $\beta$ behaves as a stable trait within a given sign--timing structure, whereas it does not constitute a single invariant object across structures. This distinction is consistent with findings in the literature, 
including the temporal stability of monetary $\beta$ \citep{meiersprenger2015} and the weaker stability of experimental measures in field samples \citep{chuang2015}.\footnote{%
%This conclusion is not invalidated by \citet{cerrone2026}, who report a correlation of $+0.78$ between their effort and money parameters. In their setting, both parameters are estimated jointly in a single regression of effort choices on work and payment timing, yielding a median sampling correlation of $+0.83$ between the two coefficients. Consequently, that correlation reflects how collinear regressors share explanatory power within a single choice, rather than preference stability across distinct elicitations.}
This conclusion is consistent with \citet{cerrone2026}, who report a rank correlation of $0.77$ between their individual effort and money parameters. Since both parameters are estimated from the same effort choices, the correlation is measured within one trade-off structure (effort exchanged for pay).} 
%PB over effort from the response of effort to the timing of the work, and PB over money from its response to the timing of the payment. 
%-- effort exchanged for pay -- rather than between an effort allocation and a money allocation. 
%It is not a by-product of joint estimation: the two estimates share a sampling correlation of $+0.83$, but uncorrelated true parameters would produce an observed correlation of only about $+0.18$, and correcting for sampling noise raises it to $+0.89$.}




%\subsection{What the effort--money gap is made of}
\subsection{Commodity Categories and Person--Reward Pairings}
% it's about person-good; also the thing of drinkers/smokers...
\label{sec:goods}

%%%%% cheung2022 %%%%%
% : 697 studenti di scuole superiori cinesi a basso/medio reddito
% 2 sessioni di scelta within subject:
%% Settimana 1 (Sessione Advance): I soggetti compiono scelte su allocazioni di ricompense da ricevere entrambe nel futuro (Settimana 2 vs. Settimana 5)
%% Settimana 2 (Sessione On-the-day): Gli stessi soggetti rifanno le stesse scelte sugli stessi due orizzonti (Settimana 2 vs. Settimana 5), ma la prima data di pagamento coincide con il giorno stesso della scelta (oggi)
%% Settimana 5: Solo erogazione finale delle ricompense ritardate
% (Money): 1 unità = 0.5 RMB; (Healthy Food): scelta individuale dello snack preferito tra noci, mandorle e uvetta; (Unhealthy Food): scelta individuale tra Skittles, M&M’s e patatine Lays
% Convex Time Budget (CTB). Per ciascun bene, i soggetti affrontano 14 decisioni di budget per sessione (variando i tassi di interesse e includendo test per violazioni dell'assioma GARP)
% Misure Descrittive: Impazienza: quota allocata alla data più vicina nella Sessione; .Present Bias descrittivo: incremento percentuale della quota allocata alla data più vicina quando la scelta avviene "oggi" (Sessione 2) rispetto a quando era "futura" (Sessione 1); Smoothness: misura dell'inclinazione a bilanciare i pagamenti tra le due date (proxy descrittiva della curvatura dell'utilità)
% Stima Strutturale ($\beta-\delta$):Assumono una funzione di utilità separabile nel tempo con curvatura CRRA $u(c) = \frac{c^{1-\alpha}}{1-\alpha}$ e sconto quasi-iperbolico $(\beta, \delta)$
%% Results
%%% Sorprendentemente, il present bias è più forte per il denaro che per il cibo ($\beta_{\text{money}} = 0.65$ contro $\beta_{\text{healthy}} = 0.69$ e $\beta_{\text{unhealthy}} = 0.71$ nello MNL aggregato) [insignif diff tra 2 cibi]
%%% L'utilità per il denaro è quasi lineare, mentre per il cibo si registra una marcata concavità (forte preferenza a spalmare il consumo di cibo nel tempo)
%%% Correlazione tra Beni: Le misure individuali di impazienza e present bias mostrano correlazioni moderate tra i vari beni ($r \approx 0.60$),

%%%%% COSA AGGIUNGO IO %%%%%
% Rigetto contemporaneo sia della lettura "Commodity" che "Temptation"Nel loro paper: Cheung et al. prendevano atto dell'ordinamento $(\beta_{\text{money}} < \beta_{\text{food}})$ spiegandolo con il contesto specifico dei giovani cinesi low-income
%  Rilassando il vincolo e permettendo un'elasticità specifica per ciascun bene ($b_{\text{good}}$), il gap si ridimensiona ma conferma in modo pulito e robusto che il denaro rimane il bene più present-biased ($\beta_{\text{money}} = 0.82$, $\beta_{\text{healthy}} = 0.85$, $\beta_{\text{unhealthy}} = 0.86$, $p < .0001$) [MA SE NE ERANO GIA RESI CONTO LORO...]
%  Il Present Bias non è una proprietà della "categoria di bene", ma della coppia (Persona, Snack)Scoperta: Scomponendo le scelte dentro le categorie di cibo in base allo specifico snack scelto da ciascuno studente (3 opzioni sane e 3 non sane), scopri che il $\beta$ varia tra $0.82$ e $0.89$ solo all'interno della voce "cibo" — un intervallo più ampio dell'intero scostamento denaro-cibo
% trovo correlazioni a livello individuale across goods anche con mio ME model... (DENTRO la stessa struttura di scelta {QUI}, ma fallisce di essere un unico tratto TRA strutture diverse (ANS2015))
% Batteria di Controlli Metodologici e Bounded Nulls (Audit di Robustezza): vedi LM.
% Dove vive davvero il gap (Drinkers & Smokers): Vai a scavare nei dati individuali e scopri per chi il denaro è più present-biased:Per la maggioranza degli studenti ($585$ su $697$), il gap è modesto ($-0.023$, $p = 0.10$).Il gap è quasi interamente concentrato nei ragazzi che fumano o bevono ($95$ e $50$ soggetti), per i quali $\beta_{\text{money}}$ crolla a $0.74$, perché per loro il denaro contante ha un uso immediato per tentazioni viscerali non osservate.



In the literature, cross-domain disparities in PB are typically interpreted as intrinsic properties of the underlying commodities, a view that potentially conflates distinct axes of variation (Figure~\ref{fig:axes_edited}). This interpretation can be tested using the data from CTW2022, which hold fixed several confounding factors present in conventional comparisons: their design uses a single protocol across $697$ subjects and two sessions, with three rewards evaluated on the same convex-time-budget grid over the same dates, delivered in identical bags and presented with matching screen representations.%
\footnote{% 
Although the design does not randomize immediacy, all three rewards are elicited from every subject in both sessions, so common session-level shocks cancel out in the cross-reward contrast. The identified object is the differential today-effect across goods, controlling for subject $\times$ reward fixed effects.} 
Under the \textit{commodity} hypothesis, money should exhibit the least PB; under the \textit{temptation} interpretation, unhealthy food should exhibit the most PB. Applying the scale-free two-step model with good-specific elasticities yields the estimates in Table~\ref{tab:goods}, consistent with the ordering reported by CTW2022.

\begin{table}[H]\centering
\caption{Present bias varies across goods and is stronger for monetary rewards.}
\label{tab:goods}
\begin{threeparttable}
{\footnotesize
\renewcommand{\arraystretch}{1.12}%altezza colonne
\setlength{\tabcolsep}{27pt}%larghezza colonne
\begin{tabular}{lccc}
\toprule
Good & elasticity $b$ & shift $\theta$ & $\beta$ \\
\midrule
\textbf{money} & $1.713$ & $-0.341$ & $\mathbf{0.820}$ $[0.784,0.855]$ \\
healthy food & $1.398$ & $-0.224$ & $0.852$ $[0.808,0.893]$ \\
unhealthy food & $1.425$ & $-0.213$ & $0.861$ $[0.819,0.900]$ \\
\midrule
\multicolumn{2}{l}{money $-$ food average} & \multicolumn{2}{l}{$\mathbf{-0.037}$, $p<.004$} \\%$[-0.064,-0.013]$
%\multicolumn{2}{l}{money $-$ healthy; money $-$ unhealthy} & \multicolumn{2}{l}{$-0.032$, $p=.052$; \ $-0.041$, $p=.004$} \\
\multicolumn{2}{l}{money $-$ healthy} & \multicolumn{2}{l}{$-0.032$, $p=.052$} \\
\multicolumn{2}{l}{money $-$ unhealthy} & \multicolumn{2}{l}{$-0.041$, $p=.004$} \\
\multicolumn{2}{l}{unhealthy $-$ healthy} & \multicolumn{2}{l}{$+0.009$, $p=.73$} \\
\bottomrule
\end{tabular}}
\begin{tablenotes}[flushleft]
\footnotesize
\item \textit{Note:} 
Three goods, one protocol, within subject and within good; good-specific price
elasticities; subject bootstrap, $B=499$. With a common elasticity imposed the money-minus-food
gap is $-0.068$; the elasticities differ by $+0.32$ and $+0.29$ ($p<.004$), and freeing
them absorbs $45\%$ of it.
\end{tablenotes}
\end{threeparttable}
\end{table}

\noindent
These findings reject both hypotheses, as PB is strongest for money and indistinguishable between the two food types.%
\footnote{This ordering is robust across five estimators -- including an interval regression that accounts exactly for the design's six-bundle discreteness -- and across both arms of the randomized presentation order. 
%Online 
Appendix~\ref{app:goods} shows that it is not driven by arbitrage, trust in the experimenters, stated valuations of the rewards, or the scale of delivery. 
%The design cannot test the position of money, which always appears second; the two foods, which swap between the first and third positions, differ by only $-0.013$, and money lies $0.036$ below the line joining them.
}
%%%
While CTW2022 suggest that monetary PB may reflect features of their non-WEIRD, low-income adolescent sample, they do not explain why it exceeds PB over food.
By disaggregating choices within each food class, I find that $\beta$ varies across specific snacks by more than the overall money--food gap: it ranges from $0.82$ to $0.89$ for healthy options and from $0.82$ to $0.88$ for unhealthy ones.%
\footnote{Specifically, PB is strongest for the snacks preferred by individuals with high sweet-craving scores.} 
This spread can be decomposed because every subject also made choices over money. In fact, grouping subjects by their chosen snack, the cross-group standard deviation of $\beta$ is $0.033$ over food and $0.034$ over money for the same individuals, with the two group means correlated at $+0.89$. This tight co-movement indicates that snack-level variation is predominantly a \emph{person} effect, %(selection),
accompanied by a \emph{good} effect of $+0.036$ on average that raises $\beta$ for food relative to money across all six groups.
Consequently, a standard commodity taxonomy fails to organize these data (Figure~\ref{fig:goods}): the broad label ``food'' spans $\beta \in [0.82, 0.89]$, while money's own $\beta$ ranges from $0.84$ among non-substance users to $0.74$ and $0.76$ among drinkers and smokers.





%%% AF backup
%Money is the most present-biased good and the two foods are indistinguishable, so both readings fail in one regression; the ordering holds under five estimators, including an interval regression that handles the design's six-bundle discreteness exactly, and in both arms of the randomized presentation order. What it is not, each with a number in the Online Appendix: not arbitrage (a background stock, simulated on these budgets, collapses the curvature estimate and biases $\beta$ \emph{toward one}, so it works against the finding); not trust in the experimenters, not stated valuation of the rewards, not the scale of the delivery (an income experiment the design contains: at a price of one the same subject faces budgets of $70$ and $56$ units, and the immediacy premium does not move). What the design cannot rule out is position: money is presented second to every subject, so its position has $\rho=0$ and no test exists inside the design --- stated as the limit it is. What bounds it is that the two foods do swap between the first and third blocks: their position contrast is $-0.013$, a third of what would be needed, and money sits $0.036$ \emph{below} the line joining the two food positions, so position could account for money only if its effect were sharply non-monotone --- flat between the first and third blocks and dipping in the second --- which no account of attention or fatigue predicts.
%\citeauthor{cheung2022} explain why money shows present bias at all --- a non-WEIRD, low-income adolescent sample --- and not why it exceeds food; the sharpest evidence on that is one the published analysis does not use. Within each food class subjects chose one of three snacks, and $\beta$ differs across those choices by more than the money--food gap: $0.82$, $0.89$, $0.87$ for the healthy options and $0.82$, $0.87$, $0.88$ for the unhealthy ones, the sugary end the present-biased end. That spread separates cleanly, because every one of these subjects also chose over money: across the six snack-choice groups the standard deviation of $\beta$ is $0.033$ for the chosen snack and $0.034$ for the same people's money, correlated at $+0.89$ --- most of it a \emph{person} effect --- with a \emph{good} effect of $+0.036$ on top that appears in all six groups and never changes sign. A commodity taxonomy does not organize these data (Figure~\ref{fig:goods}): the label ``food'' spans $0.82$--$0.89$, and money's own $\beta$ runs from $0.84$ among abstainers to $0.74$ among drinkers.

\begin{figure}[H]\centering
\caption{\footnotesize%
Money is the most present-biased good, and within food $\beta$ follows the person who chooses the snack.}
\includegraphics[width=.99\linewidth]{fig_goods.pdf}
\label{fig:goods}
% fig:goods is generated by writing_sample/TIME_RATIO/code/final/make_figs_v5.py, function fig_goods (python make_figs_v5.py goods), from output/route3_v1_summary.csv; 10-01 redesign (old: figs/diagrams/_backups/fig_goods_before_1001.pdf)
\vspace{5pt}
\begin{minipage}{.99\linewidth}
\raggedright
\scriptsize
\emph{(A):} $\beta$ by good, $95\%$ CI. \emph{(B):} the 6 snack-choice groups, $\beta$ over the chosen snack against $\beta$ over money for the same people (marker size proportional to the group): a person effect along the diagonal, and a good effect %of $+0.036$ 
displacing every group above it.
\end{minipage}
\end{figure}
\clearpage


\noindent
This contrast sheds light on a potential driver of the monetary PB disparity. 
A plausible mechanism is that, for low-income adolescents, cash provides a general claim on discretionary immediate consumption, whereas a single food item offers only a non-fungible reward tied to a specific consumption good.
In particular, subjects reporting alcohol ($N=95$) or tobacco use ($N=50$) may face immediate cash needs that no single snack can satisfy. Crucially, their estimated $\beta_{\text{money}}$ falls to $0.74$ and $0.76$ (compared to $0.84$ among the $585$ non-users), whereas their $\beta_{\text{food}}$ values change much less ($0.82/0.87$ and $0.89/0.85$, against $0.86/0.86$). %
If these individuals were simply more impatient in a domain-general sense, their discount factors would shift uniformly across all rewards. 
%The questionnaire's direct patience measure illustrates such a parallel shift: a one-standard-deviation increase is associated with a change of $+0.043$ in $\beta_{\text{money}}$ and $+0.034$ in $\beta_{\text{healthy}}$.
In fact, the questionnaire's direct patience measure moves money and healthy food together: a one-standard-deviation increase raises $\beta_{\text{money}}$ by $0.043$ and $\beta_{\text{healthy}}$ by $0.034$.
%In contrast, substance use exhibits a domain-specific tilt: it depresses $\beta_{\text{money}}$ by $-0.034$ while leaving food discount factors essentially unaffected ($-0.012$ for healthy food and $+0.003$ for unhealthy food; (Figure~\ref{fig:tilt}). %
In contrast, alcohol and tobacco use are associated with a domain-specific tilt: a lower $\beta_{\text{money}}$ and essentially unchanged food factors.
Three additional checks support this interpretation. First, the money-specific gap persists after controlling for subject $\times$ reward fixed effects -- which absorb baseline impatience and utility levels -- as well as for trait self-reports and individual covariates (wealth, BMI, age, gender, trust, and choice consistency). Second, in a reverse regression that includes PB measures from all three domains jointly, only monetary PB significantly predicts drinking ($+0.41$) and smoking ($+0.51$). Third, drinkers do not report higher subjective valuations of cash, providing evidence against differential marginal utility of money as the driver of these results.

%%% AF backup
%The honest limit is multiplicity: drinking was not chosen in advance, and under a Romano--Wolf stepdown over the twelve-moderator battery nothing survives (adjusted $p=.23$); the case rests on the stability of the pattern --- money-only in every specification, four of five snack cells, a leave-one-out range with no sign flip --- and on a randomized design, stated in Section~\ref{sec:limits}, that would turn it into evidence.


\begin{figure}[H]\centering
\caption{\footnotesize%
A trait shifts the profile; drinking and smoking tilt it toward money.}
\includegraphics[width=.75\linewidth]{d07_tilt_vs_shift.pdf}
\label{fig:tilt}
\vspace{-15pt}
\end{figure}
% fig:tilt is drawn in figs/diagrams/d07_tilt_vs_shift.tex (TikZ; pdflatex d07_tilt_vs_shift); AF's 10-01 version: figs/diagrams/_backups/d07_tilt_vs_shift_AF_1001.tex

\noindent
The findings in this section reject both the view of PB as a stable individual trait and its interpretation as an intrinsic property of the commodity class. Rather, PB attaches to the specific (person, reward) pair within a given trade-off structure and state.





% =====================================================================================
%\section{Magnitudes, and welfare}
\section{Magnitude, Welfare, and Policy Implications}
\label{sec:magn}


%\subsection{How large, and how many}\label{subsec:magnitude}
\subsection{Bounds on $\beta$ and Behavioral Type Shares}

Rejecting the stable-trait specification does not make the magnitude of PB uninformative. 
Within a given setting, $\beta$ still summarizes the weight placed on the immediate component for the states, goods, structures and people that the dataset covers. 
Two questions remain well defined: which values of $\beta$ can be empirically ruled out, and what proportion of individuals exhibits PB. 
I address the former through three analytical frameworks of decreasing structural stringency -- ranging from parametric estimation to RP sets with no functional form assumptions on utility (Online Appendix~\ref{app:magn}).%
\footnote{The \emph{parametric} estimate is the two-step estimator of Section~\ref{sec:coord}. %, which assumes a constant wage elasticity. 
The \emph{measurement-error-robust} set relies only on the wage-matching identity, a monotone supply curve for each subject, and centered measurement error; 
%it retains all subjects by extending the 
it extends the framework of \citet{aguiar2021} % -- originally formulated for exponential discounting -- 
to the quasi-hyperbolic model. 
The \emph{nonparametric} bracket assumes monotone supply alone and bounds $\beta$ between wage grid points; %, as in Section~\ref{sec:matched}; 
by construction, it cannot exclude $\beta=1$.}

\noindent
In the effort domain, PB is systematically detected: in \citet{augenblick2019}, the hypothesis that no wage quintile displays a positive plan-versus-act wedge is rejected ($p=.004$) using the joint moment-inequality test of \citet{coxshi2023}. In the same panel, the parametric estimate is $0.80$ (the state-weighted average of Section~\ref{sec:matched}), the measurement-error-robust (MER) set is $[0.75, 0.95]$, and the nonparametric bracket is $[0.72, 1.00]$. In the independent real-effort study of \citet{fedyk2024}, the corresponding figures are $0.74$ $[0.68, 0.80]$, $[0.80, 0.95]$, and $[0.67, 1.00]$. %
%\footnote{Fedyk's MER set is computed on interior allocations; including corners makes it empty as an artifact of censoring.} 
%Exponential discounting is rejected at the MER level across both experimental settings, with the identified intervals reflecting true behavioral dispersion rather than a lack of estimation precision.
Exponential discounting is rejected at the MER level in both experiments; in AR2019 the identified intervals reflect genuine behavioral dispersion rather than estimation imprecision alone. 

%%%%%%%%%%%%%%%%%%%%%%%%%%%%%%%%%%

The same empirical question applies to household consumption in real-world settings. 
The Spanish household panel (ECPF) is a standard testing ground for RP over time \citep[e.g.,][]{adams2014}. On these data, the deterministic test of \citet{browning1989} rejects exponential discounting (ED) for $78\%$ of single-person households and $82\%$ of couples, and the quasi-hyperbolic discounting (QHD) conditions of \citet{bbc2021} 
%under which a household re-plans every quarter, 
rationalize nearly half of those rejected under ED. \citet{aguiar2021} show that ED rejections disappear once each quarter's expenditure is allowed to carry centered measurement error. 
%Under that error model, however, goods can be reallocated freely within the quarter, and every household in the panel is then rationalized with exactly its observed quarterly spending, so the test cannot distinguish one discounting model from another.
Under that error model, however, goods can be freely reallocated within each quarter while preserving observed total spending, rationalizing every household and leaving the test unable to distinguish among discounting models.
Instead, centering the error within eight broad groups of goods each quarter \citep[the aggregates of][]{adams2014} preserves robustness to measurement error while restoring power.
I therefore extend the MER test for ED developed by \citet{aguiar2021} to the QHD conditions of \citet{bbc2021} (Online Appendix~\ref{app:ecpf}).
Figure~\ref{fig:fieldgrid} shows that ED is rejected at every value of $\delta$, whereas QHD with $\beta\le0.95$ and $\delta$ close to one is not.
%Specifically, rationalizing every household requires errors of at least $6\%$ of spending in each group of goods under ED, against about $1\%$ under QHD.
Specifically, the measurement error needed to rationalize every household is at least $6\%$ of spending in each goods group under ED, compared with about $1\%$ under QHD.%
\footnote{This is not direct evidence of PB; unlike experiments, observational consumption paths bound $\beta$ only from above.}
%\footnote{The test of \citet{aguiar2021}, with latent consumption satisfying the Afriat inequalities, centered measurement-error moments and the entropic statistic of \citet{schennach2014}, but with the marginal utility of wealth allowed to grow at a rate in $[1/\delta,1/(\delta\beta)]$ \citep{bbc2021} in place of $d^{-t}$. Three implementation choices depart from their code: each household's latent path starts at its best-fitting path rather than at $d=1$, the sampling density is scaled to each household's spending (an error scale of $2\%$), and the sampler can leave goods with zero consumption. Under their own error model, these choices also overturn their rejection for couples, since every couple admits a latent path with exactly its observed quarterly spending at some $d<1$. Seasonal tastes do not account for the gap: exponential discounting with three common seasonal factors still requires errors of $5$--$8\%$. All verdicts in the figure come from long chains, which give the same verdicts across seeds; exponential discounting is rejected, and $\beta\le0.95$ with $\delta\ge0.99$ is not, in every run. In simulations with the same prices, the test rejects in none of $30$ samples drawn from exponential discounters, and in none drawn from households with $\beta=0.9$. On the coarser discount grid of \citet{aguiar2021}, the deterministic rejection rates are $84\%$ and $90\%$.}
%This contrast, however, is not evidence of present bias. In simulated panels with the same prices, households with $\beta=0.9$ would make exponential discounting require errors of only about $0.5\%$: the observed gap reflects the volatility of quarterly spending, which the flexible path that quasi-hyperbolic discounting allows for the marginal utility of wealth absorbs, and the test has no power against present bias of the size measured in the experiments above. 
%More fundamentally, consumption paths bound $\beta$ only from above: a household that re-plans every quarter with a lower $\beta$ and a lower propensity to spend out of wealth behaves like one with a higher $\beta$ \citep[Corollary~3]{bbc2021}. 
%Household consumption data therefore leave the present-bias factor unidentified even when measurement error is modeled explicitly; the plan-versus-act contrast of the experiments above is what identifies it.


\iffalse % OLD WRONG
The same empirical question applies to household consumption in real-world settings. 
%On the Spanish household panel (ECPF), widely examined in the revealed-preference literature \citep[e.g.,][]{adams2014, bbc2021}, a deterministic RP test rejects exponential discounting for $82\%$ of single-person households. However, \citet{aguiar2021} demonstrate that this rejection vanishes once the model allows for centered measurement error in expenditure. 
The Spanish household panel (ECPF) is a standard testing ground for RP over time \citep[e.g.,][]{adams2014}. On this panel, \citet{bbc2021} characterize quasi-hyperbolic discounting and find that it fits the consumption of couples better than exponential discounting, with a mean $\beta$ of about $0.84$. The deterministic test of \citet{browning1989} rejects exponential discounting for $82\%$ of single-person households, but \citet{aguiar2021} show that this rejection disappears once expenditure is allowed to carry centered measurement error.
Extending their MER test from exponential to quasi-hyperbolic discounting, 
I find that the admissible set is $\beta\in[0.80,1.00]$ (Figure~\ref{fig:fieldgrid}): values below $0.80$ are rejected, time consistency fits the data, and $\beta=0.80$--$0.85$ survives only with $\delta=1$.
\fi



%AF, quindi:
%%% sezione 2 mostra come interagiscono le FOC conditions
%%% in sezione 3 e parte su effort, le FOC conditions che si combinano sono quella del piano e quella dell'atto.
%%% in partesu ECPF, le FOC che si combinano/confrontano invece sono quelle dell'atto di un quarter e quella dell'atto di un altro quarter.



%%%%%%%%%%%%%%%%%%%%%%%%%%%%%%%%%%
The second question concerns the population prevalence of PB. Using revealed and predicted PB over effort and money from \citet{cerrone2026}, I bound the population share of each behavioral type while accounting for the sampling error of each measure.%
\footnote{Each participant is placed in one of eight cells by three binary signals: PB over effort, PB over money, and whether the prediction anticipates the sign of the effort wedge. A mixing matrix gives the probability that each of five types %(time-consistent, na\"ive, sophisticated, future-biased, money-driven) 
lands in each cell, with noise calibrated to each measure's within-subject sampling error, and a linear program bounds the type shares over a simultaneous $95\%$ confidence box around the observed cell shares (Online Appendix~\ref{app:types}).}
Between $67\%$ and $79\%$ of subjects are time-consistent, between $9\%$ and $32\%$ are naïve present-biased, and at most $11\%$ are sophisticated (with residual shares bounded at $1.5\%$ for future-biased and $12\%$ for money-driven types). 
PB is thus empirically present but confined to %concentrated in
a minority of subjects, many of whom fail to anticipate it.



% previous version of this subsection (before 10-01), kept but not printed:
\iffalse
With the higher rungs rejected, the two lowest remain: does present bias exist (Rung~4),
and what does the trait model look like when imposed (Rung~3). On the effort panel the
plan-versus-today wedge is positive in every wage quintile jointly (Cox--Shi $p=.004$),
the one place where the level statistic is the right one. The magnitude is reported at
three rungs of decreasing assumptions: the two-step point ($0.66$ $[0.56,0.76]$ on
\citeauthor{augenblick2019}, $0.74$ $[0.68,0.80]$ on \citeauthor{fedyk2024}'s independent
effort design); a measurement-error-robust confidence set that imposes only the
wage-matching identity, a per-subject monotone supply curve and centered measurement error,
and drops nobody ($[0.75,0.95]$ and $[0.80,0.95]$, exponential discounting rejected on
both); and a nonparametric bracket from monotonicity alone ($[0.72,1.00]$, $[0.67,1.00]$),
which cannot exclude one by construction. Over effort $\beta$ lives in the high $0.7$s to
low $0.9$s; on the Spanish consumption panel, where the deterministic test of
\citet{aguiar2021} rejects exponential discounting for $82\%$ of households, the
noise-robust quasi-hyperbolic set is $[0.80,1.00]$ --- deep present bias excluded, time
consistency fitting exactly --- and a $\beta\le0.75$ is rejected outright. That the
effort set spans two tenths is not slack: a set for a parameter that the same data put at
$0.88$ in one state and $0.74$ in the other could not be much narrower. The population
answer is a set of type shares from \citet{cerrone2026}'s two binary margins (revealed and
predicted present bias), identified by a linear program over a mixing matrix calibrated
to each margin's own sampling error: time-consistent $[0.67,0.79]$, na\"ive present-biased
$[0.09,0.32]$, sophisticated $[0.00,0.11]$. The wedge is real and carried by a minority of
one in ten to one in three, mostly na\"ive about it --- in this snapshot, one state, one
structure.
\fi


\begin{figure}[H]\centering
\caption{\footnotesize%
The measurement error that $(\beta,\delta)$ pairs require to rationalize the ECPF data.}
\includegraphics[width=.9\linewidth]{fig_field_grid.pdf}
\label{fig:fieldgrid}
\vspace{3pt}
\begin{minipage}{.99\linewidth}
\raggedright
\scriptsize
\textit{Note:} Each cell reports the average error, in percent of spending in each group of goods and quarter, needed for every household to satisfy the QHD conditions. 
%of \citet{bbc2021} with the marginal utility of wealth growing at a rate between $1/\delta$ and $1/(\delta\beta)$; $\beta=1$ is exponential discounting, and a lower $\beta$ nests every higher one. 
The blue line encloses the pairs not rejected at the $5\%$ level by the MER test. %of \citet{aguiar2021}, with errors centered within eight groups of goods in each quarter and an error scale of $2\%$. In simulated panels with the same prices, households with $\beta=0.9$ would make exponential discounting require an error of only about $0.5\%$: the larger errors observed reflect the volatility of quarterly spending, which a flexible path for the marginal utility of wealth absorbs. 
ED needs errors of at least $6\%$ of spending, QHD with $\beta\le0.95$ and $\delta\ge0.99$ about $1\%$. %; present bias of the size measured in the experiments would not create such a gap
Without measurement error, every pair is rejected.
\end{minipage}
\end{figure}
% fig:fieldgrid (10-08) is generated by writing_sample/TIME_RATIO/code/final/ecpf/fig_field_grid.py from output/ecpf/reqerr_<kind>_adams8.npz (required_error.py) and cells_<kind>_adams8_k0.02.csv (run_cells.py); explained in FINAL_WRITING/old/ECPF_explained_1008.md. Previous figure: figs/diagrams/_backups/fig_field_grid_before_1008.pdf; previous tex: figs/diagrams/_backups/main_short_AF_1008.tex






%\subsection{Welfare: plans, outcomes, and commitment}
%\subsection{Welfare: plans, outcomes, commitment, and policy}
\subsection{Welfare Losses, Commitment, and Policy Design}
\label{sec:welfare}



Construing a plan-versus-act wedge as a self-control \emph{problem} inherently requires deciding whose preferences count. The \textit{planning} self and the \textit{acting} self rank the same options differently, and choices alone cannot determine which evaluation should prevail \citep{bernheimrangel2009}. %without arbitrarily privileging one of them
To limit arbitrary reliance on any single normative standpoint, this section adopts a deliberately conservative criterion. 
The objective is to establish a welfare benchmark that can secure broad agreement across normative perspectives -- accepting a higher detection threshold with the aim of ensuring that the resulting welfare claims remain defensible under different evaluative frameworks. %and does not rest solely on an ``ascetic paternalism''.

Specifically, a deviation that an individual did not anticipate is an error in forecasting their own behavior. Choices that rest on this forecast -- such as whether to demand commitment or which contract to accept -- are made under a misperception of their consequences, and are therefore revealed mistakes rather than evidence of preferences \citep{bernheimrangel2009}. 
I therefore define the \emph{self-control gap} $\Gamma$ as the unanticipated component of the wedge (the shortfall of the realized choice relative to the person's own forecast of it) and treat only $\Gamma$ as welfare-relevant. 

While an anticipated deviation ($e_{\rm plan}>e_{\rm act}=e_{\rm pred}$) reflects PB, it constitutes a conflict between temporal selves on which no welfare claim is asserted. 
Although this conservative filter may leave some genuine welfare losses undetected, its value lies in establishing a defensible baseline for policy action -- one whose mandate for intervention does not depend on the normative position adopted. 
Crucially, this benchmark is not empty: as shown below, most of the observed wedge is unanticipated.


\begin{figure}[H]\centering
\caption{\footnotesize%
Most of the wedge is unforeseen, and commitment demand does not track the forecast error.}
\includegraphics[width=.99\linewidth]{fig_welfare.pdf}
\label{fig:welfare}
\iffalse
\vspace{5pt}
\begin{minipage}{.99\linewidth}
\raggedright
\scriptsize
\textit{Note:} \emph{(A):} AR2019, all $100$ subjects; $\beta$ implied by the subjects' own incentivized forecasts of their choices on the day (open markers) and by those choices (filled), with the state held cold and hot; $\Gamma$ is the gap between the two. \emph{(B):} forecast minus outcome; for resolutions, belief in success minus self-rated success, both on $0$--$100$ scales; elsewhere, percentage points. \emph{(C):} correlation of the choice of sobriety incentives over cash with the stated belief about sobriety, realized sobriety and their difference ($63$ drivers in \citealt{schilbach2019}), and, in gray, of commitment demand with the taker's own measured wedge in four further datasets. All intervals are $95\%$ bootstrap intervals.
\end{minipage}
\fi
\end{figure}
% fig:welfare is generated by writing_sample/TIME_RATIO/code/final/make_figs_v5.py, function fig_welfare (python make_figs_v5.py welfare), from output/welfare_v5_summary.csv (welfare_v5.py: blocks WFORE, WSCH), exp_v3, field_v3 and alcohol summaries; 10-03 version, replacing the TikZ diagram figs/diagrams/d10_welfare.tex






\noindent
In AR2019, subjects also make incentivized predictions of the choices they will make on future work days, at five wages drawn at random.
Re-estimating them in the two-step coordinate of Section~\ref{sec:coord}, separately by assigned state, yields $\beta_{\text{cold}}^{\text{pred}} = 1.01$ and $\beta_{\text{hot}}^{\text{pred}} = 0.99$, against realized values of $\beta_{\text{cold}}^{\text{real}} = 0.88$ and $\beta_{\text{hot}}^{\text{real}} = 0.74$ (Figure~\ref{fig:welfare}A).%
\footnote{AR2019 report the same near-complete na\"ivet\'e %in their own specification 
(a perceived PB factor near $1$), and not the split by state.} 
%\footnote{AR2019 report the same near-complete na\"ivet\'e in their own specification: a perceived PB factor of $1.00$--$1.01$, with subjects anticipating at most a quarter of their PB. The split by state is new. The prediction-minus-plan gap is converted to $\beta$ with the elasticity of the act (subject bootstrap, $B=499$).}
The forecast thus captures virtually none of the realized wedge -- not even the additional gap induced by the decision-time state 
%(in the hot state, subjects anticipate only $3\%$ of it).
(in the hot state, the prediction anticipates only $3\%$ of the wedge).
%%%%%%%%%%%%%%%%%
Field data exhibit the same pattern of misprediction (Figure~\ref{fig:welfare}B).%
\footnote{The gaps, horizon profiles and updating rate in this paragraph are my computations from the replication data; unless stated otherwise, the source papers do not report them. Resolvers' belief in success and their self-rated success are on $0$--$100$ scales. \citet{carrera2022} report the over-optimism in visits ($11.5$ forecast, $7.2$ realized, without incentives). \citet{schilbach2019} describes the drivers' beliefs as somewhat overconfident but fairly accurate on average over the whole study; the gap here covers the two weeks before the commitment choice, which the test below requires.}
Forecasts exceed realized outcomes across diverse domains: by $11$ points for New Year's resolutions \citep{oscarsson2020}, $37$ for gym members' stated chance of meeting their own attendance goal \citep{carrera2022}, $10$ for the probability of working full-time past age $62$ in the Health and Retirement Study (HRS), and $28$ for the share of sober days that cycle-rickshaw drivers expect under incentives \citep{schilbach2019}.
%%%%%%%%%%%%%%%%%
%Notably, 
These %forecast 
errors are not attenuated by temporal proximity, incentives for accuracy, or experience. In the HRS, the gap does not narrow as individuals approach the target age. Smartphone users underpredict their future screen time almost as much when the forecast period is about to begin as six weeks ahead, and \citet{allcott2022} find that a larger reward for accurate prediction leaves the error level unchanged. Finally, New Year's resolvers revise their %success 
expectations by only a quarter of each month's realized surprise.

Given these welfare concerns, a possible alternative to policy intervention is self-management \citep{kruegeluhl2023}: an individual struggling with time inconsistency can voluntarily commit in advance. Yet voluntary commitment addresses the welfare gap $\Gamma$ only if those whose predictions are (most) mistaken are also those who seek commitment. 
The existing evidence is mixed: measured PB is positively associated with commitment demand in ANS2015 and \citet{kaur2015}, negatively in \citet{sadoff2020}, and largely unrelated to it in \citet{carrera2022}, where take-up even declines when awareness is raised.
The data of \citet{schilbach2019} allow a sharper test, because each driver's forecast error can be measured over the two weeks preceding the choice between sobriety incentives and unconditional payments. 
Schilbach shows that drivers' stated beliefs about their future sobriety predict demand for the incentives; among the $63$ drivers whose forecast error is measured before the choice, the correlation is $+0.20$.
Decomposing the belief into its accurate part and its error, however, I find that demand correlates positively with the former ($+0.45$ with realized sobriety) but negatively with the latter ($-0.39$; Figure~\ref{fig:welfare}C). 
Greater overestimation of sobriety is thus associated with lower commitment demand. %Consequently, drivers who most overestimate their sobriety are the least likely to choose the incentive contract.
%\footnote{The negative sign need not be causal: a driver for whom the incentives worked has both a smaller error and a stronger reason to choose them again. The conclusion rests on the absence of a positive relation.}
In four further datasets, the correlation I compute between commitment demand and the individual wedge is small and of both signs: $+0.11$ %for gym contracts 
in \citet{carrera2022}, $-0.18$ in \citet{sadoff2020}, $+0.12$ in \citet{kaur2015}, and $+0.04$ in \citet{allcott2022}, where demand for limit tools does correlate with \emph{perceived} self-control problems.
By contrast, external observers perceive the bias more accurately: \citet{fedyk2024} shows that people anticipate PB in others but not in themselves. In the two-step coordinate, subjects' actual $\beta$ over effort is $0.74$, their assessment of their own $\beta$ is $0.97$, and their belief about others' $\beta$ is $0.92$: a commitment target set by someone else on the basis of the latter belief would close $23\%$ of the gap between self-assessment and behavior. 
%\footnote{Fedyk's own structural estimates are $0.82$ (actual), $1.03$ (self) and $0.87$ (others); the counterfactual is my computation from the two-step values, with a $95\%$ interval of $[2, 49]\%$.}
Taken together, these results indicate that the decision to commit is itself made under the mistaken forecast: those who need commitment most are often the least likely to seek it. 
%In this sense, 
Thus, the %forecast 
error that defines $\Gamma$ simultaneously disables the %voluntary 
remedy that would correct it.


Three policy implications follow, independently of any normative ranking across temporal selves. 
%The criterion grants the planning self only the standing that voluntary commitment already presupposes -- the right to bind a future self through an informed choice -- and the evidence shows that this choice is often uninformed. 
%%%%%%%%%%%%%%%%%
First, voluntary take-up of commitment devices cannot reliably identify who needs them: interpreting take-up as a screening tool risks measuring who is already in control. Analogously, declining commitment is not evidence of a preference for flexibility when it rests on a forecast error. 
%%%%%%%%%%%%%%%%%
Second, the instruments that reach $\Gamma$ are those that do not depend on the person's own forecast, such as defaults, targets and schedules set by others, as well as incentives offered to everyone, which \citet{carrera2022} show to be more efficient than commitment contracts. %
%\footnote{A cohort design on the U.S. full retirement age, a population-wide increase in the return to working past $65$, points in the same direction: the plan-minus-act wedge narrows in every specification, but the biennial HRS is underpowered for the design, which is reported as exploratory in Online Appendix~\ref{app:policy}.}
%%%%%%%%%%%%%%%%%
Third, part of the wedge reflects the decision context rather than a fixed individual trait: the same people deviate further from the same plans after a block of work, while their plans and forecasts remain unchanged (Section~\ref{sec:state}). 
Properties of the choice environment -- such as the timing, the state, and the options -- %menu of available options -- 
exert a substantial influence on vulnerability to PB and realized choices. Where employers, schools, and digital platforms shape these settings, leaving them entirely unconstrained is itself a choice of environment rather than a neutral baseline. %\citep{laibson2018}
Structuring this decision context thus constitutes a powerful policy instrument.
%%%%%%%%%%%%%%%%%
Ultimately, the case for intervention is not a judgment on the person: it addresses an error that people make about themselves, under conditions that others largely set and that often supply the very cues to which the impulse responds (Section~\ref{sec:model}). 



% =====================================================================================
%\section{A framework that generates these restrictions}
\section{Microfoundations: A Multi-System Control Framework}
\label{sec:model}


While these findings do not depend on a specific structural model, collectively they motivate a unifying framework. 
PB shifts with the state, exhibits individual heterogeneity in one trade-off structure but not another, and attaches to the interaction between individual and reward type rather than to either in isolation.
This section outlines a framework that accounts for these patterns while nesting standard models as special cases. 
It is presented primarily as an organizing interpretation of the evidence: the empirical analysis tests specific restrictions implied by the framework rather than the structural model itself.
%%%
The idea that the weight attached to the present depends on the decision-maker's state dates to \citet{loewenstein1996}, and cue-triggered choice models subsequently formalized it \citep{laibson2001,bernheimrangel2004,loewensteinodonoghue2005}. 
In a related spirit, \citet{pennesi2020} derives an axiomatic quasi-hyperbolic representation in which the PB factor $\beta_\theta = 1 - \lambda(\theta)$ decreases with cue intensity $\theta$. 
The framework developed below delivers an analogous state-dependent PB factor, while incorporating the additional structure required to accommodate the empirical evidence across reward domains and trade-off structures.

In the \emph{Multi-System Control} (MSC) framework (Figure~\ref{fig:msc}), individual decision-making is governed by several parallel valuation systems, each appraising options over its own horizon.
 In the spirit of dual-process accounts \citep[e.g.,][]{thalershefrin1981,loewensteinodonoghue2005,kahneman2003},  %,loewensteinodonoghue2005,fudenberglevine2006},
these include a forward-looking, slow and costly \textit{deliberative system} $D$, and a fast, cue-triggered and myopic \textit{impulsive system} $I$.
%The state $X_t$ collects cues, visceral states, a depletable control resource, habit stocks and beliefs.
%Rather than altering fundamental valuations, the state governs the relative \textit{weight} assigned to each system. 
The state $X_t$ (cues, visceral states, depletable control, habits, beliefs) governs the relative \textit{weight} assigned to each system rather than the fundamental valuations.
A fast and automatic \textit{arbitrator} selects the action
\[
a_t=\arg\max_a \sum_i \omega_i(X_t)\, v_i(a), \qquad \omega_i(X_t)=\frac{\exp \eta_i(X_t)}{\sum_j \exp \eta_j(X_t)},
\]
where $v_i(a)$ is the value that system $i \in \{D, I, \dots\}$ assigns to action $a$ and the weights $\omega_i$ depend on the state.
An \textit{executive control} process can increase the deliberative weight $\omega_D$ by drawing on a finite control resource. It is recruited only when a monitoring system -- whose precision degrades under cognitive load, visceral heat, and fatigue -- judges that the benefits of control outweigh its costs.
Beyond immediate choices, the executive can act dynamically by modifying the state under which a subsequent decision takes place -- the core mechanism of commitment.
%%%%%%%%%%%

\begin{figure}[H]\centering
\caption{\footnotesize%
The MSC --- state-dependent weights combine parallel valuations into one action and determine PB.}
\includegraphics[width=.99\linewidth]{d09_msc.pdf}
\label{fig:msc}
\end{figure}
% fig:msc is drawn in figs/diagrams/d09_msc.tex (TikZ; pdflatex d09_msc); 10-01 wide version (old: figs/diagrams/_backups/d09_msc_before_1001.tex)


\clearpage
\noindent
This architecture is consistent with neuroeconomic evidence: immediate rewards predominantly engage limbic valuation areas (the analogue of $I$), whereas the lateral prefrontal cortex (closer to $D$) is engaged by all intertemporal choices \citep{mcclure2004}.
Crucially, \citet{figner2010} show that disrupting the left lateral prefrontal cortex raises choices of immediate rewards while leaving choices between two delayed rewards unchanged -- mirroring the decision-time state effect identified in Section~\ref{sec:state}.
Whether these neural dynamics reflect active inhibitory control or integrated value comparison remains an active debate \citep{bulley2022}.



%\noindent
%PB can be derived as an equilibrium outcome of state-dependent competition between parallel valuation systems.
PB follows from the state-dependent weighting of parallel valuation systems.
For instance, in an intertemporal trade-off where the impulsive system $I$ evaluates only the immediate payoff, the arbitrator ranks options under an effective PB factor:
\[
\beff(X)=\frac{\omega_D(X)}{\omega_D(X)+\omega_I(X)\,\pi}\in(0,1],
\qquad \pi\equiv\frac{\text{$I$'s value of the immediate component}}{\text{$D$'s value of it}}.
\]
When the impulsive weight is zero ($\omega_I(X)=0$), $\beff(X)=1$ and decision-making reduces to standard exponential discounting; conversely, as cues, visceral heat, or temporal proximity raise $\omega_I(X)$, $\beff(X)$ decreases.
%%%%%
An evolutionary perspective suggests why such a fast, cue-driven system exists: self-control failures can arise from well-learned strategies that are adaptive in the broader settings in which they evolved, for instance where waiting carries a risk of losing the reward \citep{hayden2018}. However, they misfire in modern environments that manufacture such cues faster than genes can adapt \citep{burnhamphelan2020desire}.

Standard models can be represented as special cases of this framework.
With a single system and no control cost, the agent discounts exponentially. 
Under two systems whose weights are invariant to the state, $\beff$ is constant and the framework reduces to quasi-hyperbolic discounting (the special case that Sections~\ref{sec:state}--\ref{sec:person} reject). 
A single self-control cost, common to all temptations, corresponds to the dual-self model of \citet{fudenberg2006}, and a habit stock with a cue-triggered impulsive weight to the addiction model of \citet{bernheimrangel2004}.%
\footnote{
The MSC framework differs from the dual-self model on two testable margins. First, the single willpower cost of dual-self models implies one $\beta$ per person, which the evidence of Section~\ref{sec:person} rejects. Second, a cue that raises salience without changing resources moves choices under this framework, but leaves them unaffected under a pure willpower model.} %existing micro-randomized trials cannot settle this margin, because a reminder also moves a time-consistent agent.}
Relative to \citet{pennesi2020}, incorporating the executive, the monitor, and the domain-specific allocation of control yields predictions that a scalar cue-indexed weight does not generate: the distortion attaches specifically to the immediate component of a tempting reward, and the PB factor can vary across individuals in one trade-off structure while remaining uniform in another.
Finally, the state-dependent factor that the framework implies is characterized over budgets in Section~\ref{subsec:RPtheory2}, so that it has empirical content on the data the literature collects.

The framework further provides an interpretive lens for the monetary findings in Section~\ref{sec:person}.
As a relatively recent innovation in evolutionary terms, liquid cash possesses no intrinsic visceral appeal; PB over monetary payoffs instead reflects what funds can purchase at the time of choice.
Where cash implies no immediate consumption -- as may be the case for the university subjects in ANS2015 -- the model predicts little PB over monetary rewards and minimal interpersonal dispersion. Conversely, where liquidity offers immediate access to visceral goods -- as may occur among the adolescent smokers and drinkers in CTW2022 -- the MSC architecture is consistent with stronger PB over cash than over food.
Findings that conflict under a commodity taxonomy may thus be reconciled within a single state-dependent model.
\bigskip

%Three predictions find no support in the data: the effect of executive load is a bounded null, appetite does not move PB over food relative to money in the one design that can test it, and a prediction about subjects already engaged in a task was not supported in the form in which it could be tested. For these reasons, the claims of this paper rest on the randomized evidence rather than on the framework.



% =====================================================================================
\section{Conclusions}
\label{sec:limits}

PB is often treated as a property of the person: a single $\beta$ that persists within an individual across goods, contexts and moments. 
%This paper has tested that view across its implied invariances, finding that it fails along every observable dimension.
This paper has tested that view against the invariances it implies: it fails across states, trade-off structures and goods, while relative individual rankings remain stable within a single structure.
Within the same individuals, a randomly assigned state at the moment of choice moves the act and leaves the plan unchanged. Across people, individual PB varies over effort but not over money. 
In a protocol that varies only the reward, money is the commodity exhibiting the strongest PB -- yet this distortion attaches to the person--reward pair rather than to either component in isolation.
Uncovering these patterns required a valid testing framework: statistics based on the level of the wedge, conventionally used to document heterogeneity in $\beta$, cannot distinguish genuine state dependence from stochastic variation under a state-independent model. In contrast, design-recentered tests of equality across assigned states successfully separate the two.
The welfare consequences follow: because the state moves the act while leaving the plan and the forecast unchanged, most of the wedge is unforeseen, and the decision to commit, guided by that forecast, does not track the error. The MSC framework derives this reaction function from state-dependent weights on parallel valuation systems and nests quasi-hyperbolic and dual-self models as restricted cases.

These findings and their limitations point toward actionable avenues for future experiments.
The state variation in AR2019 is also consistent with a transient shock to task costs. The main qualitative lesson holds under either interpretation -- the aggregate $\beta$ averages across states -- but investigating the underlying mechanism in greater depth requires a purpose-built design.
Moreover, the results on goods rest on a single protocol in which money always appears in the same position. 
Finally, relying on existing datasets naturally restricts the observable range of states and commodities. 
%%%
These considerations motivate three concrete design features.
First, factorially combining a randomized consumption cue with dated trade-offs over effort and money would cleanly separate a state-dependent $\beta$ from a transient cost shock, while testing the MSC framework against a pure willpower model.
Second, randomizing the display position of each reward and systematically varying the consumption opportunities afforded by cash at fixed dates would isolate the option value of money from the good itself. Finally, repeated elicitations for each individual within each trade-off structure would distinguish true heterogeneity from measurement imprecision.

%These conclusions have limits. The state randomized in the canonical panel is the state of having just worked, rather than a consumption cue, and its effect on immediate choices is also consistent with a transient increase in the cost of the next tasks (\textcolor{red}{AF info, riguardo alla cosa di alternative reading for finding 1 di state hold-cot: non lo sto più citando sopra.}); the published $\beta$ averages over states under either reading, but separating the two requires a different design. Several margins yield bounded nulls at the intensities available in existing data. The results on goods rest on a single protocol in which money always appears in the same position, and the concentration of the money--food gap among drinkers and smokers does not survive a correction for the number of characteristics examined. Finally, all the evidence comes from public data collected for other purposes. This has an advantage, since the main randomization was designed and run by authors who interpreted it differently, but it limits the states and goods that can be studied.
%These limits point to specific designs. Crossing an intense consumption cue, assigned in random order through a published assignment process, with dated trade-offs over both effort and money would separate a state-dependent $\beta$ from a transient cost, and the framework of Section~\ref{sec:model} from a pure willpower model. Randomizing the position of each good, and varying what money can buy at fixed amounts and dates, would separate the option value of cash from the good itself. Repeated elicitations for each person within each trade-off structure would distinguish the absence of heterogeneity from the absence of precision.


Ultimately, these findings reframe how estimated discount factors should be interpreted. A single estimate of $\beta$ reflects not an immutable trait of the person, but a conditional summary of a state-dependent reaction function, averaged over the support of states, commodities, and trade-off structures sampled in a given dataset. 
In this view, PB is less a fixed property of who people are than a systematic response to the context in which they choose.
This reinterpretation underscores the critical role of targeted policy design.
Across domains as diverse as retirement saving, health behaviors, digital engagement, and substance use, present-biased choices can be amplified by an evolutionary mismatch between ancestral decision heuristics and modern environments engineered to capture attention and exploit impulse.
A static personal trait invites identification and targeting; a state-dependent reaction function, by contrast, calls for attention to the settings in which decisions are made. 
%A static personal trait can only be identified or targeted; a state-dependent reaction function, by contrast, must be accounted for by designing choice settings that are mindful of decision timing, present cues, and choice context.
Transcending normative debates on whether to fix or penalize the individual, the design of choice environments provides a directly actionable margin for protecting long-run welfare while respecting human agency.
% AF, addressing claude's comment on welfare: I'm not putting the self before, I'm following the sconservative criterion of Section 7.2. Proprio perché lo stato muove l'atto ma non la previsione, lo stato è la causa diretta del misprediction error non anticipato ($\Gamma = e_{act} - e_{pred}$). AF: maybe the objection of Claude was only on the word "long-run", that makes it a bit ambiguous - but I think it's fine, just to put it together / in constras with "human agency".




\iffalse
{\tiny
\medskip
\textcolor{red}{AF; sign of where I arrived: from here onwards, I didn't read (and didn't reword) yet.}
%https://chatgpt.com/c/6ab2edf4-2128-83eb-a76e-7dec6fb69c04
%https://notebook.google.com/notebook/9d2c72ae-54d8-44e0-9d00-9a5dd80f7f99?authuser=1

\textcolor{red}{AF: from here, should still ask claude to review/check that everything is correct, as I edited it manually following LM's insights (which may be wrong in content and/ord grammar/style)}
}
\fi

\iffalse
\begin{figure}[H]\centering
\caption{\footnotesize%
The argument in one picture.}
\includegraphics[width=.99\linewidth]{d12_roadmap_1007.pdf}
\label{fig:roadmap}
\end{figure}
\fi

\clearpage

\iffalse % smaller
{\fontsize{8}{9.6}\selectfont
\setlength{\bibsep}{1.5pt}
\begin{multicols}{2}
\interlinepenalty=10000
\bibliography{../refs}
\end{multicols}
}
\fi

\section*{References}

{\footnotesize
\setlength{\bibsep}{1.5pt}
\renewcommand{\bibsection}{}
\begin{multicols}{2}
\interlinepenalty=10000
\bibliography{../refs}
\end{multicols}
}

% =====================================================================================
\clearpage
\appendix
\renewcommand{\thesection}{\Alph{section}}
\numberwithin{table}{section}
\numberwithin{figure}{section}
\renewcommand{\thetable}{\Alph{section}.\arabic{table}}
\renewcommand{\thefigure}{\Alph{section}.\arabic{figure}}
\begin{center}
{\Large\textbf{Online Appendix}}\\[4pt]
{\large Beyond the Trait: Present Bias as a Reaction Function}\\[2pt]
Alessandro Facchini --- 10 October 2026
\end{center}
\medskip
\noindent This appendix contains the proofs (Appendix~\ref{app:proofs}); the recentered state test in full
(\ref{app:recentered}); the robustness of the state result and the other state margins (\ref{app:robust}); the
measurement-error model (\ref{app:me}); the three-good protocol (\ref{app:goods}); the magnitudes, the household
panel and the type shares (\ref{app:magn}); the tests used throughout, with their Monte Carlo certificates
(\ref{app:inference}); two exploratory policy designs (\ref{app:policy}); and the datasets and the code
(\ref{app:data}). A longer version of the paper with every specification is available on request.

% -------------------------------------------------------------------------------------
\section{Proofs}
\label{app:proofs}

\paragraph{Proposition~\ref{prop:sqh2}.} \emph{Necessity.} Interior optimality gives
$D_{c(k)}(t)u'(x^{k}_{t})=\lambda^{k}p^{k}_{t}$; at a corner the condition holds with inequality, which only
strengthens what follows. Along a balanced sequence with $x^{k_i}_{t_i}>x^{k'_i}_{t'_i}$ for every $i$, concavity
gives $1\ge\prod_i u'(x^{k_i}_{t_i})/u'(x^{k'_i}_{t'_i})$. Balance cancels the multipliers, and collecting the
discount terms the product equals $\delta^{-\tau}\prod_c\beta(c)^{n^-_c-n^+_c}$ times the price product, as in
\eqref{eq:chainRP2}. Condition~(ii) sets every exponent of $\beta$ to zero and condition~(i) makes the $\delta$ factor
at least one, so the price product is at most one.

\emph{Sufficiency.} Writing marginal utilities as $v^{k}_{t}=\lambda^{k}p^{k}_{t}/D_{c(k)}(t)$, rationalization
requires numbers $\log\lambda^{k}$, $\log\beta(c)$ and $\log\delta\le0$ such that $v$ is non-increasing in quantity,
that is, whenever $x^{k}_{t}>x^{k'}_{t'}$,
\[
\log\lambda^{k}-\log\lambda^{k'}
-\mathbf 1\{t>0\}\log\beta\big(c(k)\big)+\mathbf 1\{t'>0\}\log\beta\big(c(k')\big)
-(t-t')\log\delta
\;\le\;\log p^{k'}_{t'}-\log p^{k}_{t}.
\]
By a theorem of the alternative, this linear system is infeasible if and only if a nonnegative integer combination
of its rows eliminates every $\log\lambda^{k}$ and every $\log\beta(c)$, leaves a nonpositive multiple of
$\log\delta$, and has a negative right-hand side. Eliminating the multipliers is balance, eliminating each
$\log\beta(c)$ is condition~(ii), the sign on $\log\delta$ is condition~(i), and a negative right-hand side is
$\prod_i\rho_i>1$. An infeasible system therefore exhibits a balanced sequence that satisfies (i) and (ii) and
violates the downward-sloping demand property (integer weights for rational prices, irrational prices by
approximation, as in \citet{eis2020}). SAR-SQH excludes such sequences, so the system is feasible, and Afriat's
construction on the resulting supergradients delivers an increasing concave $u$. For fixed $(\beta(\cdot),\delta)$
the system reduces to difference constraints in the multipliers, which are feasible if and only if the associated
weighted graph has no negative cycle; this is how the program is computed and certified on synthetic agents
(Appendix~\ref{app:mc}, block T15).

\emph{Nesting.} The pooled count neutralizes a single $\beta$ and gives \citet[Theorem~2]{eis2020}'s axiom; deleting
the count requires $\beta=1$. Per-cell equality implies pooled equality, so SAR-SQH restricts a subset of the
sequences that SAR-QHD restricts and is the weaker axiom. Ties impose no constraint, since a concave $u$ may have a
kink at the common quantity; imposing one would make the test reject far too often. \qed

\paragraph{Proposition~\ref{prop:rung0}.} Holding $(u,\delta)$ fixed, each cell statistic depends only on that
cell's factor and is strictly monotone in it ($m_c=b\log\beta_c$ in the log-wedge coordinate). The map from
$(\beta_1,\dots,\beta_C)$ to $(m_1,\dots,m_C)$ is therefore a bijection onto its range, with as many parameters as
moments, and every observed vector in the range is fitted exactly. \qed

\paragraph{Proposition~\ref{prop:class}.} (i) With $g_k=\bar g+\eta_k$ and $\xi\perp\eta\mid w$,
$\E[z_kg_k]=\E\big[\E[z_k\mid w]\,\E[g_k\mid w]\big]=0$ because $\E[z_k\mid w]=0$; nothing is assumed about the
dependence of the $\eta_k$ across observations. (ii) Let $z_k=h_k(\xi,w)$ with $m_k\equiv\E[h_k\mid w]$ not almost
surely zero. If $m_k$ is a nonzero constant, $\E[T]=m_k(\bar g+\E[\eta_k])\ne0$ under any Rung-1 null with PB and a
mean-zero error. Otherwise take $\eta_k=c\,(m_k-\E[m_k])$, a function of $w$ alone: $\xi\perp\eta\mid w$ holds and
every Rung-1 null is respected, while $\E[T]=\bar g\,\E[m_k]+c\Var(m_k)$, which differs from zero for every $c$ but
one \citep[Proposition~1]{borusyakhull2026}. (iii) With $z_k\equiv1$, $\E[T]=\bar g+\E[\eta_k]$. \qed

\paragraph{The feasibility share.} Since $\Var(s)=\E[\Var(s\mid w)]+\Var(\E[s\mid w])$, the recentered instrument
$s-\E[s\mid w]$ has variance $\rho\Var(s)$, with $\rho=\E[\Var(s\mid w)]/\Var(s)$. When $\rho=0$ the recentered
instrument is identically zero and every member of the valid class is the zero statistic. In the effort panel
$\rho=0.81$ for the assigned state at the day $\times$ composition stratum; in the three-good protocol $\rho=0$ for
the position of money, which is always second, and $\rho=1$ for each food.

% -------------------------------------------------------------------------------------
\section{The recentered state test in full}
\label{app:recentered}

\paragraph{The expected state.} The $614$ subject-days of \citet{augenblick2019} contain $2{,}245$ decision sets.
Six day compositions occur, the full day with today's set, two future sets and two prediction sets being the most
common ($313$ subject-days). The expected state $\mu_k$ of a decision is the share of subject-days in its stratum in
which the decision's set came after the mandatory block. At the baseline stratum, the calendar day of the experiment
crossed with the day's composition ($13$ strata), it is $0.46$ for today's set, $0.60$ and $0.83$ for the first and
second future sets, and $0.59$ and $0.83$ for the two prediction sets. The expected state explains $19\%$ of the
variance of the raw indicator at this stratum, $18\%$ at the composition alone, $22\%$ with the session added and
$33\%$ with the slot. Conditioning on the slot is over-conditioning, since the slot is an outcome of the draw rather
than a feature of the day. The assignment process is not published, so $\mu$ rests on the exchangeability of the
realized orderings within a stratum, which the falsification test below examines.

\paragraph{Estimators.} The recentered instrument is $s_k-\mu_k$, with its interaction with the today indicator;
``control for $\mu$'' adds $\mu_k$ and $\mu_k\times$today to the regression and instruments with the raw state; the
reduced form regresses the outcome on the recentered instruments. In the control-function regression on all $100$
subjects, the coefficients on $\mu$ and $\mu\times$today are $+0.14$ and $-0.19$: exposure predicts the outcome
directly, which is the channel the adjustment removes. Intervals and $p$-values in Table~\ref{tab:rcstrata} come from
the subject bootstrap ($B=499$).

\begin{table}[H]\centering
\caption{$\Delta\log\beta$ across the assigned state by stratum and estimator; $p$ in parentheses.}
\label{tab:rcstrata}
{\footnotesize
\begin{tabular}{llcccc}
\toprule
Sample & Stratum & unadjusted & recentered IV & control for $\mu$ & reduced form \\
\midrule
all $100$ & day $\times$ composition & $-0.180$ ($.008$) & $-0.122$ ($.12$) & $-0.122$ ($.11$) & $-0.121$ ($.12$) \\
all $100$ & composition & $-0.180$ ($.008$) & $-0.117$ ($.12$) & $-0.118$ ($.14$) & $-0.117$ ($.12$) \\
all $100$ & day $\times$ composition $\times$ slot & $-0.180$ ($.008$) & $-0.097$ ($.30$) & $-0.097$ ($.30$) & $-0.097$ ($.30$) \\
keep $72$ & day $\times$ composition & $-0.191$ ($<.004$) & $-0.161$ ($.02$) & $-0.162$ ($.016$) & $-0.161$ ($.02$) \\
keep $72$ & composition & $-0.191$ ($<.004$) & $-0.136$ ($.056$) & $-0.136$ ($.056$) & $-0.136$ ($.056$) \\
keep $72$ & day $\times$ composition $\times$ slot & $-0.191$ ($<.004$) & $-0.140$ ($.13$) & $-0.139$ ($.13$) & $-0.139$ ($.13$) \\
\midrule
\multicolumn{6}{l}{\emph{with distance to the work date as a control (day $\times$ composition)}} \\
all $100$ & & $-0.182$ ($<.004$) & $-0.119$ ($.12$) & $-0.119$ ($.13$) & \\
keep $72$ & & $-0.192$ ($<.004$) & $-0.161$ ($.047$) & $-0.160$ ($.054$) & \\
\bottomrule
\end{tabular}}
\end{table}

\paragraph{Randomization inference and falsification.} Under the sharp null that the state moves no decision, every
reordering of a day's components within its stratum leaves the choices unchanged. The orderings are redrawn $999$
times, the state and its recentered version are recomputed with the choices held fixed, and the $p$-value is the
share of redraws with an estimate at least as large in absolute value as the observed one. The same redraws give the
falsification test of \citet[Section~3.5]{borusyak2023}: the recentered state must be orthogonal, within subject,
to every predetermined variable (Table~\ref{tab:rcfals}). The joint statistic is the sum of squared fitted values from
a regression of the recentered state on the five variables, compared with its distribution over the redraws.

\begin{table}[H]\centering
\caption{Falsification: covariance of the recentered state with predetermined variables, within subject, against its
distribution over $999$ redraws; the raw indicator for comparison.}
\label{tab:rcfals}
{\footnotesize
\begin{tabular}{lcccc}
\toprule
Variable & cov$(s-\mu,r)$ & null s.d. & RI $p$ & cov$(s,r)$, raw \\
\midrule
log wage drawn & $-0.002$ & $0.004$ & $.58$ & $+0.002$ \\
distance to the work date & $+0.085$ & $0.035$ & $.011$ & $+0.550$ \\
net distance to the payment date & $-0.029$ & $0.048$ & $.53$ & $+0.108$ \\
decision day & $-0.008$ & $0.005$ & $.16$ & $-0.104$ \\
the expected treatment $\mu$ & $-0.000$ & $0.000$ & $.82$ & \\
joint (sum of squared fitted values) & & & $.28$ & \\
\bottomrule
\end{tabular}}
\end{table}

\noindent The residual imbalance in the distance to the work date reflects which future dates a day's decision sets
refer to, a feature of the protocol that day and composition do not capture; distance therefore enters as a control
in the last two rows of Table~\ref{tab:rcstrata}, with little effect on any estimate. Exact randomization inference
for the total effect of the state on today's decisions gives $p=.018$ on all subjects ($-0.077$ against a null
standard deviation of $0.034$) and $p\le.001$ on the source sample ($-0.135$ against $0.041$); for the interaction
alone, $p=.091$ and $.021$; for the plan, $.73$ and $.30$. On all $100$ subjects and all $11{,}405$ decisions the cross-vantage test gives present $-0.072$
under both estimators, future $+0.020$ unadjusted and $-0.011$ recentered, and present minus future $-0.092$
($p=.004$) and $-0.061$ ($p=.064$); in the coordinate $\log(e+20)$ on the source sample the recentered contrast is
$-0.095$ ($p=.013$). The wage-quintile shares are flat across components, days, vantages and states ($\chi^2$ $p$
from $.40$ to $.87$), so the stakes margin needs no recentering.

% -------------------------------------------------------------------------------------
\section{Robustness of the state result, and the other margins}
\label{app:robust}

\paragraph{The matched-state estimate.} Table~\ref{tab:robmatched} reports the estimate under every specification
tried. In this panel the recorded effort variable counts the ten mandatory tasks, so $\log e$ is the source paper's
cost coordinate $\log(e+10)$ and a choice of no additional tasks is not dropped; the $21.8\%$ of choices at the
ceiling of $100$ tasks are censored, and the source paper's own Tobit analysis shows that handling the censoring
strengthens the before/after contrast.

\begin{table}[H]\centering
\caption{The matched-state estimate under every specification; within-subject fit in the source paper's cost
coordinate.}
\label{tab:robmatched}
{\small
\begin{tabular}{>{\raggedright\arraybackslash}p{7.0cm}cc}
\toprule
Specification & $\Delta\log\beta$ & $p$ \\
\midrule
All $100$ subjects, subject FE & $\mathbf{-0.180}$ $[-0.310,-0.062]$ & $.008$ \\
Keep-list ($72$), subject FE & $-0.191$ $[-0.326,-0.060]$ & $<.004$ \\
$+$ set order & $-0.200$ $[-0.328,-0.088]$ & $<.004$ \\
$+$ set order $+$ decision day & $-0.159$ $[-0.290,-0.055]$ & $.016$ \\
$+$ set order $+$ distance to work & $-0.204$ $[-0.334,-0.093]$ & $<.004$ \\
Coordinate $\log(e+20)$ & $-0.191$ $[-0.318,-0.085]$ & $<.004$ \\
Source paper's exact estimation sample & $-0.235$ $[-0.385,-0.088]$ & $<.004$ \\
Within subject-\emph{day} ($157$ cells) & $-0.112$ $[-0.258,+0.039]$ & $.153$ \\
\bottomrule
\end{tabular}}
\end{table}

\paragraph{Nonparametric brackets by state.} On the $96$ subjects observed in both states and at both vantages, each
today-choice is matched into the subject-averaged advance schedule of the same state. At $K=4$ wage cells the cold
bracket is $[0.706,1.000]$ against $[0.589,0.706]$ hot; at $K=5$, $[0.787,1.000]$ against $[0.707,0.787]$; at $K=6$
the cold bracket is empty, the familiar small-sample failure of monotone matching, in which one cell mean out of order
voids the intersection. Once each cell mean carries its sampling error the widened brackets overlap, so the brackets
corroborate the direction of the state contrast without establishing it.

\paragraph{Stakes.} Flatness of $\beta$ across wage cells is rejected at terciles ($p=.0050$), quartiles ($.0039$),
quintiles ($.00004$) and sextiles ($<10^{-5}$), and across specifications at quintiles: $.00004$ on the $72$ subjects
with all choices, $.00007$ on all $100$, $.030$ on interior choices only, and $.0021$ on interior choices with
work-date fixed effects. A $\beta$ that decreases with the stake is rejected at every binning ($p\le.0015$); an
increasing one is not rejected at terciles ($p=1.00$) or quartiles ($.29$) but is at quintiles ($.0005$), where a dip
in the $7$--$12$ cent cell breaks it, and dropping that cell lifts flatness to $p=.109$. The $72$ leave-one-subject-out
refits give $p$ between $4\times10^{-6}$ and $1.5\times10^{-4}$.

\paragraph{Proximity.} Among future-dated decisions, the allocation to a given work date moves as the date
approaches. The comparison has to be made within work date: with subject and work-date fixed effects on interior
choices, the wage-adjusted log allocation by quartile of distance, from nearest to farthest, is $-0.042$, $-0.043$,
$+0.035$ and $+0.045$, and equality is rejected ($\chi^2_3$, $p=.0004$). Without the work-date effects the comparison
runs across calendar targets and the gradient disappears ($p=.39$). An allocation that rises with distance is not
rejected as a monotone profile ($p=.97$); one that falls is ($p<.001$).

\paragraph{Hunger.} \citet{marchgluth2025} induce hunger within person: $70$ participants, who come to the laboratory after
an overnight fast, choose between pairs of foods once without eating and once after a protein shake matched to their
daily caloric needs, in counterbalanced order.
In the $10{,}595$ trials in which one item is tastier and the other healthier, taste wins $61.4\%$ of the time when
sated and $66.7\%$ when hungry; the within-subject shift is $+4.5$ percentage points $[0.3,8.8]$ ($t=2.1$, Cox--Shi
$p=.035$), and $56\%$ of participants move in that direction. A food choice trades an immediate reward, taste,
against a delayed cost, health, so the share of conflicts won by the immediate attribute is an analogue of $\beff$
rather than $\beff$ itself.

\paragraph{Intoxication.} In \citet{schilbach2019}, $475$ cycle-rickshaw drivers in Chennai are breathalyzed daily
($5{,}342$ driver-days). At the choice interviews, drivers who are intoxicated choose sobriety incentives over
Rs~$90$ of unconditional pay $43\%$ of the time, against $81\%$ when sober ($-0.38$, $t=-8.0$). Within driver, among
the $27$ drivers observed both sober and intoxicated at a choice interview, the effect is $-0.11$ $[-0.34,+0.11]$
($t=-1.0$), and $-0.13$ and $0.00$ against the alternatives of Rs~$120$ and Rs~$150$. The cross-sectional gap
describes who is intoxicated at an interview rather than what intoxication does.

\paragraph{Executive load.} \citet{bessone2021} measure present bias with real-effort choices in the design of
\citet{augenblick2019} among $425$ low-income workers in Chennai, together with objectively measured nightly sleep
($3.4$ to $8.0$ hours). By sleep quartile, $\beff$ is $0.89$, $0.88$, $0.85$ and $0.95$; the slope is $+0.016$ per
hour ($t=0.9$) and the joint test of a departure from monotonicity gives $p=.36$. The randomized night-sleep treatments
do not move present bias either, as the source paper reports. At the doses these data reach, the executive-resource
channel is a bounded null.

\paragraph{Two restrictions that survive.} The same machinery does not reject two restrictions that the theory
imposes. First, the two-step divides one $\theta$ by one $b$, which presumes that the wage elasticity of effort is
the same at both vantages and in both states. Estimated separately, $b=0.668$ $[0.604,0.732]$ on future-dated choices
and $0.633$ $[0.569,0.695]$ on today's, and $0.657$ cold against $0.654$ hot; in the joint model the interactions of
$\log W$ with the vantage, with the state and with both are $-0.002$ ($p=.93$), $+0.016$ ($p=.67$) and $-0.065$
($p=.16$), so the two denominators the two-step would use are $0.658$ and $0.656$, a ratio of $0.997$. Second, a
$(\beta,\delta)$ agent values a dated payment by its amount and its date, so the medium of payment cannot move the
intertemporal switching point. \citet{hoel2022} randomize payment in cash or through a mobile payment application
among $338$ participants. The switching row in a seven-row price list moves by $+0.29$ raw and by $+0.52$
$[-0.11,+1.14]$ once the $37\%$ of participants who never switch are treated as censored (rank test $p=.14$), which
excludes effects larger than about a sixth of the elicited range. The risk task of the same experiment does respond
to the medium ($+0.58$ $[0.09,1.06]$, rank $p=.03$), so the null is specific to the intertemporal margin.

\paragraph{A prompt assigned at the moment of choice.} Two micro-randomized trials assign a prompt within person, with
known probability, at each decision point. In the Drink Less trial \citep{bell2023}, an evening notification
($10{,}470$ decision points, $349$ users) raises the probability that users record their drinking in the app by $9.1$
percentage points on a base of $3.6$ ($t=11.2$), and the effect is no smaller in the second half of the trial than
in the first ($+9.1$ against $+8.0$). In HeartSteps \citep{klasnja2019}, the
average effect of an activity suggestion on steps in the next thirty minutes is positive but imprecise ($+3.9\%$,
$t=0.7$, $37$ participants), and it is concentrated in moments when the user has been inactive ($+26.5$ steps,
one-sided Cox--Shi $p=.026$, although the two-sided bootstrap interval $[-1.6,+59.7]$ includes zero). The trials show
that behavior on these margins moves with a prompt at unchanged resources, the margin on which the framework of
Section~\ref{sec:model} departs from a pure willpower model. They are not tests of $\beta$: the outcomes are not dated
trade-offs, and a time-consistent agent with imperfect memory also responds to a reminder \citep{ericson2017}.

% -------------------------------------------------------------------------------------
\section{The measurement-error model}
\label{app:me}

\paragraph{Individual estimates.} In \citet{ans2015} each subject divides work between two work dates a week apart at
five task rates, deciding both a week in advance and on the earlier work date, and in separate convex-time-budget
choices divides money between two payment dates at five interest rates, the earlier date being either the day of the
decision or a later one. In each domain the log ratio of the two allocations is regressed on subject $\times$ vantage
effects, one common elasticity to the log task or interest rate and, over money, the delay between the two dates; the
difference between a subject's two vantage effects is the subject's today-effect $\hat\theta_i$, and
$\log\hat\beta_i=\pm\hat\theta_i/\hat b$ (the sign reflects that effort is a cost). This is the restriction of
\citeauthor{ans2015}'s own individual analysis, one curvature per domain, and it reproduces their mean individual
factors ($0.916$ and $0.984$, against their $0.91$ and $0.99$). The common elasticities are $\hat b=-1.15$ over effort
and $-15.9$ over money. The sampling variance of $\hat\theta_i$ is the subject's own residual variance times
$1/n_{i1}+1/n_{i0}$, the numbers of choices at the two vantages, plus the contribution of the common coefficients
(clustered by subject); dividing by $\hat b^2$ gives $V_{id}$. A subject needs both vantages, three rates and eight
decisions in a domain, and $96$ subjects have estimates in both. In \citet{cheung2022}, $\hat r_{id}$ is the logarithm
of the published individual estimate of $\beta$ over good $d$ and $V_{id}=(\mathrm{se}_{id}/\hat\beta_{id})^2$, for
$0.05<\hat\beta_{id}<5$ and $V_{id}<1$ ($409$, $413$ and $389$ subjects in the three pairs of goods).

\paragraph{Likelihood.} With $\hat r_i=(\hat r_{i1},\hat r_{i2})'$ and $V_i=\mathrm{diag}(V_{i1},V_{i2})$, the model
$\hat r_i=r_i+\nu_i$, $r_i\sim N(\bar r,\Sigma)$, $\nu_i\sim N(0,V_i)$, with $\nu_i$ independent of $r_i$ and all
terms independent across subjects, gives $\hat r_i\sim N(\bar r,\Omega_i)$ with $\Omega_i=\Sigma+V_i$, and the
log-likelihood
\[
\ell(\bar r,\Sigma)=-\tfrac12\sum_{i}\Big[\log\det\Omega_i+(\hat r_i-\bar r)'\,
\Omega_i^{-1}(\hat r_i-\bar r)\Big].
\]
$\Sigma$ is parametrized by $(\log\sigma_1,\log\sigma_2,\operatorname{artanh}\rho)$, which keeps it positive
semidefinite, and $\ell$ is maximized by Nelder--Mead from several starting values. An estimate with a large $V_{id}$
has a large $\Omega_i$ and therefore a small weight, so imprecisely estimated subjects remain in the sample and lose
influence. The error covariance is diagonal because the two estimates come from separate choices. In
\citet{cheung2022} the three goods are elicited in the same sittings, so a shock common to the occasion enters
$\Sigma$ and raises the estimated correlations, the caveat stated in Section~\ref{sec:person}.

\paragraph{Estimates and inference.} On \citet{ans2015} the mean true factors are $0.876$ over effort and $0.987$
over money, the aggregate gap of Section~\ref{sec:person}. A true standard deviation of zero lies on the boundary of
the parameter space: fixing $\sigma_d=0$ also removes the covariance, so the likelihood-ratio statistic for no
dispersion has the null distribution $\tfrac12\chi^2_1+\tfrac12\chi^2_2$ \citep{stramlee1994}. Over effort the true
standard deviation of $\log\beta$ is $0.237$, no dispersion is rejected (likelihood ratio $22.6$, $p<10^{-5}$), the
profile-likelihood $95\%$ interval is $[0.160,0.320]$, and a subject bootstrap gives $[0.067,0.375]$. Over money it
is $0.003$, the likelihood ratio for no dispersion is $0.05$ ($p=.89$), and the profile interval is $[0,0.038]$. With
$\sigma_2\approx0$ the correlation $\rho$ is not identified and sits at a boundary; the two individual estimates are
uncorrelated ($+0.030$ $[-0.096,+0.145]$, subject bootstrap), as \citeauthor{ans2015} find for their own individual
estimates ($-0.05$). Read at one elasticity per domain, a difference in today-effects is a difference in $\beta$;
subjects whose own effort responds steeply to the task rate show little of the dispersion over effort, so part of it
may reflect differences in curvature, which a design with more price points per subject could separate.

\paragraph{Robustness.} Leaving out any one subject keeps the effort standard deviation between $0.178$ and $0.242$
and the money one below $0.005$. With the pooled residual variance in place of each subject's own, the two are
$0.219$ ($p=.009$) and $0.021$; with Taylor-free variances from a within-subject bootstrap of each today-effect,
$0.251$ and $0.036$; without the least precise quartile of effort estimates, $0.260$ and $0.017$. A moment check that
uses no likelihood agrees on the money side: in the most precise half of money estimates their standard deviation is
$0.059$, below the median sampling standard deviation of $0.077$. The design would have detected a real dispersion
over money: in $200$ samples simulated at the subjects' own sampling variances, the test rejects a true zero in $3.5\%$
of samples at the $5\%$ level, a true standard deviation of $0.05$ in $84\%$, and one of $0.10$ in every sample.

\paragraph{Dispersion and correlation within a structure.} On \citet{cheung2022} both members of every pair have real
dispersion (true standard deviations from $0.199$ to $0.232$), and the true correlations are $+0.92$ $[0.81,1.00]$,
$+0.89$ $[0.76,1.00]$ and $+0.71$ $[0.55,0.87]$ for money--healthy, money--unhealthy and healthy--unhealthy food
(subject bootstrap, $399$ resamples), against Spearman correlations of $0.47$ to $0.60$ between the raw individual
estimates reported by the source paper. No subject drives them: dropping the subject furthest from the fitted
distribution, the furthest $1\%$ or $5\%$, or the five most precise estimates keeps every true standard deviation
between $0.15$ and $0.24$ and every correlation between $0.70$ and $0.99$.

% -------------------------------------------------------------------------------------
\section{The three-good protocol}
\label{app:goods}

With subject fixed effects only, the today-effects are $-0.291$, $-0.250$, $-0.235$ (money, healthy, unhealthy);
with subject $\times$ good fixed effects, $-0.341$, $-0.222$, $-0.213$. The second is the specification the argument
needs, since the today indicator is the session indicator: with subject effects only, each good's today-effect also
absorbs that good's first-session level. At a common elasticity $b=1.514$ the implied $\beta$'s are $0.799$, $0.863$
and $0.869$; at the good-specific elasticities $1.713$, $1.398$ and $1.425$ they are $0.820$, $0.852$ and $0.861$.

\paragraph{Discreteness.} Every budget offers its six bundles at log-ratios $\log P+\log(j/(7-j))$, $j=1,\dots,6$, a
lattice whose steps are $1.52$ times wider at the ends, where money's choices sit. The interval regression with the
lattice's own thresholds (which the source data's \texttt{lnC\_lo}/\texttt{lnC\_up} were built for) gives $0.800$,
$0.836$ and $0.845$, with curvatures of $0.52$, $0.67$ and $0.65$; on synthetic choices with $\beta=0.80$ snapped to the
lattice, the two-step returns $0.794$ and the interval regression $0.803$. The source paper's own estimators agree on
the ordering (nonlinear least squares $0.846/0.885/0.882$; multinomial logit $0.657/0.696/0.716$).

\paragraph{Position, arbitrage, trust, scale.} By randomized presentation order, money minus food is $-0.041$
($p=.016$) and $-0.032$ ($p=.088$), and food minus food is null in both arms. Money is always presented second, so its
position cannot be separated from the good inside the design ($\rho=0$, Appendix~\ref{app:proofs}); the two foods swap
between the first and third positions, their position contrast is $-0.013$, and money lies $0.036$ below the line
joining the two food positions, so a position effect would have to be sharply non-monotone to account for it. A
background stock of money, the usual explanation for the absence of present bias over money \citep{cohen2020}, works
against the finding: with utility over $c+w$ on these budgets, as $w$ runs from zero to seven times the budget, the
estimated curvature collapses from $0.67$ to $0.15$ while $\hat\beta$ rises monotonically toward one ($+0.015$ at a
third of the budget, $+0.071$ at the extreme). Trust in the experimenters interacted with the today indicator moves
the gap by $+0.001$ $[-0.022,+0.025]$ per standard deviation, the relative valuation of the goods by $-0.002$
$[-0.030,+0.022]$, and the $70$-versus-$56$ delivery moves no good ($p>.19$).

\paragraph{Who carries the gap.} By self-reported behavior (good-specific elasticities, $B=999$): everyone
$0.820/0.852/0.861$, gap $-0.037$ ($p=.002$); neither behavior ($585$ subjects) $0.838/0.857/0.865$, $-0.023$
($p=.10$); drinks ($95$) $0.737/0.822/0.868$, $-0.108$ ($p<.002$); smokes ($50$) $0.762/0.890/0.847$, $-0.106$
($p=.002$). The elasticity-free interaction is $-0.029$ for drinking ($p=.024$), $-0.026$ ($.036$) with both trait
reports and $-0.027$ ($.028$) with every control; smoking alone does not survive the full controls ($.21$). Over the
twelve-moderator battery the Romano--Wolf adjusted $p$ for drinking is $.232$, and no other moderator falls below
$.45$, so the pattern is reported as exploratory. \citet{cheung2022} relate each domain's measure to these behaviors
one at a time and find that present bias over money predicts smoking and, together with healthy food, drinking.
Entered jointly in a logit, only money survives: $+0.41$ ($0.16$), $-0.09$ and $-0.22$ for drinking; $+0.51$
($0.21$), $-0.37$ and $-0.04$ for smoking. At the advance session, where $\beta$ does not enter, drinkers allocate
more to the sooner date by $0.058$, $0.052$ and $0.076$ log points on the three goods, a domain-general difference that
the subject $\times$ good effects remove.

\paragraph{Appetite.} Per standard deviation of the re-signed six-item appetite index, $\Delta\beta$ over money /
food / food minus money is $-0.035$ / $-0.039$ / $-0.004$ $[-0.027,+0.020]$ at the on-the-day level, $-0.020$ /
$-0.015$ / $+0.005$ for the change between sessions, $-0.015$ / $-0.025$ / $-0.010$ for the advance level (a placebo
timing) and $-0.009$ / $-0.020$ / $-0.011$ for fatigue (a placebo state). The attenuation-corrected bound on the
differential is $0.042$ per standard deviation of true appetite. On the six-item index, subjects are hungrier in the
second session than in the first ($3.21$ against $2.98$, paired $p=2\times10^{-5}$), although the single hunger item
reported by the source paper does not differ between sessions.

% -------------------------------------------------------------------------------------
\section{Magnitudes, the household panel, and the type shares}
\label{app:magn}

\subsection{Three frameworks on the effort panels}

\paragraph{Existence.} Whether PB is present at all is the one question the level of the wedge answers (Rung~4).
For each subject and wage quintile in \citet{augenblick2019}, the wedge $G_q$ is the mean log allocation chosen in
advance minus the mean log allocation chosen on the day, on interior choices ($87$ subjects). The null that no
quintile has a positive wedge, $\E[G_q]\le0$ for every $q$, is rejected by the test of \citet{coxshi2023}
($p=.004$); generalized moment selection gives $p=.001$, and the test of \citet{arp2023}, which is driven by the most
violated inequality, $p=.07$.

\paragraph{The two-step estimate.} The parametric estimate is the two-step of Section~\ref{sec:coord} with the states
pooled: $0.80$ on all $100$ subjects of \citet{augenblick2019} and $0.74$ $[0.68,0.80]$ on \citet{fedyk2024}. It
assumes an isoelastic cost.

\paragraph{The measurement-error-robust set.} The set drops the cost function and keeps three assumptions: each subject
has a nondecreasing advance supply curve $e_A(W)$; the choice on the day at wage $W$ equals the advance choice at the
devalued wage, $e_T(W)=e_A(\beta W)$; and the recorded cell means equal the true ones up to a centered measurement
error. For a candidate $\beta_0$, the latent object is each subject's supply curve at the five quintile wages and at
$\beta_0$ times them, and the moments are the ten differences between recorded and latent cell means, five at each
vantage. The test asks whether some distribution of nondecreasing latent curves makes these differences mean-zero.
It uses the entropic latent-variable integration of \citet{schennach2014}, as in \citet{aguiar2021}: candidate curves
are drawn around each subject's isotonized advance curve, tilted exponentially so that the averaged moments are as
close to zero as possible, and the statistic $\mathit{TS}=n\min_\gamma\bar h(\gamma)'\hat\Omega^{-1}\bar h(\gamma)$ is
compared with the $95\%$ quantile of a $\chi^2$ with as many degrees of freedom as moments, which is conservative. The
set collects the values of $\beta_0$, on a grid of step $0.05$ from $0.55$ to $1$, that are not rejected. In
\citet{augenblick2019} it keeps all $100$ subjects, including corner choices and subjects whose choices are not
monotone; a missing today cell is filled from the subject's own advance curve, which is neutral with respect to
present bias. The set is $[0.75,0.95]$, and $\beta=1$ is rejected. In \citet{fedyk2024}, where each subject faces two
tasks on different scales, the unit is a subject--task pair; on the $120$ units with no cell mean at the bounds of $0$
or $60$ rounds the set is $[0.80,0.95]$. On all $396$ units it is empty: at a bound the recorded choice can err in one
direction only, so many choices piled at the bounds make centered errors impossible, whatever $\beta_0$ is.

\paragraph{The nonparametric bracket.} The bracket keeps only the first two assumptions. Each subject's mean advance
allocation is removed, cell means of the advance and today allocations are computed at the five wage quintiles, and
each today mean is matched into the isotonized advance schedule: if it lies between the advance means at $W_k$ and
$W_{k+1}$, then $\beta\in[W_k/W_j,W_{k+1}/W_j]$, and the bracket is the intersection over $j$. Each bracket is widened
by $1.96$ times the sum of the standard errors of the two cell means before matching, so that it is a conservative
confidence region for the identified set. It is $[0.72,1.00]$ on \citet{augenblick2019} and $[0.67,1.00]$ on
\citet{fedyk2024}. Excluding $\beta=1$ would require some today mean to lie surely below the advance mean at the next
lower wage on the grid, and none does in either panel. The bracket need not contain the measurement-error-robust
set, which pools subjects and imposes the matching identity at every wage.

\subsection{The household panel}
\label{app:ecpf}

\paragraph{Data and deterministic tests.} The data are the extract of the Spanish Continuous Family Expenditure
Survey (ECPF) prepared by \citet{adams2014} and distributed by \citet{aguiar2021}: $185$ single-person households and
$2{,}004$ couples, each observed for four consecutive quarters, with expenditure on seventeen nondurable goods and
prices discounted by the interest rate. Exponential discounting requires utility levels $v_t$ and a discount factor
$d$ such that $v_s\le v_t+d^{-t}\rho_t'(c_s-c_t)$ for every pair of quarters \citep{browning1989}. On the coarse grid
of $d$ used by \citet{aguiar2021} the deterministic test rejects it for $84.3\%$ of singles and $89.8\%$ of couples,
which the pipeline reproduces exactly; on a fine grid the rates are $77.8\%$ and $82.1\%$. For a household that
re-plans every quarter, \citet[Proposition~1]{bbc2021} show that quasi-hyperbolic discounting replaces $d^{-t}$ with
an increasing path of multipliers $m_t$; given $\beta$ and $\delta$, each ratio $m_t/m_{t-1}$ lies in
$[1/\delta,1/(\delta\beta)]$, its position depending on the household's propensity to spend out of wealth. The model
with this band, QHD($\beta_0$), is exponential discounting at $\beta_0=1$, and a lower $\beta_0$ allows everything a
higher one does. On the recorded data it rationalizes $58\%$ of singles and $53\%$ of couples with $\beta_0$
unrestricted, against $22\%$ and $18\%$ under exponential discounting with a household-specific $\delta$, and at
most $9\%$ and $4\%$ at any common $\delta$. Because the band always contains $1/\delta$, consumption paths bound
$\beta$ only from above \citep[Corollary~3]{bbc2021}.

\paragraph{Where the error is centered.} \citet{aguiar2021} allow recorded consumption to differ from true
consumption by an error that is centered across households, and test whether a distribution of true consumption that
satisfies the model's inequalities makes the errors centered. Their moments center the error in each quarter's total
expenditure, a condition they support with validation data. Under it, true consumption can be reallocated freely
across goods within a quarter, and every household in the panel, singles and couples alike, admits a true path with
exactly its recorded quarterly spending under exponential discounting at some household-specific $d$. The test then
rejects no model, which is why it does not reject exponential discounting for singles. Their rejection for couples
comes from the sampler: each household's chain starts at $d=1$, and a sampling density scaled in the units of the
couples file keeps it there. Centering the error within each of the eight groups of goods of \citet{adams2014} in each
quarter ($32$ moments) restricts reallocation to goods of the same group and gives the test power. It is a stronger
assumption than theirs and a weaker one than centering the error good by good, under which \citet[footnote~52]{aguiar2021}
reject exponential discounting strongly.

\paragraph{The test.} Each cell $(\beta_0,\delta_0)$ of Figure~\ref{fig:fieldgrid} fixes a common $\delta_0$ and
restricts every multiplier ratio to $[1/\delta_0,1/(\delta_0\beta_0)]$. Latent paths are drawn from the set that
satisfies the inequalities, by hit-and-run with Gibbs updates of each household's multiplier path, and the statistic
of \citet{aguiar2021} is compared with a $\chi^2_{32}$ ($5\%$ critical value $46.2$). Three implementation choices
differ from their code. Each household's chain starts at its best-fitting path in the cell, found by a quadratic
program, rather than at $d=1$; the sampling density is scaled to each household's own spending (a scale of $2\%$), which
makes the test invariant to the units of the data; and the chain starts from an interior point and
mixes in coordinate moves, so that it can leave goods with zero consumption. Verdicts next to the boundary of the
non-rejected region depend on the length of the chains: with $100{,}000$ steps and $6{,}000$ retained draws they move in
both directions, so the figure uses chains of $200{,}000$ steps with $10{,}000$ draws, whose verdicts a second seed
reproduces at every boundary cell. Cells away from the region are closed by nesting: rejecting $\beta_0=0.7$ at a given
$\delta_0$ rejects every higher $\beta_0$ there. The number printed in each cell is deterministic: the smallest error,
as a share of a household's spending in a group and quarter, at which its recorded consumption is consistent with the
cell, averaged over households.

\paragraph{Results.} For single-person households, exponential discounting is rejected at every $\delta_0$, and it
requires errors of at least $6.0\%$ ($\delta_0=0.98$; $9.3\%$ at $\delta_0=1$). The pairs that are not rejected are
$\beta_0\le0.99$ at $\delta_0=0.98$, $\beta_0\le0.975$ at $\delta_0=0.99$ and $0.995$, and $\beta_0\le0.95$ at
$\delta_0=1$; every pair with $\delta_0\le0.97$ is rejected. For couples, exponential discounting is rejected at
every $\delta_0$ and requires at least $6.3\%$; the pairs not rejected are $\beta_0\le0.975$ at $\delta_0=0.99$ and $0.995$ and $\beta_0\le0.95$ at
$\delta_0=1$, and every pair with $\delta_0\le0.98$ is rejected. In both samples, $\beta_0\le0.95$ with
$\delta_0\ge0.99$ requires errors of $0.7$ to $1.1\%$. Without measurement error, no pair rationalizes every
household.

\paragraph{What the contrast can and cannot show.} Seasonal tastes do not produce the gap: exponential discounting
with three common seasonal taste factors still requires errors of $5.4$ to $8.4\%$ for singles and $6.1$ to $8.1\%$
for couples. Present bias of the size measured in the experiments does not produce it either. In simulated panels
with the same prices and a $3\%$ error in each good, households with $\beta=0.9$ make exponential discounting require
errors of $0.1$ to $0.5\%$, and households with $\beta=0.7$ errors of $0.5$ to $2.8\%$, against $6.0$ to $9.3\%$ in the
data; exponential households hit by quarterly spending shocks with a standard deviation of $25\%$ require $0.1$ to
$2.2\%$. And the test has no power against such present bias: across $30$ simulated samples of $185$ households with
$\beta=0.9$ it rejects exponential discounting in none, both with the distribution of curvature of
\citet{aguiar2021} and with more elastic households, while it also rejects in none of $30$ samples of exponential
households. A free concave utility absorbs the smooth extra decline in consumption that present bias produces. What
exponential discounting cannot absorb in the data is the volatility of quarterly spending across groups of goods, which
the flexible path of multipliers that quasi-hyperbolic discounting allows absorbs as it would any variation of that
form. The contrast in Figure~\ref{fig:fieldgrid} is therefore a statement about how much measurement error each model
needs, not evidence of present bias. Income uncertainty, which also lets the growth of the marginal utility of wealth
vary across quarters \citep[Appendix~E.1]{aguiar2021}, was not tested; the results hold under perfect foresight, as
in \citet{bbc2021}.

\subsection{The type shares}
\label{app:types}

\paragraph{Margins.} In \citet{cerrone2026}, participants decide over three days how many tasks to complete at given
wages, on the day of the decision or on a later workday, and predict their own later decisions. Every payment is made
on the third day, so a decision taken on that day pairs the work with an immediate payment and a decision taken
earlier with a delayed one. For each of the $80$ participants in the authors' main sample, three margins come from
these choices, all measured in tasks: the revealed wedge over effort (the number of tasks chosen for a later workday
minus the number chosen for the same day), the revealed wedge over money (the number chosen when payment is delayed
minus the number chosen when it is immediate), and the predicted wedge over effort (the same difference in the
predictions). Each margin becomes a binary signal: present bias over effort if the effort
wedge exceeds half its cross-subject standard deviation; present bias over money if the money wedge falls below minus
half its standard deviation; and an indicator that the predicted wedge has the sign of a revealed effort wedge that
exceeds the same threshold in absolute value. The three signals define the eight cells of Table~\ref{tab:mixing}.

\paragraph{Types and the mixing matrix.} Five latent types are considered: time-consistent; na\"ive present-biased (an
effort wedge, none predicted); sophisticated present-biased (an effort wedge, predicted in full); future-biased (a
negative effort wedge, predicted); and money-driven (a money wedge only). The wedge of each biased type equals the mean
wedge of the subjects classified as biased in the data. Entry $M_{b\tau}$ is the probability that a subject of type
$\tau$ lands in cell $b$, computed from $6{,}000$ simulated subjects of each type with normal noise on each margin,
whose standard deviation is the median within-subject sampling error of that margin ($0.67$, $0.64$ and $0.79$ tasks,
against cross-subject standard deviations of $0.99$, $1.29$ and $1.41$). The noise is what makes the types overlap
across cells: a na\"ive subject, whose predicted wedge is pure noise, matches the sign of its revealed wedge about half
the time.

\begin{table}[H]\centering
\caption{Observed cell shares and the mixing matrix for the type shares ($80$ participants of \citet{cerrone2026}).}
\label{tab:mixing}
{\footnotesize
\begin{tabular}{ccccccccc}
\toprule
\multicolumn{3}{c}{Cell} & & \multicolumn{5}{c}{$M$: probability of the cell for each type} \\
\cmidrule(lr){1-3}\cmidrule(lr){5-9}
\shortstack{effort\\PB} & \shortstack{money\\PB} & \shortstack{sign\\predicted} & \shortstack{observed\\share} &
\shortstack{time-\\consistent} & na\"ive & \shortstack{sophis-\\ticated} & \shortstack{future-\\biased} &
\shortstack{money-\\driven} \\
\midrule
0 & 0 & 0 & $0.400$ & $0.553$ & $0.105$ & $0.120$ & $0.148$ & $0.061$ \\
0 & 0 & 1 & $0.025$ & $0.093$ & $0.002$ & $0.000$ & $0.685$ & $0.011$ \\
0 & 1 & 0 & $0.075$ & $0.111$ & $0.019$ & $0.024$ & $0.033$ & $0.590$ \\
0 & 1 & 1 & $0.050$ & $0.020$ & $0.000$ & $0.000$ & $0.128$ & $0.105$ \\
1 & 0 & 0 & $0.175$ & $0.093$ & $0.366$ & $0.045$ & $0.004$ & $0.011$ \\
1 & 0 & 1 & $0.237$ & $0.093$ & $0.366$ & $0.681$ & $0.000$ & $0.011$ \\
1 & 1 & 0 & $0.037$ & $0.020$ & $0.069$ & $0.008$ & $0.001$ & $0.106$ \\
1 & 1 & 1 & $0.000$ & $0.017$ & $0.073$ & $0.122$ & $0.000$ & $0.104$ \\
\bottomrule
\end{tabular}}
\end{table}

\paragraph{Bounds.} With $f$ the vector of observed cell shares and $\pi$ the vector of type shares, each bound solves
$\min\pi_\tau$ or $\max\pi_\tau$ subject to $\mathrm{lo}\le M\pi\le\mathrm{hi}$, $\pi\ge0$ and $\mathbf 1'\pi=1$, where
$[\mathrm{lo},\mathrm{hi}]$ is a simultaneous $95\%$ Bonferroni box around $f$ (sixteen one-sided bounds, $z=2.73$),
whose width adds the simulation error of $M$ to the sampling error of $f$. The program is feasible at the nominal
level. Time-consistent types are bounded in $[0.67,0.79]$, na\"ive present-biased types in $[0.09,0.32]$,
sophisticated ones in $[0,0.11]$, future-biased in $[0,0.015]$ and money-driven in $[0,0.12]$. The certificate T2
(Appendix~\ref{app:mc}) gives coverage of at least $0.975$ for such bounds on a five-type mixture, with widths from
$0.18$ to $0.47$.

\paragraph{Reading.} The bounds locate present bias over effort in a minority but do not separate na\"ive from
sophisticated types: the two intervals overlap, and since the sophisticated type predicts its full wedge, partially
sophisticated subjects are split between the two. In the raw cells, $19$ of the $36$ subjects with an effort wedge
above the threshold predict its sign, close to the half that noise alone produces for a na\"ive subject.
\citet{cerrone2026}, who estimate present bias and sophistication as continuous individual parameters on all
participants, find that $59$ of $77$ anticipate the direction of their bias, and that predictions recover about three
quarters of the average effort wedge ($3.09$ against $4.17$ tasks). The shares describe one elicitation, in one state
and one trade-off structure.

% -------------------------------------------------------------------------------------
\section{Inference: the tests, and their Monte Carlo certificates}
\label{app:inference}

Throughout, the subject is the unit of inference: decisions made by one person share that person's preferences, days
and noise, and every procedure below resamples, redraws or clusters at the level of the subject.
Table~\ref{tab:inference} lists the procedures and where each is used.

\begin{table}[H]\centering
\caption{The procedures, the hypothesis or object each delivers, and where it is used.}
\label{tab:inference}
{\footnotesize
\renewcommand{\arraystretch}{1.08}
\begin{tabular}{@{}>{\raggedright\arraybackslash}p{4.0cm}>{\raggedright\arraybackslash}p{6.6cm}>{\raggedright\arraybackslash}p{4.4cm}@{}}
\toprule
Procedure & Null hypothesis or object & Used in \\
\midrule
Subject bootstrap & intervals and two-sided $p$ for any estimate & Tables~\ref{tab:matched}--\ref{tab:goods}; App.~\ref{app:recentered}--\ref{app:goods} \\
Exact randomization inference & the assigned state moves no decision & Table~\ref{tab:recentered}, Figure~\ref{fig:recentered}; App.~\ref{app:recentered} \\
Ladder (adjacent-increment) Wald & $\beta$ equal across state cells (Rung~1) & \S\ref{sec:matched}, \S\ref{sec:margins}; App.~\ref{app:robust} \\
Distance to the monotone cone & $\beta$ monotone across cells (Rung~2) & \S\ref{sec:margins}; App.~\ref{app:robust} \\
Cox--Shi (with GMS and ARP) & no positive wedge in any cell; no shift toward taste & \S\ref{sec:margins}, \S\ref{subsec:magnitude}; App.~\ref{app:robust}, \ref{app:magn} \\
Likelihood ratio at the boundary & no true dispersion of individual $\log\beta$ & \S\ref{sec:person}; App.~\ref{app:me} \\
Romano--Wolf stepdown & no moderator moves the money--food gap & App.~\ref{app:goods} \\
Two-limit Tobit; interval regression & censored and discrete allocations & \S\ref{sec:person}; App.~\ref{app:goods} \\
Measurement-error-robust test & a latent law consistent with the model and with centered errors exists & \S\ref{subsec:magnitude}; App.~\ref{app:magn} \\
Linear program on a confidence box & bounds on type shares & \S\ref{subsec:magnitude}; App.~\ref{app:types} \\
\bottomrule
\end{tabular}}
\end{table}

\paragraph{Subject bootstrap.} Whole subjects are drawn with replacement, each bringing all of their decisions, and the
full estimation, including both steps that produce $\beta$, is repeated on each of $B$ resamples ($B=499$ unless stated
otherwise). Intervals are the $2.5$th and $97.5$th percentiles of the bootstrap distribution, and the two-sided
$p$-value is twice the share of draws on the far side of zero. With $B=499$ the smallest attainable $p$ is $2/499$,
reported as $p<.004$.

\paragraph{Exact randomization inference.} Where the assignment process is known, the sharp null that the state moves
no decision can be tested with the experiment's own randomness. Holding every choice fixed, the orderings of the
session components are redrawn $999$ times within the strata of Appendix~\ref{app:recentered}, the recentered state is
recomputed, and the estimate is re-run; the $p$-value is the share of redraws whose estimate is at least as large in
absolute value as the observed one. No large-sample approximation and no model of the errors is involved.

\paragraph{Equality across cells.} State independence (Rung~1) is tested on the cell-specific today-effects
$\theta_c$, estimated jointly in one within-subject regression with today $\times$ cell indicators. The statistic is a
Wald test on the $C-1$ adjacent increments $\theta_{c+1}-\theta_c$, with covariance from the subject bootstrap,
compared with a $\chi^2_{C-1}$. Each increment is a difference of two cell levels, so any component common to all
cells (the average present bias, a transformation bias, a fixed misprediction) cancels, which makes the test valid
under the random-$\beta$ null; and the increments span every pairwise contrast, so the test detects a response that
peaks in the middle cells, which a linear slope misses. The raw recentered product moment
$\E[(s_k-\mu_k)g_k]$ is an equally valid but less efficient member of the same class, because it leaves each
subject's mean and the wage effect in the wedge: on the stakes margin it gives $p=.29$, where the joint fit gives
$p<10^{-4}$. The difference is the optional residualization step of \citet[Algorithm~1]{borusyakhull2026}.

\paragraph{Monotonicity.} Rung~2 is tested with a distance statistic in the spirit of \citet{wolak1987}: the squared
distance, in the metric of the inverse bootstrap covariance, from the vector of adjacent increments to the cone of
monotone profiles, computed by nonnegative least squares. Its null distribution is simulated at the least favorable
point, where every increment is zero. Running it in both directions separates increasing, decreasing, and neither.

\paragraph{Moment inequalities.} Statements about the sign of a wedge in every cell, such as the existence test of
Appendix~\ref{app:magn} and the hunger margin of Appendix~\ref{app:robust}, are moment inequalities, $\E[G_j]\le0$
for every $j$. The test of \citet{coxshi2023} projects the vector of sample means onto the set of means that satisfy the
inequalities and compares the studentized distance with a $\chi^2$ whose degrees of freedom equal the number of
inequalities binding at the projection; it needs no tuning parameter. Generalized moment selection
\citep{andrewssoares2010} and the conditional test of \citet{arp2023} are reported as checks. The latter tests the most
violated inequality, correcting for its selection, and stays valid when the moments are linearly redundant, which is
why it is the engine for the paired-cell contrasts of the Monte Carlo, where Cox--Shi fails.

\paragraph{Variance components at the boundary.} In the measurement-error model, the hypothesis of no true dispersion
puts a variance on the boundary of the parameter space, so the likelihood-ratio statistic follows the mixture
$\tfrac12\chi^2_1+\tfrac12\chi^2_2$ \citep{stramlee1994} rather than a $\chi^2_2$. Intervals for the standard
deviations come from the profile likelihood and from the subject bootstrap (Appendix~\ref{app:me}).

\paragraph{Multiple moderators.} The twelve characteristics interacted with the money--food gap in the three-good
protocol are tested jointly by the stepdown procedure of \citet{romanowolf2005}, on $999$ common bootstrap draws, which
controls the family-wise error rate while allowing the statistics to be correlated (Appendix~\ref{app:goods}).

\paragraph{Censoring and discreteness.} When many allocations sit at a bound, as $86\%$ of the money allocations in
\citet{ans2015} do, a two-limit Tobit treats them as censored. In the three-good protocol each budget offers six
bundles, so a choice reveals an interval of the latent log ratio rather than a point, and the interval regression uses
the thresholds the lattice implies (Appendix~\ref{app:goods}).

\paragraph{Sets.} Three procedures deliver sets rather than tests of a point: the measurement-error-robust sets, which
collect the parameter values at which a latent distribution consistent with the model and with centered errors exists
\citep{schennach2014,aguiar2021}; the nonparametric brackets; and the linear program for type shares
(Appendix~\ref{app:magn}).

\subsection{Monte Carlo certificates}
\label{app:mc}

No test touched data before its size and power were certified on synthetic populations whose truth is known: a
twelve-configuration population (exponential; na\"ive, sophisticated and partially sophisticated present bias;
fixed-cost and temptation agents; exponential agents confounded by risk or by liquidity; projection-biased;
state-dependent $\beff$; myopic; domain-specific; aggregated households), always with measurement error and
preference heterogeneity, and built to fool the tests rather than flatter them. Table~\ref{tab:mc} reports the blocks
quoted in the text. Block T12 computes the recentered product moment subject by subject and tests its mean with a
$t$-statistic across subjects, which is the subject-clustered form of the test.

\begin{table}[!ht]\centering
\caption{The blocks quoted in the text ($200$--$500$ replications, $n=120$--$185$, nominal size $5\%$).}
\label{tab:mc}
{\small
\begin{tabular}{>{\raggedright\arraybackslash}p{3.0cm}>{\raggedright\arraybackslash}p{7.6cm}cc}
\toprule
Block & Decisive scenario & Size & Power \\
\midrule
T1 decomposition & risk- and liquidity-confounded agents fool the raw wedge test; the controlled-cell test clears them & .035--.045 & 1.00 \\
T2 type shares & linear-program bounds on a five-type mixture; widths .18--.47 & coverage $\ge.975$ & --- \\
T3 trait vs state & the state ladder $(.90,.75,.60)$ recovered; projection bias mimics the whole gradient and only a matched-state comparison separates them & .055--.07 & 1.00 \\
T7 noise layer & deterministic revealed preference falsely rejects a true exponential $83\%$ of the time under $10\%$ measurement error & --- & --- \\
T10 robust sets & truth inside the set in $.92$--$1.00$ of replications, corner choices included; exponential discounting rejected in every replication under present bias and in none under exponential discounting & --- & --- \\
T12 random-$\beta$ null & the level statistic rejects under every state-independent $\beta$; the recentered product holds $.054/.044/.056$ & 1.00/.05 & 1.00 \\
T13 which statistic & level / linear slope / ladder Wald / paired cells (ARP) / paired cells (Cox--Shi): size $1.00/.035/.045/.043/\mathbf{.83}$; power against a mid-peaked response $1.00/.04/1.00/.74$ & & \\
T14 moment forms & partition / product / directed oracle / one-sided / basis: power against a monotone response $.815/.859/.859/.919/.873$, against a mid-peaked one $.599/.047/.854/.047/.792$ & .046--.052 & \\
T15 the revealed-preference program & concave utility: identified set $\{0.80\}$, $\{1.00\}$; convex cost, five-price two-vantage design: $[0.65,0.85]$, $[0.80,1.20]$ & --- & --- \\
T16 background stock & curvature collapses from $0.67$ to $0.15$ while $\hat\beta$ rises toward one; the six-bundle lattice costs the two-step $0.006$ of $\beta$ & --- & --- \\
\bottomrule
\end{tabular}}
\end{table}

% -------------------------------------------------------------------------------------
\section{Two policy designs}
\label{app:policy}

The two designs below answer questions the paper raises, what liquidity does to effort and whether a policy moves the
plan or the act, with the methods the replication data allow. Neither carries a claim in the main text.

\paragraph{Liquidity on effort.} In \citet{kaur2025}, $408$ piece-rate workers in Odisha were paid their accumulated
earnings at the end of the round or, in the treatment group, on an earlier day announced in advance, with the cash day
varying across $28$ cohorts (round $\times$ wave). The source paper's specification reproduces at $+0.097$ (cash
$\times$ post-payment, with a baseline-productivity control; cash $\times$ announced $-0.002$), but the raw gap between
cash and control groups is positive at every event time, including the days before the announcement ($+0.13$ to
$+0.19$). Within worker, a stacked event study by cohort (worker and cohort $\times$ day fixed effects, with the
pre-announcement days as reference) and the imputation estimator of \citet{bjs2024} give: announced but unpaid days
$-0.000$ $[-0.065,+0.057]$ and $+0.015$ $[-0.065,+0.086]$; paid days $+0.026$ $[-0.059,+0.100]$ and $+0.048$
$[-0.048,+0.149]$ ($+1.7\%$ and $+3.1\%$ of baseline output); the day after receipt $+0.080$ and $+0.126$
$[+0.005,+0.242]$. There is no anticipation, and a liquidity channel can move a plan-versus-act wedge over effort by at
most about $0.15$ in a design that randomizes it. Over money the same channel is first order: randomized windfalls move
measured time inconsistency \citep{cassidy2019}.

\paragraph{The full retirement age.} The Social Security full retirement age is $65$ for cohorts born up to $1937$,
rises by two months per birth year to $66$ for $1943$, stays there to $1954$, and rises again toward $67$, a
cohort-staggered dose on the value of working past $65$ \citep{mastrobuoni2009}. In the RAND HRS, $13{,}158$ people born
$1931$--$1957$ both stated at ages $51$--$61$ a probability of working full-time past $65$ (the plan) and were observed
at $65$--$66$ (the act). The plan exceeds the act in every birth year (group means: born $1931$--$37$,
$0.219/0.176$; $1938$--$42$, $0.223/0.199$; $1943$--$54$, $0.258/0.215$; $1955$--$57$, $0.312/0.234$). Per year of
additional retirement age, with a linear cohort trend: plan $-0.047$ ($p=.006$), act $+0.005$ ($.79$), wedge $-0.052$
($.010$); with a quadratic trend: $+0.001$, $+0.033$, $-0.032$ ($.22$); in a stacked difference-in-differences of the
first phase-in cohorts against $1935$--$37$: $-0.030$, $+0.015$, $-0.045$ ($.19$). The direction is consistent and the
precision insufficient: about five hundred people per birth year, interviewed every two years, against a dose of two
months per year.

% -------------------------------------------------------------------------------------
\section{Data and code}
\label{app:data}

Every dataset is a public replication package or an openly distributed research dataset; no new data were collected.
Table~\ref{tab:datasets} lists them: the three designs of Table~\ref{tab:timing}, which carry the main results, and the
further datasets, each with the feature of its design that the paper uses and the place where it is used.

\begin{table}[H]\centering
\caption{The datasets, what each contributes, and where it is used.}
\label{tab:datasets}
{\scriptsize
\renewcommand{\arraystretch}{1.08}
\begin{tabular}{@{}>{\raggedright\arraybackslash}p{3.3cm}>{\raggedright\arraybackslash}p{10.25cm}>{\raggedright\arraybackslash}p{2.05cm}@{}}
\toprule
Dataset & Design, and what the paper uses & Used in \\
\midrule
\multicolumn{3}{@{}l@{}}{\textit{The three designs of Table~\ref{tab:timing}}} \\
\citet{augenblick2019} & $100$ participants allocate real-effort transcription tasks on seven dates over six weeks, at wages from $1$ to $35$ cents, in advance and on the day, with incentivized predictions; whether a decision comes before or after the ten mandatory tasks of a session is randomized, which assigns the state & \S\ref{sec:state}, \S\ref{sec:magn}; App.~\ref{app:recentered}, \ref{app:robust}, \ref{app:magn} \\
\citet{ans2015} & $101$ participants allocate real effort and money over seven weeks in parallel protocols, each with choices made in advance and choices that involve the present & \S\ref{sec:coord}, \S\ref{sec:memodel}, \S\ref{sec:person}, \S\ref{sec:welfare}; App.~\ref{app:me} \\
\citet{cheung2022} & $697$ high-school students make convex-time-budget choices over money, healthy food and unhealthy food with one price grid, delay and delivery protocol; appetite measured at both sessions; published individual estimates with standard errors & \S\ref{sec:margins}, \S\ref{sec:person}; App.~\ref{app:me}, \ref{app:goods} \\
\addlinespace[4pt]
\multicolumn{3}{@{}l@{}}{\textit{Further state margins, and two restrictions that survive}} \\
\citet{marchgluth2025} & $70$ participants choose between foods once hungry and once sated: hunger manipulated within person & \S\ref{sec:margins}; App.~\ref{app:robust} \\
\citet{schilbach2019} & cycle-rickshaw drivers in Chennai, breathalyzed daily, with randomized incentives for sobriety; beliefs about their own sobriety and a later choice of commitment & \S\ref{sec:margins}, \S\ref{sec:welfare}; App.~\ref{app:robust} \\
\citet{bessone2021} & low-income adults in Chennai with randomized sleep interventions and measured sleep; a real-effort discounting task in the design of \citet{augenblick2019} & \S\ref{sec:margins}; App.~\ref{app:robust} \\
\citet{hoel2022} & payment in cash or through a mobile payment application, randomized between subjects, in an intertemporal-choice task & App.~\ref{app:robust} \\
HeartSteps and Drink Less & two micro-randomized trials \citep{klasnja2019,bell2023}: a prompt randomized within person, with known probability, at each decision point & App.~\ref{app:robust} \\
\addlinespace[4pt]
\multicolumn{3}{@{}l@{}}{\textit{The person, magnitudes and types}} \\
\citet{cerrone2026} & present bias over effort and over money elicited in the same effort choices, revealed and predicted, in the same participants & \S\ref{sec:person}, \S\ref{sec:magn}; App.~\ref{app:types} \\
\citet{fedyk2024} & an independent real-effort design with beliefs about one's own and others' present bias & \S\ref{subsec:magnitude}, \S\ref{sec:welfare}; App.~\ref{app:magn} \\
Spanish ECPF, as prepared by \citet{adams2014} and \citet{aguiar2021} & a household consumption panel: deterministic tests of exponential and quasi-hyperbolic discounting, and the measurement-error-robust test extended to quasi-hyperbolic discounting & \S\ref{sec:magn}; App.~\ref{app:ecpf} \\
\addlinespace[4pt]
\multicolumn{3}{@{}l@{}}{\textit{Plans, beliefs and commitment}} \\
\citet{oscarsson2020} & New Year's resolutions: monthly beliefs about success against realized success over a year & \S\ref{sec:welfare} \\
\citet{carrera2022} & gym members: attendance goals and predictions against gym records, recalled past attendance, and the choice of commitment contracts & \S\ref{sec:welfare} \\
\citet{allcott2022} & smartphone users with a randomized screen-time bonus and limit tool: predicted against measured use, and willingness to pay for the limit tool & \S\ref{sec:welfare} \\
\citet{sadoff2020} & grocery deliveries in two food deserts: food chosen in advance and at delivery, and demand for commitment & \S\ref{sec:welfare} \\
\citet{kaur2015} & data-entry workers on piece rates, with the option of dominated commitment contracts & \S\ref{sec:welfare} \\
RAND HRS & the longitudinal file ($1992$--$2022$): stated probabilities of working past $62$ and full-time past $65$, against labor-force status at those ages & \S\ref{sec:welfare}; App.~\ref{app:policy} \\
\addlinespace[4pt]
\multicolumn{3}{@{}l@{}}{\textit{Policy, background and code checks}} \\
\citet{kaur2025} & $408$ piece-rate workers in Odisha; the treatment group is paid its accumulated earnings early, on an announced day that varies across $28$ cohorts & App.~\ref{app:policy} \\
Compilation of \citet{lockwood2020} & published estimates of $\beta$, collected study by study & \S\ref{sec:intro}; Fig.~\ref{fig:three} \\
Two panels of \citet{eis2020} & convex-time-budget choices over money from \citet{andreoni2012a} and \citet{carvalho2016}; the revealed-preference code reproduces the published pass rates exactly & code check \\
\bottomrule
\end{tabular}}

\smallskip
{\scriptsize\raggedright The pipeline also processes \citet{andreoni2012b} and \citet{dertwinkel2021}, whose results are not reported in this version.\par}
\end{table}

\noindent All estimates are produced by one pipeline around a shared inference engine (GARP and CCEI, Afriat
feasibility, two-limit Tobit, cluster bootstrap, GMS, Cox--Shi, ARP, entropic latent-variable integration), with a
separate module for the household panel; every module writes a log and a summary file, and a verifier checks every
number quoted in the paper against those files ($402$ checks, $0$ mismatches on the run that produced them). Four data
traps are worth recording. The three-good study's online appendix prints its craving items with anchors reversed
relative to the delivered data (signed from three behavioral checks, six items give $\alpha=.76/.80$); its dictionary
codes gender and the drinking and smoking indicators the other way round from the data; its estimation file carries
the session-1 hunger response on both sessions' rows; and the effort panel's randomized order exposes decisions
unequally, which is what Section~\ref{sec:recentered} corrects.

\end{document}
